\chapter{Introduction}
This chapter first discusses how macroscopic matter is built from its microscopic constituents and how it interacts with its surroundings. The discussion leads to the concepts of internal energy and equilibrium state, to the definitions for different boundary types, and to concepts of heat and work. The chapter then describes the state variables that are needed for a macroscopic description of matter. While doing so, the zeroth law of thermodynamics is encountered. Finally, the chapter closes by providing a roadmap to these lecture notes and a brief discussion on how thermodynamics and continuum mechanics are positioned with respect to each other and with respect to other classical fields of engineering.

\section{Internal energy and equilibrium state}
Under close inspection, ordinary matter is observed to consist of particles such as atoms, molecules, and ions. All of these particles interact through forces that have origins in the Coulomb interactions between their charged constituents, namely electrons and nuclei. Even when particles carry no net electric charge, the internal distribution of positive and negative charge gives rise to effective interactions that can be attractive at moderate separations and strongly repulsive at very short separations. As a result, the interaction between two individual particles can be understood as a consequence of potential energy that has a minimum at a characteristic distance.

For a collection of particles, the collective effect of different interaction strengths and geometries leads to different macroscopic states. In a solid, particles typically reside close to each other and their preferred separations vibrating about equilibrium positions. In a liquid, particles reside close together and near their preferred separations but continually rearrange over time. In a gas, the average separation is large and interactions are weak except when particles briefly pass each other close by in an event that can be described as an instantaneous collision. For an illustration, the sketches in Fig.~\ref{fig:water-solid-liquid-vapor-molucules} show collections of water particles in various configurations. The macroscopic state of matter is therefore governed by a competition between the interparticle potential energy, which favours certain separations and arrangements, and the kinetic energy of the particles, which includes contributions from translational, rotational, and vibrational motion. The microscopic kinetic and interaction potential energies of the particles, excluding the center-of-mass kinetic energy, gravitational energy, and nuclear binding energy, are referred to as internal energy.
\begin{figure}[!htb]
    \centering
    \includegraphics[width=0.18\linewidth]{\rootfolder/figs/fig-water-frozen.png}
    \includegraphics[width=0.18\linewidth]{\rootfolder/figs/fig-water-melting.png}
    \includegraphics[width=0.18\linewidth]{\rootfolder/figs/fig-water-liquid.png}
    \includegraphics[width=0.18\linewidth]{\rootfolder/figs/fig-water-evaporating.png}
    \includegraphics[width=0.18\linewidth]{\rootfolder/figs/fig-water-gaseous.png}
    \caption{Left to right: solid, melting, liquid, evaporating, and gaseous water with hydrogen bonds (cyan dashed lines) illustrated.}% Snapshots from molecular dynamics simulations.}% Vrbka, L. and Jungwirth, P.,  Journal of Physical Chemistry B 110: 18126 (2006).}
    \label{fig:water-solid-liquid-vapor-molucules}
\end{figure}

As the motion of the particles in classical sense is governed by Newton's laws, the particles continually exchange energy and momentum through their interactions, the forces between them. A collection of particles is said to be in a thermodynamic equilibrium state when this exchange has led to a distribution of energies and configurations that appears time-independent when averaged over a short period of time. Although individual particles continue to move, interact, and exchange energy and momentum, the overall distribution of energies and configurations in an equilibrium state remains approximately steady in time.

\section{Boundary types, heat, and work}
Consider two collections of particles, each of which has settled into its own equilibrium state. Regard each such collection as a system. A boundary between systems is any surface that separates them and determines the mechanisms by which energy can be transferred between the systems.

A diathermal boundary allows energy to be transferred in a way that can change the energy or configuration distributions of the particles within systems, without requiring macroscopic motion of the boundary in system's rest frame. When two systems are separated by a diathermal boundary, they are said to be in thermal contact. Typical examples include two solid objects touching each other, a thin metal wall separating two volumes of liquid, and a solid object surrounded by air or liquid water. The portion of net energy transfer across a diathermal boundary that is mediated by microscopic motion of particles is referred to as heat. Heat therefore describes energy in transfer and is not a property contained within a system. Energy transfer across a diathermal boundary can be mediated also by, e.g, electromagnetic fields. By contrast, an adiabatic boundary prevents all energy transfer that would alter the energy or configuration distributions of the particles within the systems without macroscopic motion of the boundary between them. Practical examples include thick insulation, vacuum gaps, Faraday cages, or thin highly reflective material layers.

When two systems are separated by a movable boundary, they are said to be in mechanical contact. A typical example is a movable piston in a cylinder that separates the volume of gas inside the cylinder chamber from the outsides of it. By contrast, a stationary rigid boundary does not move, preventing mechanical contact. The energy transferred across a mechanical contact through macroscopic motion of the boundary is referred to as mechanical work. Like heat, mechanical work is energy in transfer and is not a property contained within a system.

In general, work is any energy transfer across a boundary that is not heat, including mechanical and electromagnetic field work. A boundary can also combine different types, allowing both heat and work to be exchanged, or neither.


\section{Density, specific volume, and molar volume}
A macroscopic description of matter should be consistent with experimental observations while remaining as simple as possible and expressible in terms of a small number of variables. Although the state of matter is determined from the microscopic energy and configuration distributions, a macroscopic description should not depend on the detailed positions and motions of individual particles, but instead capture their collective effect.

As the forces between particles depend strongly on the distances between them, and the motion of each particle in response to these forces depends on its mass, the macroscopic state of matter can be expected to depend on the mass--volume ratio of the system. This motivates introducing macroscopic variables referred to as density and specific volume
\begin{align}
\rho &= \frac{m}{V},\\
v &= \frac{V}{m} = \frac{1}{\rho}.
\end{align}
Density measures how closely packed the particles are within the boundary, while specific volume indicates how much space is available per unit mass of matter. 

As different materials consist of different types of particles with different particle masses, it is sometimes useful to describe the ratio of volume $V$ to the number of particles $N$ rather than to the total mass. The typical number of particles in a macroscopic volume, however, is large, so one uses the Avogadro constant $N_A=6.02214076\times 10^{23}\,\text{1}/\text{mol}$ and the number of moles $n$ instead of the absolute number of particles. This leads to a macroscopic variable definition referred to as molar volume 
\begin{align}
    V_m=\frac{V}{n}=N_A\frac{V}{N}=Mv,
\end{align}
where the link to specific volume is provided by the molar mass $M=m/n$ of the matter.

Although macroscopic quantities, density, specific volume, and molar volume can be defined for a collection of particles whether or not it is in an equilibrium state. The SI-units for the space and mass related quantities are listed in Table~\ref{tab:volume-mass-related-units}.
\begin{table}[!htb]
    \centering
    \begin{threeparttable}
    \caption{Units of volume and mass related quantities.}
    \begin{tabular}{lll}
         \toprule 
         Symbol & Name & SI-unit \\
         \midrule 
         $V$ & Volume & m$^3$\\
         $m$ & Mass& kg\\
         $\rho$ & Density & $\text{kg}/\text{m}^3$\\
         $v$ & Secific volume & $\text{m}^3/\text{kg}$\\
         $M$ & Molar mass & $\text{kg}/\text{mol}$\\
         $V_m$ & Molar volume & $\text{m}^3/\text{mol}$\\
         \bottomrule
    \end{tabular}
    \end{threeparttable}
    \label{tab:volume-mass-related-units}
\end{table}

\section{Pressure}
Consider then a cylinder filled with gas and closed with a piston, as in Fig.~\ref{fig:pressure}. A gas particle that interacts with the piston changes its velocity over a short time interval and, according to Newton’s laws, imparts a small impulse to the piston. Such collisions occur extremely frequently, and the cumulative effect of many impulses produces a force that, on macroscopic time scales, appears continuous. To keep the piston from moving, an opposing force must be applied normal to the piston surface.
\begin{figure}[!htb]
    \centering
    \includegraphics[width=0.45\linewidth]{\rootfolder/figs/fig-piston.pdf}
    \includegraphics[width=0.45\linewidth]{\rootfolder/figs/fig-cylinder.pdf}
    \caption{(Left) A gas particle imparts impulse on a piston. (Right) A force is need to keep the piston in place.}% Adapted from \url{https://tikz.net/piston/} under CC~BY-SA~4.0. Original by Izaak Neutelings.}
    \label{fig:pressure}
\end{figure}

The magnitude of this force depends on both the state of the gas and the area of the boundary over which the particles interact. This motivates defining pressure as a macroscopic property to characterize the mechanical effect of a system on its boundary. For a system, pressure is therefore defined as the normal component of the macroscopic force per unit area exerted by the particles of the system on a surface, namely
\begin{align}
    \lim_{A\to 0}\frac{F_n}{A}=p.
\end{align}
The SI-units for the pressure related quantities are listed in Table~\ref{tab:pressure-related-units}. While the SI-unit for pressure is $\text{Pa}$, in engineering applications a more common unit is $1\,\text{bar}=10^5\,\text{Pa}$.
\begin{table}[!htb]
    \centering
    \begin{threeparttable}
    \caption{Units of pressure related quantities.}
    \begin{tabular}{cll}
         \toprule 
         Symbol & Name & SI-unit \\
         \midrule 
         %$p$ & Paine & Pa \\
         %$T$ & Lämpötila & K\\
         $F$ & Force & $\text{N}$\\
         $A$ & Area& $\text{m}^2$\\
         $p$ & Pressure & $\text{Pa}\,(=\text{N}/\text{m}^2)$\\
         \bottomrule
    \end{tabular}
    \end{threeparttable}
    \label{tab:pressure-related-units}
\end{table}

If the piston is pushed inward by increasing the applied force, the volume of the gas decreases. For a sufficiently large compression, part of the gas may condense into a liquid, so that the cylinder contains both gas and liquid. Because of different densities for liquid and gas, gravity ultimately separates them, forming an interface. In equilibrium, the pressure in the gas on a flat interface is equal to the pressure in the liquid at the same location. Otherwise, a net force would act on the interface and it would move until equilibrium is restored. On a strongly curved liquid-gas interfaces, as in small bubbles or dropplets, also the surface tension playes a role, allowing for a small pressure difference between the liquid and gas phases. On both flat and curved surfaces, this shows that the surface on which pressure is defined does not need to be a solid wall like the surface of a piston. An imaginary surface, drawn within gas or liquid, suffices as well. In thermal equilibrium, the particles on either side of such a surface exchange momentum through the interparticle forces, giving rise to a normal force per unit area. Pressure is therefore a macroscopic property, not a property of a particular container wall enclosing the gas or liquid. In gases and liquids, it can nevertheless be applied and measured directly by compressing them in a cylinder with a piston. 

\section{A remark on solids and stress}
While the interparticle forces acting across an imaginary surface within an equilibrium gas or liquid are strictly normal to that surface, regardless of the surface orientation, this is generally not true for solids. In a solid body, the forces transmitted across an internal imaginary surface through a point depend explicitly on the surface orientation. Consequently, the force per unit area at a given point is no longer invariant with direction and cannot be fully described by a single scalar value. Instead, a tensor-valued quantity $\bm{\sigma}$ called stress is introduced to account for both the magnitude and the direction of the force $\bm{F}$ relative to the orientation of the surface normal $\bm{\hat{n}}$.
\begin{align}
    \lim_{A\to 0}\frac{\bm{F}}{A}=\bm{\sigma}\cdot\bm{\hat{n}}
\end{align}

This directional dependence of internal forces is a result of a non-isotropic microscopic structure. For instance, extruded plastics possess molecular arrangements that are preferentially aligned with the direction of extrusion. Similarly, in cold-rolled steels, the microscopic metallic grains are mechanically forced to orient along the rolling direction. This microstructural realignment induces macroscopic anisotropy, rendering the material's mechanical strength, elasticity, and thermal response direction-dependent.

For an isotropic solid in an unloaded equilibrium state, such as unworked cast iron, the stress is isotropic and represented by pressure. However, the experimental measurement of this pressure differs from how one measures pressure in liquids and gases. Because a solid cannot rearrange its particles freely to transmit forces uniformly, one cannot submerge a sensor inside it. Instead, a surrounding gas or liquid is typically used as an intermediary. The gas or liquid transmits a uniform force over the solid's external boundary, ensuring that the induced internal force per unit area corresponds uniformly to the applied gas or liquid pressure. 

\section{The zeroth law and temperature}
Where pressure, more generally stress, characterizes macroscopic description of the mechanical interactions within a system and across a boundary, one can ask whether a macroscopic characterization exists for interactions across a thermal contact.  To explore this, suppose two systems are separated by a diathermal boundary and are otherwise isolated from external influences. The systems are able to exchange heat. After some time, the interaction across the boundary no longer produces any net heat exchange and the systems settle into a mutual equilibrium. The two systems are then said to be in thermal equilibrium. 

The zeroth law of thermodynamics states that thermal equilibrium is transitive: if system A is in thermal equilibrium with system C, and system B is also in thermal equilibrium with system C, then A and B are in thermal equilibrium with each other. 
The zeroth law enables defining a system property referred to as temperature: two systems are in thermal equilibrium if and only if they have the same temperature. The symbol usually chosen for the temperature defined in this way is $T$.

\section{A gas thermometer}
While the zeroth law enables defining the concept of temperature, it does not directly provide a way to measure or calculate it. A detour is needed. To quantify the temperature $T$ of a system, one may place a gas-filled bulb in thermal contact with the system, connect the gas to a manometer, as illustrated in Fig.~\ref{fig:manometer}, and make experimental observations.

\begin{figure}[!htb]
    \centering
    %\includegraphics[width=0.4\linewidth]{\rootfolder/figs/fig-manometer.pdf}
    \input{\rootfolder/code/fig-manometer}
    \caption{A gas-filled bulb in thermal contact with a system and connected to a manometer. The manometer measures the gas pressure and volume.}
    \label{fig:manometer}
\end{figure}

According to the zeroth law, the gas in the bulb will reach thermal equilibrium with the system and thus have the same temperature $T$. By changing the volume of the gas by adjusting the amount of liquid in the manometer, while keeping the state of the system unchanged, one observes many different pairs of gas pressure $p$ and volume $V$ that all correspond to the same system temperature. This demonstrates that the temperature of the gas, which is the same as that of the system, depends on its pressure and volume in a reproducible way, justifying the existence of a functional relation
\begin{align}
    T = f(p,V).
\end{align}
Further experiments with sufficiently dilute gases show that, while the gas remains in thermal equilibrium with the system, its pressure and volume are inversely proportional to each other. This empirical observation implies the existence of a relation of the form
\begin{align}
    pV = g(T),
\end{align}
where $g(T)$ depends only on the temperature.

One may now choose a reference state for the gas and label its temperature, pressure, and volume as $T_\text{ref}$, $p_\text{ref}$, and $V_\text{ref}$, respectively. For any other state of the gas with pressure $p$ and volume $V$, the temperature $T$ then satisfies
\begin{align}
    \frac{g(T)}{g(T_\text{ref})} = \frac{pV}{p_\text{ref}V_\text{ref}}.
\end{align}
Since temperature has not yet been assigned a numerical scale, the functional form of $g(T)$ may be chosen by convention. A convenient choice is
\begin{align}
    g(T) = \text{const}\times T,
\end{align}
where the constant depends on the chosen gas and the specific experimental setup. This choice defines a temperature scale such that the temperature of the system can be determined by measuring the pressure and volume of the gas and comparing them to the reference values
\begin{align}
    \frac{T}{T_\text{ref}} = \frac{pV}{p_\text{ref}V_\text{ref}}.
\end{align}

It is standard to choose the reference temperature to correspond to the state at which water can exist simultaneously in solid, liquid, and gas phases, and to assign it the value $T_\text{ref} = 273.16\,\mathrm{K}$. With this choice, the gas thermometer measures temperature on the Kelvin scale (K) which serves also as the SI-unit for temperature. Another commonly used temperature scale is the Celsius scale that is related to the Kelvin scale by a constant
\begin{align}
    T[^\circ\text{C}]=T[\text{K}]-273.15.
\end{align}

\section{Extensive and intensive quantities}
The volume, mass, and internal energy of a system are extensive quantities, meaning they depend on the amount of matter in the system. In contrast, density, specific volume, molar volume, pressure, specific internal energy, and temperature are intensive quantities, which do not depend on the system size. If two systems with identical properties are combined, the extensive variables sum, while the intensive variables remain unchanged. Theories describing the behaviour of matter usually rely on intensive quantities and variables. That in mind, plain upper-case symbols are reserved for extensive variables, upper-case symbols with a subscript $m$ are reserved for intensive variables expressed relative to the number of moles, and plain lower-case symbols are reserved for intensive variables expressed relative to the mass, as in Table~\ref{tab:intensive-extensive-convention}. Temperature ($T$), pressure ($p$), and mass ($m$) are the primary exceptions to this symbolic convention.
\begin{table}[!htb]
    \centering
    \begin{threeparttable}
    \caption{Extensive and intensive symbol conventions for matter-related quantities.}
    \begin{tabular}{lcc}
         \toprule 
         Type & Symbol & Unit\\
         \midrule 
         Extensive & $A$ & $[A]$\\
         Intensive (specific to mass) & $a$ & $[A]/\text{kg}$\\
         Intensive (specific to amount of substance) & $A_m$ & $[A]/\text{mol}$\\
         \bottomrule
    \end{tabular}
    \end{threeparttable}
    \label{tab:intensive-extensive-convention}
\end{table}

\section{Matter in motion}
Matter, be it gas, liquid, or solid, doesn't have to remain still. Also its shape can change. Liquids and gases flow, hence they are called fluids, while solids may bend elastically and oscillate back and forth, or deform inelastically. Typically the macroscopic scale of the volume the matter occupies is vast in comparison to the interparticle distances, and the time scale of motion one observes is orders of magnitudes longer than the average time between inter-particle collisions. Under such conditions, different parts of macroscopic volumes of matter can be treated as individual systems connected by boundaries, with each system remaining in thermodynamic equilibrium state but potentially in a different one from the neighbouring system's state. The matter then is said to remain in a local thermodynamic equilibrium and thermodynamic relations between state variables are assumed to hold locally.

In such macroscopic continuum approximation of matter, each subsystem of the macroscopic volume is assumed infinitesimally small. The state properties are described as continuous fields that smoothly depend on the point-wise location $\bm{x}$ and time $t$. To track matter in motion, one supplements the thermodynamic state fields with a velocity field, $\bm{v}(\bm{x},t)$, which describes the macroscopic speed and direction of matter at any given point and time. Evolution equations for the state fields are derived by applying conservation laws of mass, momentum, and energy to these moving subvolumes.

In classical mechanics of point particles and rigid bodies, one is accustomed to tracking the motion of a fixed mass subjected to a force. In continuum mechanics, this corresponds to tracking a volume containing a fixed mass, referred to as material volume, subject to a distribution of force through out the volume and its surface. In case of solids, the deformation of the object is limited by the material strength and the approach of tracking the deformation of the volume from its initial state is the golden standard. Such view is generally referred to as Lagrangian view. In case of fluids that may flow indefinitely, tracking the deformation becomes cumbersome and eventually impossible. A continuous flow may deform an initially well described cubic volume into an infinitely long and infinitesimally thin volume. As a remedy, one switches to observing how the properties of the fluid change when it flows through a volume chosen to make the observations convenient. Such view is referred to as Eulerian. In Lagrangian view, the velocity field $\bm{v}(\bm{x},t)$ is interpreted to move a point $\bm{x}(t)$ according to an ordinary differential equation
\begin{align}
    \frac{d\bm{x}(t)}{dt}=\bm{v}(\bm{x}(t),t).
\end{align}
The connection to Eulerian view is established through the acceleration of the point $\bm{x}(t)$
\begin{align}
    \frac{d^2\bm{x}(t)}{dt^2}=\frac{d\bm{v}(\bm{x}(t),t)}{dt}=\partial_t\bm{v}(\bm{x}(t),t)+\bm{v}(\bm{x}(t),t)\cdot\nabla\bm{v}(\bm{x}(t),t).
\end{align}
The first term $\partial_t\bm{v}$ on the right describes the local Eulerian rate of change of the velocity at a point while the second term $\bm{v}\cdot\nabla\bm{v}$ describes the local acceleration the moving point $\bm{x}(t)$ is subjected to in the Lagrangian view if the velocity field were stationary in time.

\section{A roadmap for the course}
The continuum mechanics framework can be built on three core principles: the measurement of observable quantities; the fundamental conservation laws of mass, momentum, and energy; and the natural tendency of local differences in macroscopic quantities to smooth out over time in the absence of external influences. 

For fluids, measurements at stationary equilibrium provide values not only for pressure, volume, and temperature, but also for how these state variables are coupled to one another, and how they change differentially with respect to each other.
\begin{center}
\begin{forest} downright
[Stationary time-independent measurements:
  [{$p$: pressure}]
  [{$v$, $\rho$: specific volume, density}]
  [$T$: temperature]
  [$\beta_p$: differential dependency of volume on temperature at constant pressure]
  [$\kappa_T$: differential dependency of volume on pressure at constant temperature]
]
\end{forest}
\end{center}
Applying conservation laws over material volumes leads to a set of governing partial differential equations that describe the evolution of the density field, the velocity field, and the specific internal energy field:
\begin{center}
\begin{forest} downright
[Conservation laws over material volume:
  [$\partial_t\rho$: evolution equation for density]
  [$\partial_t\bm{v}$: evolution equation for velocity]
  [$\partial_tu$: evolution equation for specific internal energy]
]
\end{forest}
\end{center}
However, this system of partial differential equations is fundamentally unclosed—and therefore mathematically intractable—without knowing how the fluid's internal energy depends on temperature and density (or specific volume). This crucial closing information is obtained by studying energy conservation in slow, quasistatic flows. In these idealized processes, the heat added relates directly to temperature changes via heat capacities, the macroscopic gradients driving dissipation are negligible, and the mechanical work performed on the fluid is governed strictly by pressure via volume changes. Through this analysis, also the state function known as enthalpy is encountered.
\begin{center}
\begin{forest} downright
[Measurements and energy conservation in quasistationary motion:
  [{$c_p, c_v$}: specific heat capacities at constant pressure and volume]
  [$u$: state function for specific internal energy]
  [$h$: state function for specific enthalpy]
]
\end{forest}
\end{center}

Further investigation of quasistatic flows reveals that certain thermodynamic states are completely inaccessible to without the addition or removal of heat. This observation leads to the discovery of entropy and the formulation of the Helmholtz and Gibbs free energies.
\begin{center}
\begin{forest} downright
[Quasistationary motion and existence of inaccessible states without heat transfer:
  [$s$: state function for specific entropy]
  [$f$: state function for specific Helmholtz free energy]
  [$g$: state function for specific Gibbs free energy]
]
\end{forest}
\end{center}
Once either the Helmholtz or Gibbs free energy has been modeled or quantified for a given substance, all other thermodynamic properties can be computed through differentiation, completing the theoretical framework. The web of interdependencies between the quantities is illustrated in Fig.~\ref{fig:roadmap-thermo}.

\begin{figure}[!htb]
    \centering
    \begin{tikzpicture}
    
    \node[draw,circle,minimum size=2.5cm,inner sep=0pt,align=center] (A) at (0,0) {$\begin{aligned}
        &\text{Experiments:}\\
        &p,v,T,\\
        &\beta_p,\kappa_T,\\
        &c_p,c_v
        \end{aligned}$};

    \node[draw,circle,minimum size=1.8cm,inner sep=0pt,align=center] (B) at (6,0) {$\begin{aligned}
        &\text{Computable:}\\
        &u,h,s\end{aligned}$};
    
    \node[draw,circle,minimum size=1.8cm,inner sep=0pt] (C) at (3,4) {$\begin{aligned}
    &\text{Computable}\\
    &f,g\end{aligned}$};
   
    \draw [-latex,postaction={decorate,decoration={raise=2ex,text along path,text align=center,text={Integrate}}}] (A) to[bend right=30] (B);

    \draw [-latex,postaction={decorate,decoration={raise=-2ex,text along path,reverse path,text align=center,text={Sum and multiply}}}] (B) to[bend left=30] (C);

    \draw [-latex,postaction={decorate,decoration={raise=2ex,text along path,reverse path,text align=center,text={Differentiate}}}] (C.210) to[bend right=30] (A.90);

    \draw [-latex,postaction={decorate,decoration={raise=2ex,text along path,text align=center,text={Differentiate}}}] (C.-30) to[bend left=30] (B.90);
    \end{tikzpicture}
    \caption{Relations of thermodynamic quantities.}
    \label{fig:roadmap-thermo}
\end{figure}

With the thermodynamic closure established, the fluid equations coupled to thermodynamical relations become a useful predictive tool. The Helmholtz and Gibbs free energies serve as the mathematical foundation for numerical property libraries, allowing evaluation of fluid's thermodynamic state using only localized measurements of temperature, pressure, or volume. Once the state relations are embeded analytically to the fluid equations, macroscopic balance equations of mass, linear and angular momentum, energy, and entropy enable analyzis of open engineering domains. Together, fluid mechanics and thermodynamics provides the tools necessary to design, optimize, and understand the fundamental working principles of many devices from pumps, fans, compressors, and turbines to heat exchangers, heat pumps, power plants, and internal combustion engines. The joint framework of is illustrated in Fig.~\ref{fig:roadmap-joined}.

\begin{figure}
    \centering
    \begin{tikzpicture}

    \node[draw,circle,minimum size=2.5cm,inner sep=0pt,align=center] (A) at (0,0) {$\begin{aligned}
        &\text{Fluid mechanics:}\\
        &\text{Evolution equations}\\
        &\text{for}\,\partial_t\rho,\, \partial_t\bm{v},\,\partial_tu
        \end{aligned}$};

    \node[draw,circle,minimum size=2.5cm,inner sep=0pt,align=center] (B) at (8,0) {$\begin{aligned}
        &\text{Thermodynamics:}\\
        &\text{State relations for}\\
        &\rho, T, p, u, h, s, g, f
        \end{aligned}$};

    \node[draw,circle,minimum size=2.5cm,inner sep=0pt,align=center] (C) at (4,7) {$\begin{aligned}
        &\text{Applications:}\\
        &\text{heat rate}\,\dot{Q},\\
        &\text{work rate}\,\dot{W},\\
        &\text{and efficiency in}\\
        &\text{power plants}\\
        &\text{heat pumps,}\\
        &\text{and engines}
        \end{aligned}$};
    
    \draw [latex-,postaction={decorate,decoration={raise=2ex,text along path,text align=center,text={Embed}}},postaction={decorate,decoration={raise=-3.5ex,text along path,text align=center,text={analytically}}}] (A) to[bend right=30] (B);

    \draw [-latex,
    postaction={decorate,decoration={raise=2ex,text along path,text align=center,text={Solve or integrate}}},
    postaction={decorate,decoration={raise=-3.5ex,text along path,text align=center,text={over a volume}}}] 
    (A) to[bend left=30] (C);

    \draw [latex-,postaction={decorate,decoration={raise=2ex,text along path,text align=center,text={Calculate for}}},postaction={decorate,decoration={raise=-3.5ex,text along path,text align=center,text={assessment}}}] (C) to[bend left=30] (B);

    \end{tikzpicture}
    \caption{Framework of thermodynamics and fluid mechanics for applications}
    \label{fig:roadmap-joined}
\end{figure}

\chapter{State of Pure Substances}
In this chapter, relations between the state variables of pure substances are first illustrated with phase diagrams. Then mathematical expressions and numerical libraries are introduced to model the experimentally observed state relations with equations of state. Finally, computation of quantities needed in developing the thermodynamic theory is discussed.

\section{Learning goals}

\begin{itemize}
    \item[\elementary] The student points locations of critical point, triple line, and phase transition lines in phase diagrams of pure substances.    \item[\elementary] The student recalls the different phases of ordinary matter, the transitions between them, and is able to sketch phase diagrams.
    \item[\elementary] The student calculates vapor quality of a state point on a given $p$-$v$ or $v$-$T$ diagram.
    \item[\elementary] The student uses \verb|pyromat| or \verb|coolprop| to calculate pressure, density, or temperature for a given state point.
    \item[\elementary] The student calculates the compressibility factor from given data to estimate the validity of ideal gas model.
    \item[\intermediate] The student estimates the isobaric thermal expansion and the isothermal compressibility coefficients from $p$-$v$ or $v$-$T$ diagrams for a given state point.
    \item[\intermediate] The student estimates the coefficients of a Van der Waals model or another two-parameter model from measured data.
    \item[\advanced] The student manipulates differential forms to obtain relations between partial derivatives.
    \item[\advanced] The student programs a python script utilizing \verb|pyromat| or \verb|coolprop| to create a $p$-$v$ or a $v$-$T$ diagram.
\end{itemize}

\section{The $p$-$v$-$T$ phase diagram}
In mechanics, matter is classified into solid states that can sustain shape under shearing forces and into fluid states that cannot. In thermodynamics, fluid states are classified not only into liquid and gas states but subclassified further into superheated vapor and supercritical fluid states. In short, superheated vapor is gaseous while supercritical fluid has characteristics of both liquid and gas. Also mixtures of solid, liquid, and gas states in common thermodynamic equilibrium are encountered. In a three-dimensional space with pressure $p$, specific volume $v$, and temperature $T$ as coordinates, thermodynamic states of a pure substance form a two-dimensional, so-called $p$-$v$-$T$ surface. Schematic illustrations of such surfaces are provided in Fig.~\ref{fig:pvT-diagrams}, with temperature and pressure isocontours, i.e., isotherms and isobars, drawn with black and white dashed curves.
\begin{figure}[!htb]
    \centering
    \includegraphics[width=0.95\linewidth]{\rootfolder/figs/fig-pvT-expanding.pdf}
    \includegraphics[width=0.95\linewidth]{\rootfolder/figs/fig-pvT-contracting.pdf}
    \caption{Schematic $p$-$v$-$T$ diagrams for pure substances that expand (top) or contract (bottom) during melting.}
    \label{fig:pvT-diagrams}
\end{figure}

Important classifiers are observed in Fig.~\ref{fig:pvT-diagrams}. First is the critical point with pressure $p_\text{crit}$, temperature $T_\text{crit}$, and specific volume $v_\text{crit}$. At temperatures above that of the critical point, a gaseous state is referred to as gas and cannot be condensed into a liquid by increasing the pressure alone. This has important consequences for the design of engineering applications that often rely on phase transitions. At temperatures below that of the critical point, a gaseous state is labelled as superheated vapor and can be condensed into liquid by increasing the pressure alone. If both pressure and temperature are above those of the critical point, the fluid state is labelled as supercritical. Second is the triple line that, if projected onto the $p$-$T$ plane, is observed to reduce into a single point. This indicates that the triple line is characterized by a single pressure value $p_\text{trp}$, a single temperature value $T_\text{trp}$, and a range of values for specific volume. For pressure values below that of the triple line, a stable liquid phase does not exist. Further important classifiers in Fig.~\ref{fig:pvT-diagrams} are the thick colored curves that separate single-phase regions from two-phase regions and, when projected onto the $p$-$T$ plane, form the two-colored co-existence curves representing mixtures of two phases. For convenience, the labels for the single-phase regions, the mixture-phase regions, and the names of the transitions between them are illustrated also in Fig.~\ref{fig:pvT-labels}.
\begin{figure}[!htb]
    \centering
    \includegraphics[width=0.9\linewidth]{\rootfolder/figs/fig-pvT-labels.pdf}
    \caption{Labels of the single-phases regions and the transitions between them.}
    \label{fig:pvT-labels}
\end{figure}

Following the isobars and isotherms on the $p$-$v$-$T$ surface, one may build understanding of the substance behaviour. For pressure values below that of the triple line, keeping the pressure constant while decreasing the temperature of a gas results in reduction of the gas volume, until vapor begins to deposit into a solid. A phase-mixture of gaseous vapor and solid in a thermodynamic equilibrium then forms, with pressure and temperature common for both phases. If volume is then reduced while keeping the temperature constant, also pressure remains constant until all of the vapor is deposited into solid. After that, further reduction of temperature at constant pressure can continue while the substance state remains in solid phase, leading typically to minor reduction in the substance volume. 

For pressure values above that of the triple line but below that of the critical point, keeping the pressure constant while decreasing the gas temperature again leads to decrease in the gas volume, until the vapor starts condensing into liquid. A phase-mixture of liquid and vapor then forms, with pressure and temperature common for both phases. If the volume is then reduced while keeping temperature constant, pressure remains constant until all vapor has condensed into liquid. Reducing temperature at constant pressure can then continue, until liquid at some point begins to freeze into solid. For many substances, the phase-mixture of solid and liquid under constant temperature and pressure reduces in volume upon freezing but for some, like water, freezing leads to expansion in substance volume instead. This behaviour is observed also in Fig.~\ref{fig:pvT-diagrams}. Once all of the liquid has frozen into solid, temperature can be further reduced at constant pressure, leading again typically to a reduction of the substance volume. At pressures above that of the critical point, no phase boundary nor coexistence curve exists between gas- and liquid-like phases. Instead, the transition is continuous and happens in a single phase. 

The phase-mixture regions might seem somewhat exotic at first but are not. A region where two phases coexist in equilibrium can be obtained, e.g., by placing a kettle filled with liquid water on a stove and turning the heat on. Once the water starts boiling, water vapor forms and fills the remaining space under the kettle lid, pushing air out. A mass of water then exists within a volume in both liquid and vapor states, sharing a common temperature and pressure. Because liquid water has higher density than water vapor, gravity separates the liquid and vapor phases into physically different regions within the kettle. As a system, the water in the kettle nevertheless has one mass and one volume, and therefore, as a system, can be assigned a single specific volume or density value, corresponding to a state. In this respect the triple line is a unique collection of states. As the line separates two phase-transition regions, matter on the line can exist simultaneously in all three primary phases, as solid, liquid, and gaseous vapor.

\section{The $p$-$v$ and $v$-$T$ diagrams, vapor quality}
Most engineering applications involve fluids at pressures above that of the triple line. To represent the fluid state and the associated isotherms and isobars quantitatively, the $p$–$v$ and $v$–$T$ projections replace the schematic $p$–$v$–$T$ surface, as shown for water in Fig.~\ref{fig:pv-vT-H2O}. Particular attention is given to the two-phase region, representing mixtures of liquid and vapor, and to the behaviour of temperature and pressure isocontours in the $p$–$v$ and $v$–$T$ planes.
\begin{figure}[!htb]
    \centering
    \begin{tikzpicture}
    \node[above right,inner sep = 0] (image) at (0,0) {
        \includegraphics[width=0.85\linewidth]{\rootfolder/figs/fig-pv-T-H2O.pdf}
        %\includegraphics[width=0.85\linewidth]{\rootfolder/figs/fig-pv-H2O.pdf}
    };
    \begin{scope}[relative to node=image]
        % \foreach \ratio in {-1,-0.9,...,1}{
        %    \draw[help lines,purple] (-1,\ratio) -- (1,\ratio);
        %    \draw[help lines,purple] (\ratio,-1) -- (\ratio,1);
        % }
        % \draw (-1,0) -- (1,0);
        % \draw (0,-1) -- (0,1);
        \coordinate (crit) at (-0.56,0.65);
        \coordinate (trpLiq) at (-0.675,-0.772);
        \coordinate (trpVap) at (0.694,-0.772);
        \fill (crit) circle (2.5pt);
        \draw[latex-,thick,shorten <=3pt] (crit) to[bend left=30] ($(crit)+(0.6,0)$) node[anchor= west,align=left] {Critical point:\\$\left\{\begin{aligned}&T=373.95^\circ\text{C}\\
        &p=220.64\,\text{bar}\\
        &v=0.0031056\,\text{m}^3/\text{kg}
        \end{aligned}\right.$};
        
        \draw[very thick] (trpLiq)--(trpVap) coordinate[pos=0.4] (trp);
        \draw[latex-,thick,shorten <=1pt] (trp) to[bend left=30] ($(trp)+(0.3,1.0)$) node[anchor= west,align=left] {Triple line:\\$\left\{\begin{aligned}&T=0.01^\circ\text{C}\\
        &p=6.12\times 10^{-3}\,\text{bar}\\
        &v\in[0.001,206.00]\,\text{m}^3/\text{kg}
        \end{aligned}\right.$}; 
        
    \end{scope}
    \end{tikzpicture}
    \begin{tikzpicture}
    \node[above right,inner sep = 0] (image) at (0,0) {
        %\includegraphics[width=0.85\linewidth]{\rootfolder/figs/fig-vT-H2O.pdf}
        \includegraphics[width=0.85\linewidth]{\rootfolder/figs/fig-vT-p-H2O.pdf}
    };
    \begin{scope}[relative to node=image]
        % \foreach \ratio in {-1,-0.9,...,1}{
        %    \draw[help lines,purple] (-1,\ratio) -- (1,\ratio);
        %    \draw[help lines,purple] (\ratio,-1) -- (\ratio,1);
        % }
        % \draw (-1,0) -- (1,0);
        % \draw (0,-1) -- (0,1);
        \coordinate (crit) at (0.284,-0.575);
        \coordinate (trpLiq) at (-0.75,-0.7);
        \coordinate (trpVap) at (-0.75,0.7);
        \fill (crit) circle (2.5pt);
        \draw[latex-,thick,shorten <=3pt] (crit) to[bend left=30] ($(crit)+(-0.3,0.8)$) node[anchor= south,align=left] {Critical point:};
        \draw[very thick] (trpLiq)--(trpVap) coordinate[pos=0.4] (trp);
        \draw[latex-,thick,shorten <=0pt] (trp) to[bend right=30] ($(trp)+(0.5,0.6)$) node[anchor= south,align=left] {Triple line:}; 
    \end{scope}
    \end{tikzpicture}
    
    \caption{The $p$-$v$ (top) and $v$-$T$ (bottom) diagrams for water.}
    \label{fig:pv-vT-H2O}
\end{figure}

In the $p$-$v$ and $v$-$T$ diagrams of Fig.~\ref{fig:pv-vT-H2O}, the red curve denotes the saturated liquid states, while the green curve represents the saturated vapor states. The two curves start from the opposing ends of the triple line and meet at the critical point. Together they form the saturation curve which, together with the triple line, encloses the liquid-vapor two-phase region. Elsewhere the fluid exists in single phase. Although Fig.~\ref{fig:pv-vT-H2O} is specific to water, the same definitions and qualitative behavior apply to all pure substances. 

As disccussed in connection with the $p$-$v$-$T$ surface, temperature and pressure both remain constant during a phase transion when volume change is not constrained. This is reflected in both the $p$-$v$ and the $v$-$T$ diagram. In the $p$-$v$ diagram, isotherms are horizontal in the two-phase region, and the phase transition appears as a horizontal line between the saturated liquid and saturated vapor states. In the $v$-$T$ diagram, the roles are reversed. Isobars are vertical within the two-phase region, and the phase transition appears as a vertical line between the saturated liquid and saturated vapor states.

What changes during a phase transition is the specific volume of a substance, and a quantity referred to as vapor quality. In a volume $V$ containing mass $m$ in a two-phase state, the liquid will occupy a volume $V_\text{liq}$ and have mass $m_\text{liq}$ while the vapor will occupy a volume $V_\text{vap}$ and have mass $m_\text{vap}$. Naturally, the volumes and masses sum together according to
\begin{align}
    V_\text{liq}+V_\text{vap}&=V,\\
    m_\text{liq}+m_\text{vap}&=m.
\end{align}
On the saturated liquid curve, no vapor exists at all while on the saturated vapor curve, only vapor exists. To reflect this behaviour, vapor quality is defined as the ratio of the vapor mass to the total fluid mass
\begin{align}
    x=\frac{m_\text{vap}}{m}=\frac{m_\text{vap}}{m_\text{liq}+m_\text{vap}}.
\end{align}
The liquid and vapor parts have specific volumes corresponding to the respective states at the saturated liquid and vapor curves 
\begin{align}
    v_\text{liq}&=\frac{V_\text{liq}}{m_\text{liq}},\\
    v_\text{vap}&=\frac{V_\text{vap}}{m_\text{vap}}, 
\end{align}
and a state in the two-phase region, having vapor quality $x$, can be shown to have a specific volume of
\begin{align}
    v=(1-x)v_\text{liq}+x\,v_\text{vap}.
\end{align}
Conversely, the vapor quality in the two-phase region can be determined from the specific volume with
\begin{align}
    x=\frac{v-v_\text{liq}}{v_\text{vap}-v_\text{liq}}.
\end{align}
Chosen isocontours of the vapor quality defined in this manner are drawn in the $p$-$v$ and $v$-$T$ diagrams of water in Fig.~\ref{fig:pv-vT-H2O} with black dashed lines. One should remember that the vapor quality is specified only within the two-phase region and with values $x\in[0,1]$.

\begin{example}
    Hot steam at $T=150^\circ\text{C}$ is produced for an industrial process by boiling water with an electric resistance heater within a pressure vessel. The pressure vessel has a volume of $1000\,\ell$ and the liquid water inside is observed to occupy a volume of $200\,\ell$. Estimate the mass of water vapor inside the pressure vessel and calculate the vapor quality for the water in the vessel.

    From the $p$-$v$ diagram in Fig.~\ref{fig:pv-vT-H2O}, one can estimate the specific volumes of saturated liquid and saturated vapor using the $150^\circ\text{C}$ isotherm and the saturated liquid and vapor curves:
    \begin{align*}
        v_\text{liq}&\approx 1\times 10^{-3}\,\text{m}^3/\text{kg}, &
        v_\text{vap}&\approx 4\times 10^{-1}\,\text{m}^3/\text{kg}.
    \end{align*}
    Using the definition of specific volume, the masses of liquid water and water vapor are then calculated from the volumes and the specific volumes
    \begin{align*}
        m_\text{liq}&=\frac{V_\text{liq}}{v_\text{liq}}\approx 200\,\text{kg}, &
        m_\text{vap}&=\frac{V_\text{vap}}{v_\text{vap}}=\frac{V-V_\text{liq}}{v_\text{vap}}\approx 2\,\text{kg},
    \end{align*}
    and the total mass of water inside the pressure vessel can be calculated to be $m\approx202\,\text{kg}$. The specific volume for the water inside the vessel is therefore $v=V/m\approx 5\times 10^{-3}\,\text{m}^3/\text{kg}$ and the vapor quality is calculated to be
    \begin{align*}
        x=\frac{v-v_\text{liq}}{v_\text{vap}-v_\text{liq}}\approx 0.01.
    \end{align*}
    The points from which the values used are obtained are illustrated below.
    \begin{center}
    \begin{tikzpicture}
    \node[above right,inner sep = 0] (image) at (0,0) {
        \includegraphics[width=0.9\textwidth]{\rootfolder/figs/fig-pv-T-H2O.pdf}
        %\includegraphics[width=0.9\textwidth]{\rootfolder/figs/fig-pv-H2O.pdf}
    };
    \begin{scope}[relative to node=image]
        %\foreach \ratio in {-1,-0.9,...,1}{
        %    \draw[help lines,purple] (-1,\ratio) -- (1,\ratio);
        %    \draw[help lines,purple] (\ratio,-1) -- (\ratio,1);
        %}
        %\draw (-1,0) -- (1,0);
        %\draw (0,-1) -- (0,1);

        \coordinate (vliq) at (-0.667,0.13);
        \coordinate (vvap) at (-0.01,0.13);
        \coordinate (vpoint) at (-0.497,0.13);
        \fill (vliq) circle (2pt);
        \fill (vvap) circle (2pt);
        \fill (vpoint) circle (2pt);
        \draw[dashed,purple,thick] (vliq)--++(0,-1.1);
        \draw[dashed,purple,thick] (vvap)--++(0,-1.1);
        \draw[dashed,purple,thick] (vpoint)--++(0,-1.1);

        \draw[latex-,thick,shorten <=5pt] (vvap) to[bend right=10] ($(vvap)+(0.3,0.5)$) node[anchor= west] {$\left\{\begin{aligned}&T=150^\circ\text{C}\\
        &v_\text{vap}\approx0.4\,\text{m}^3/\text{kg}
        \end{aligned}\right.$};

        \draw[latex-,thick,shorten <=5pt] (vliq) to[bend right=25] ($(vliq)+(1.2,-0.4)$) node[anchor=west] {$\left\{\begin{aligned}&T=150^\circ\text{C}\\
        &v_\text{liq}\approx1\times 10^{-3}\,\text{m}^3/\text{kg}
        \end{aligned}\right.$};

        \draw[latex-,thick,shorten <=5pt] (vpoint) to[bend right=25] ($(vpoint)+(1,0.1)$) node[anchor=west] {$\left\{\begin{aligned}&T=150^\circ\text{C}\\
        &v\approx5\times 10^{-3}\,\text{m}^3/\text{kg}
        \end{aligned}\right.$};
        
    \end{scope}
    \end{tikzpicture}
    \end{center}
\end{example}
\begin{example}
    A pressure vessel with a volume of $30\,\text{m}^3$ contains $1000\,\text{kg}$ of water at a pressure of $p=10\,\text{bar}$. What is the temperature of water in the vessel? At what temperature would all water in the vessel become vapor?

    The specific volume of water in the vessel is $v=V/m=0.03\,\text{m}^3/\text{kg}$. From the $v$-$T$ diagram in Fig.~\ref{fig:pv-vT-H2O}, one finds the temperature at this specific volume and pressure of $10\,\text{bar}$ to be $T\approx 180^\circ\text{C}$. Both the volume of the pressure vessel and the mass of water inside it remain fixed even if temperature were to change. The temperature at which all of the water inside would be vapor, is then found on the green saturated vapor curve at the given specific volume, corresponding to $T\approx 280^\circ\text{C}$. Both state points are illustrated below.
    \begin{center}
    \begin{tikzpicture}
    \node[above right,inner sep = 0] (image) at (0,0) {
        \includegraphics[width=0.9\textwidth]{\rootfolder/figs/fig-vT-p-H2O.pdf}
        %\includegraphics[width=0.9\textwidth]{\rootfolder/figs/fig-vT-H2O.pdf}
    };
    \begin{scope}[relative to node=image]
        % \foreach \ratio in {-1,-0.9,...,1}{
        %    \draw[help lines,purple] (-1,\ratio) -- (1,\ratio);
        %    \draw[help lines,purple] (\ratio,-1) -- (\ratio,1);
        % }
        % \draw (-1,0) -- (1,0);
        % \draw (0,-1) -- (0,1);

        \coordinate (vpoint) at (-0.255,-0.31);
        \coordinate (vvap) at (0.024,-0.31);
        
        \fill (vvap) circle (2pt);
        \fill (vpoint) circle (2pt);
        \draw[dashed,purple,thick] (vvap)--++(0,-0.6);
        \draw[dashed,purple,thick] (vpoint)--++(0,-0.6);
        \draw[dashed,,purple,thick] ($(vpoint)+(-0.6,0)$)--($(vvap)+(0.2,0)$);

        \draw[latex-,thick,shorten <=5pt] (vvap) to[bend left=20] ($(vvap)+(0.4,0.3)$) node[anchor= west] {$\left\{\begin{aligned}&v=0.03\,\text{m}^3/\text{kg}\\
        &T\approx280^\circ\text{C}
        \end{aligned}\right.$};

        \draw[latex-,thick,shorten <=5pt] (vpoint) to[bend left=25] ($(vpoint)+(0.3,0.7)$) node[anchor=west] {$\left\{\begin{aligned}&v=0.03\,\text{m}^3/\text{kg}\\
        &p=10\,\text{bar}\\
        &T\approx180^\circ\text{C}
        \end{aligned}\right.$};
        
    \end{scope}
    \end{tikzpicture}    
    \end{center}
\end{example}

\section{Equations of state}
Mathematically, a general two-dimensional surface in a three-dimensional space with coordinates $p$, $v$, and $T$, can be represented by an equation of the form
\begin{align}
    f(p,v,T)=0.
\end{align}
Experimental observations show that, in the two-phase regions, pressure and temperature can both stay constant while the specific volume varies. Real phase transitions may exhibit delayed boiling or condensation due to kinetic effects and the absence of nucleation sites. The equilibrium phase boundaries discussed here nevertheless provide an accurate approximation for most engineering applications. In these regions, pressure and temperature therefore satisfy a relation of the form
\begin{align}
    g(p,T)=0,
\end{align}
while, for a fixed temperature or pressure, the specific volume parametrizes the two-phase mixture states. Across the boundaries between single- and two-phase regions, the derivatives of the state surface change abruptly. This is observed in the $p$-$v$ and $v$-$T$ diagrams for water in Fig.~\ref{fig:pv-vT-H2O}, where the isotherms and isobars vary smoothly outside the liquid–vapor region but change direction abruptly on the saturation curves. Because of this, it is practically impossible to represent the entire state surface by a single analytic equation that behaves smoothly everywhere. Therefore, mathematical expressions referred to as equations of state are developed with constrained applicability domains. 

The arguably simplest equation of state, found experimentally and applicable for dilute gases, is referred to as the ideal gas equation 
\begin{align}
    pv=R_sT.
\end{align}
The constant $R_s=R/M$ is gas specific, with $R=8.3145\,\text{J}/\text{mol}/\text{K}$ the universal gas constant and $M$ the molar mass of the gas. The mathematical form of the equation implies that the ideal gas isotherms in a $p$-$v$ diagram look like hyperbolas on a linear-linear scale. Taking logarithms from both sides of the equation and rearranging leads to $\log_{10}p=\log_{10}(R_sT)-\log_{10}v$. In a $p$-$v$ diagram using a log-log scale, the ideal gas isotherms would therefore appear as straight downward trending lines with a slope of -1. This behavior is observed in the $p$-$v$ diagram of water in Fig.~\ref{fig:pv-vT-H2O}. The slope of the $700^\circ\text{C}$ isotherm is very close to -1 throughout while for lower temperatures, the slope is quite close to -1 if the pressure is low enough. 

Close to the saturated vapor curve or near the critical point, the real isotherms deviate significantly from the hyperbolas of the ideal gas model (the straight downward sloping lines in log-log scale). Such deviations become important in engineering applications operating at high pressures or near phase transitions, as ideal gas models may produce large errors. Therefore more accurate models for the gaseous phase have been developed. Historically the most influential has been the Van der Waals equation of state
\begin{align}
    p&=\frac{R_sT}{v-b_s}-\frac{a_s}{v^2}, & &\left\{\begin{aligned}
    a_s&=\frac{27}{64}\frac{R_s^2T_\text{crit}^2}{p_\text{crit}},\\
    b_s&=\frac{R_sT_\text{crit}}{8p_\text{crit}}.
\end{aligned}\right.
\end{align}
Further improvements include for example the Dieterici equation of state
\begin{align}
    p&=\frac{R_sT}{v-b_s}\exp\left(-\frac{a_s}{vR_sT}\right), & &\left\{\begin{aligned}
    a_s&=4e^{-2}\frac{R_s^2T_\text{crit}^2}{p_\text{crit}},\\
    b_s&=e^{-2}\frac{R_sT_\text{crit}}{p_\text{crit}},
    \end{aligned}\right.
\end{align}
and the Redlich-Kwong equation of state
\begin{align}
    p&=\frac{R_sT}{v-b_s}-\frac{a_s}{\sqrt{T} v(v+b_s)}, & &\left\{\begin{aligned}
    a_s&=\frac{1}{9(2^{1/3}-1)}\frac{R_s^2 T^{5/2}_\text{crit}}{p_\text{crit}},\\
    b_s&=\frac{2^{1/3}-1}{3}\frac{R_sT_\text{crit}}{p_\text{crit}},
    \end{aligned}\right.
\end{align}
Of these the Van der Waals and Redlich–Kwong models belong to the class of cubic equations of state while all three belong to the family of two-parameter models. The equations are modifications of the ideal gas equation to account for molecular interaction effects and finite molecular size, and reduce to the ideal gas model at appropriate limits. The expressions for the coefficients in the two-parameter models are found by requiring both the first and the second derivative of pressure with respect to volume to be zero at temperature and pressure corresponding to the critical point. This reflects the behaviour of measured data. 
\begin{example}
    Writing the Van der Waals model as
    \begin{align*}
        (p+a_s\rho^2)(v-b_s)=R_sT,
    \end{align*}
    one can interpret the parameter $a_s$ as a correction to pressure associated with intermolecular interactions that scale quadratically with respect to gas density. The parameter $b_s$ can be interpreted to account for reduction in accessible volume due to finite molecular size. Writing the equation further as
    \begin{align*}
        pv(1+a_s\rho^2/p)(1-b_s/v)=R_sT,
    \end{align*}
    one observes the Van der Waals model to reduce to the ideal gas equation when both $p\gg a_s\rho^2$ and $v\gg b_s$ hold true.
\end{example}
\begin{example}
    Using the Van der Waals model, the first and second derivatives of pressure with respect to specific volume at constant temperature become
    \begin{align*}
        \left(\frac{\partial p}{\partial v}\right)_T&=-\frac{R_sT}{(v-b_s)^2}+\frac{2a_s}{v^3}, &
        \left(\frac{\partial^2 p}{\partial v^2}\right)_T&=\frac{2R_sT}{(v-b_s)^3}-\frac{6a_s}{v^4}.
    \end{align*}
    Requiring both derivatives to vanish when pressure and temperature correspond to the critical point values, one can solve for
    \begin{align*}
        a_s&=\frac{9R_sT_\text{crit}v_\text{crit}}{8}, & b_s&=\frac{v_\text{crit}}{3}.
    \end{align*}
    Here $v_\text{crit}=v(p_\text{crit},T_\text{crit})$ is not the measured value but is implicitly determined from the equation of state
    \begin{align*}
        p_\text{crit}=\frac{R_sT_\text{crit}}{v_\text{crit}-b_s}-\frac{a_s}{v_\text{crit}^2}
    \end{align*}
    leading to 
    \begin{align*}
        v_\text{crit}=\frac{3}{8}\frac{R_sT_\text{crit}}{p_\text{crit}}
    \end{align*}
    Substituting back, one finds the expressions for the Van der Waals coefficients in terms of the critical point temperature and pressure
    \begin{align*}
        a_s&=\frac{27}{64}\frac{R_s^2T_\text{crit}^2}{p_\text{crit}}, &
        b_s&=\frac{R_sT_\text{crit}}{8p_\text{crit}}.
    \end{align*}
\end{example}

The Van der Waals equation, the Dieterici equation, and the Redlich-Kwong equation all describe state surfaces of type $f(p,v,T)=0$ and improve upon the ideal gas equation. They nevertheless fall short of correct behaviour when the nature of the state surface changes abruptly from the gaseous state to the two-phase liquid-vapor mixture state where pressure and temperature follow a relation of the form $g(p,T)$ instead. The shortfall is demonstrated in Fig.~\ref{fig:pv-model-comparison-CO2} for carbon dioxide using the Redlich-Kwong equation. For temperatures above $30.98^\circ\text{C}$, the critical temperature of CO$_2$, the Redlich-Kwong isotherms align quite well with the measured data. The isotherm corresponding to the critical temperature and the predicted location of the critical point are slightly off, since a point with three coordinates generally cannot be exactly matched with only two parameters. The isotherms below that of the critical temperature, however, show unphysical behaviour between the saturated liquid and vapor states. A region exists where pressure increases with increasing volume, whereas the correct isotherms remain flat.
\begin{figure}[!htb]
    \centering
    \begin{tikzpicture}
    \node[above right,inner sep = 0] (image) at (0,0) {
        \includegraphics[width=0.82\linewidth]{\rootfolder/figs/fig-pv-model-comparison-CO2.pdf}
    };
    \begin{scope}[relative to node=image]
        % \foreach \ratio in {-1,-0.9,...,1}{
        %    \draw[help lines,purple] (-1,\ratio) -- (1,\ratio);
        %    \draw[help lines,purple] (\ratio,-1) -- (\ratio,1);
        % }
        % \draw (-1,0) -- (1,0);
        % \draw (0,-1) -- (0,1);
        \coordinate (crit) at (-0.64,-0.054);
        \coordinate (critRK) at (-0.58,-0.054);
        \fill (crit) circle (1.5pt);
        \fill (critRK) circle (1.5pt);
        \draw[latex-,thick,shorten <=5pt] (crit) to[bend left=20] ($(crit)+(0.7,0.5)$) node[anchor= west,align=left] {Correct critical point:};

        \draw[latex-,thick,shorten <=5pt] (critRK) to[bend left=20] ($(crit)+(0.8,0.2)$) node[anchor= west,align=left] {Redlich-Kwong\\critical point:};
        
    \end{scope}
    \end{tikzpicture}
    
    \caption{Selected isotherms for CO$_2$ using different equations of state.}
    \label{fig:pv-model-comparison-CO2}
\end{figure}

The predictive capability of the two-parameter equations of state is improved when they are used together with the so-called Maxwell construct which replaces the unphysical portion of the isotherm with a constant-pressure segment, producing a relation of the form $g(p,T)=0$. The associated constraints for the construction are
\begin{align}
    \int_{v_\text{liq}(T)}^{v_\text{vap}(T)} p(v,T)dv&=p_\text{sat}(T)(v_\text{vap}(T)-v_\text{liq}(T)),\\
    p(v_\text{liq}(T),T)&=p_\text{sat}(T),\\
    p(v_\text{vap}(T),T)&=p_\text{sat}(T),
\end{align}
where $p_\text{sat}(T)$ is the model-specific saturation pressure at a temperature $T$ below the critical temperature, reflecting the $g(p,T)=0$ condition. The corresponding model-specific saturated liquid and vapor volumes $v_\text{liq}(T)$ and $v_\text{vap}(T)$ define the phase boundaries between which the two-phase mixture is predicted to exist.

\begin{example}
    Applying the Maxwell construct starts by identifying a plausible range for $p_\text{sat}(T)$. In Fig.~\ref{fig:pv-model-comparison-CO2}, one observes the $T=12^\circ\text{C}$ Redlich-Kwong isotherm to have a local minimum and maximum between which the prediction for $p_\text{sat}(12^\circ\text{C})$ must lie. Generally, any local minimum or maximum of a $T<T_\text{crit}$ Redlich-Kwong isotherm can be found by calculating the derivative of the pressure with respect to specific volume
    \begin{align*}
        \left(\frac{\partial p}{\partial v}\right)_T=-\frac{R_sT}{(v-b_s)^2}+\frac{a_s}{\sqrt{T}}\frac{2v+b_s}{v^2(v+b_s)^2},
    \end{align*}
    setting the derivative to zero, and solving the roots of the resulting quartic polynomial 
    \begin{align*}
        v^4+2\left(b_s-\frac{a_s}{R_sT\sqrt{T}}\right)v^3+b_s^2\left(1+\frac{3a_s}{b_sR_sT\sqrt{T}}\right)v^2-\frac{a_sb_s^3}{R_sT\sqrt{T}}=0.
    \end{align*}
    A polynomial of degree four has four roots. The physically relevant ones are real-valued and larger than $b_s$, the limit where the Redlich-Kwong model diverges. 
    
    Once an appropriate range has been identified for the saturation pressure, the Redlich-Kwong equation of state is rearranged into a cubic polynomial equation for specific volume
    \begin{align*}
        v^3-\frac{R_sT}{p}v^2+\left(\frac{a_s}{p\sqrt{T}}-\frac{R_sTb_s}{p}-b_s^2\right)v-\frac{a_sb_s}{p\sqrt{T}}=0,
    \end{align*}    
    one guesses the value of $p_\text{sat}(T)$, and finds the roots of the cubic polynomial, of which the outer most correspond to the values $v_\text{liq}(T)$ and $v_\text{vap}(T)$. Then one uses the predicted saturated liquid and vapor values to check if the constraint of the Maxwell construction holds. The guess is then refined by applying any standard bracketing method. Below is an example implementation using polynomial root finding from numpy and a bracketing method from scipy library:    \lstinputlisting[language=Python,firstline=1,lastline=60,firstnumber=1]{\rootfolder/code/example-redlich-kwong-saturation.py} 
    Running the script results in Redlich-Kwong estimates $p_\text{sat}(12^\circ\text{C})=50.9\,\text{bar}$, $v_\text{liq}(12^\circ\text{C})=1.47\,\ell/\text{kg}$, and $v_\text{vap}(12^\circ\text{C})=6.28\,\ell/\text{kg}$ that are reasonably close to the measured values.
    \begin{center}
    \begin{tikzpicture}
    \node[above right,inner sep = 0] (image) at (0,0) {
        \includegraphics[width=0.85\linewidth]{\rootfolder/figs/example-redlich-kwong-saturation.pdf}
    };
    \begin{scope}[relative to node=image]
        % \foreach \ratio in {-1,-0.9,...,1}{
        %    \draw[help lines,purple] (-1,\ratio) -- (1,\ratio);
        %    \draw[help lines,purple] (\ratio,-1) -- (\ratio,1);
        % }
        % \draw (-1,0) -- (1,0);
        % \draw (0,-1) -- (0,1);
        \coordinate (v_liq) at (-0.77,-0.368);
        \coordinate (v_liqrk) at (-0.715,-0.28);
        \coordinate (v_vap) at (0.315,-0.368);
        \coordinate (v_vaprk) at (0.2,-0.28);
        \fill (v_liq) circle (2pt);
        \fill (v_liqrk) circle (2pt);
        \fill (v_vap) circle (2pt);
        \fill (v_vaprk) circle (2pt);

        \node[anchor=south,align=left] (DATA) at (-0.35,0.25) {Measured data:\\$\left\{\begin{aligned}&p_\text{sat}(12^\circ\text{C})=47.3\,\text{bar}\\
        &v_\text{liq}(12^\circ\text{C})=1.18\,\ell/\text{kg}\\
        &v_\text{vap}(12^\circ\text{C})=6.91\,\ell/\text{kg}
        \end{aligned}\right.$}; 

        \node[anchor=south,align=left] (RK) at (0.5,0.24) {Redlich-Kwong prediction:\\$\left\{\begin{aligned}&p_\text{sat}(12^\circ\text{C})=50.9\,\text{bar}\\
        &v_\text{liq}(12^\circ\text{C})=1.47\,\ell/\text{kg}\\
        &v_\text{vap}(12^\circ\text{C})=6.28\,\ell/\text{kg}
        \end{aligned}\right.$};
            
        \draw[-latex,thick,shorten >=5pt] (DATA) to[bend left=25] (v_liq);
        \draw[-latex,thick,shorten >=5pt] (DATA) to[bend right=25] (v_vap);

        \draw[-latex,thick,shorten >=5pt] (RK) to[bend left=15] (v_liqrk);
        \draw[-latex,thick,shorten >=5pt] (RK) to[bend left=10] (v_vaprk);

        \draw[purple, dashed] (v_liq)--++(0,-0.5);
        \draw[purple, dashed] (v_vap)--++(0,-0.5);
        \draw[purple, dashed] (v_liqrk)--++(0,-0.6);
        \draw[purple, dashed] (v_vaprk)--++(0,-0.6);
        
    \end{scope}
    \end{tikzpicture}
        
    \end{center}
\end{example}

Although the two-parameter models predict values for the saturation properties that do not exactly agree with experimental measurements, the historical importance of the models is significant. The understanding of substance properties has developed gradually over time, and even rough estimates of the pressure required to liquefy a gas are valuable when mapping the thermodynamic properties of a substance. 
With the help of two-parameter models, it is possible to estimate the parameters $a_s$ and $b_s$ from measurements performed in the single-phase gaseous region. These parameters can then be used to estimate the critical point and the pressure required to liquefy the gas at a given temperature. Such estimates can then be used to design the devices with which the experimental measurements are carried out.
\begin{example}
    Assuming a gas to obey the Redlich–Kwong equation of state, one may determine the model coefficients by measuring the pressure of the gas inside a rigid container at two different temperatures. Given state points $(p_1,v,T_1)$ and $(p_2,v,T_2)$, one finds
    \begin{align*}
        &\left.\begin{aligned}
        p_1&=\frac{R_sT_1}{v-b_s}-\frac{a_s}{\sqrt{T_1} v(v+b_s)}\\
        p_2&=\frac{R_sT_2}{v-b_s}-\frac{a_s}{\sqrt{T_2} v(v+b_s)}
        \end{aligned}\right\}\implies 
        \begin{aligned}
        b_s&=v-\frac{R_sT_1\sqrt{T_1}-R_sT_2\sqrt{T_2} }{p_1\sqrt{T_1}-p_2\sqrt{T_2}}\\
        a_s&=\frac{(p_1T_2-p_2T_1)\sqrt{T_2}\sqrt{T_1}}{T_1\sqrt{T_1}-T_2\sqrt{T_2}} v(v+b_s)
        \end{aligned}
    \end{align*}
    Once the values for $a_s$ and $b_s$ are known, the critical temperature and pressure predicted by the Redlich-Kwong model are obtained from
    \begin{align*}
        T_\text{crit}&=\left(\frac{3(2^{1/3}-1)^2a_s}{R_sb_s}\right)^{2/3}, &
        p_\text{crit}&=\frac{2^{1/3}-1}{3}\frac{R_sT_\text{crit}}{b_s}
    \end{align*}
    
    At a specific volume of $v=0.10\,\text{m}^3/\text{kg}$, the pressure of methane at $T_1=-20^\circ\text{C}$ has been measured at $p_1=12.61\,\text{bar}$. At the same volume at temperature $T_2=50^\circ\text{C}$, the pressure has been measured at $p_2=16.40\,\text{bar}$. The Redlich-Kwong model then predicts the critical temperature and pressure of methane to be $T_\text{crit}=-91.54^\circ\text{C}$ and $p_\text{crit}=40.50\,\text{bar}$. The estimates are reasonably close to the experimentally measured values of $T_\text{crit}=-82.59^\circ\text{C}$ and $p_\text{crit}=45.99\,\text{bar}$.
\end{example}

\section{Numerical libraries}
Engineering applications such as steam power plants and heat pumps involve not only the gaseous phase of substances but also the liquid phase and the liquid-vapor phase mixture. While the two-parameter models have historical value in predicting parameter regimes needed for designing experimental measurements, they lack the accuracy needed in the most demanding of the applications. Fortunately, the theoretical framework of thermodynamics enables constructing multiparameter substance models that can be fitted to experimental data with good accuracy. Several libraries and software implementations of such models exist. The ones introduced here are \verb|pyromat| and \verb|coolprop|, which both are open source, offer a user-friendly python interface, and have extensive documentation. Both libraries can be installed using \verb|pip|, the package installer for python.

To use \verb|pyromat|, one imports the library and reverts to SI-units, if desired. By default, \verb|pyromat| uses non-SI units for pressure (bar), molar quantities (kmol), and energy (kJ). In engineering applications, such units are often common, though. Then one instantiates a multiphase model for the desired substance and computes thermodynamic quantities by calling individual quantity methods such as \verb|p| (pressure), \verb|v| (specific volume), \verb|d| (density), \verb|T| (temperature), and \verb|x| (vapor quality) or by calling the \verb|state| method that evaluates all quantities from a given input quantity pair: 
\lstinputlisting[language=Python,firstline=1,firstnumber=1,lastline=22]{\rootfolder/code/example-pyromat-pvTx.py}
The library has separate methods for critical point and triple line information, and for evaluation of saturation properties given only temperature (or pressure):
\lstinputlisting[language=Python,firstline=24,firstnumber=24,lastline=34]{\rootfolder/code/example-pyromat-pvTx.py}
The acceptable input parameter pairs for the methods, that can be formed from the state variables $p$, $v$, $\rho$, $T$, and $x$, are listed in Table~\ref{tab:pyromat-pvTx-input-pairs} together with their applicability domains. The full documentation is available at \href{http://pyromat.org}{http://pyromat.org}.
\begin{table}[!htb]
    \centering
    \begin{threeparttable}
    \caption{Acceptable pyromat input parameter pairs formable from the variables $p$, $v$, $\rho$, $T$, and $x$.}
    \begin{tabular}{llcc}
    \toprule 
    Variable pair & Input pair & Single-phase & Two-phase\\
    \midrule 
    $(p, T)$ & \texttt{p} and \texttt{T}  & Yes & No \\
    $(\rho, T)$ & \texttt{d} and \texttt{T} & Yes & Yes \\
    $(v, T)$ & \texttt{v} and \texttt{T} & Yes & Yes \\
    $(\rho, p)$ & \texttt{d} and \texttt{p} & Yes & Yes \\
    $(v, p)$ & \texttt{v} and \texttt{p} & Yes & Yes \\
    $(x, T)$ & \texttt{x} and \texttt{T} & No & Yes \\
    $(x, p)$ & \texttt{x} and \texttt{p} & No & Yes \\
    \bottomrule
    \end{tabular}
    \end{threeparttable}
    \label{tab:pyromat-pvTx-input-pairs}
\end{table}

The \verb|pyromat|-library accepts input as scalars, as \verb|numpy| arrays, or as a mixture of the two. The output is always a \verb|numpy| array, even in case of scalar input. This enables convenient data queries for multiple states with one call but must be kept in mind if outputs from scalar input queries are used in further calculations that do not accept or accept \verb|numpy| arrays. The array inputs and outputs are particularly useful for visualization of the data.
\begin{example}
    With the help of \verb|pyromat|, selected isotherms and vapor quality curves can be drawn in a $p$-$v$ diagram for example in the following manner:
    \lstinputlisting[language=Python,firstline=36,firstnumber=1]{\rootfolder/code/example-pyromat-pvTx.py}
    Running the script results in the following figure, showing selected isotherms and vapor quality contours:
    \begin{center}
        \includegraphics[width=0.7\linewidth]{\rootfolder/figs/example-pyromat-pv.pdf}
    \end{center}
\end{example}

To use \verb|coolprop|, one imports both the library and a specific \verb|AbstractState| class from it. The SI-units will be used by default, needing no specific configuration. A multiphase model for the desired substance is then instantiated from the \verb|AbstractState| class. Before any thermodynamic quantities can be computed, though, the substance state has to be specified. This happens by calling the method \verb|update| and passing it not only the variables but also an input key that informs the library which variables have been passed. Once the state has been updated, quantity values are accessed through separate methods such as \verb|p| (pressure), \verb|T| (temperature), \verb|rhomass| (density), and \verb|Q| (vapor quality).
\lstinputlisting[language=Python,firstline=1,lastline=14,firstnumber=1]{\rootfolder/code/example-coolprop-pvTx.py}
Like \verb|pyromat|, also \verb|coolprop| has methods for critical point and triple line information, although with different syntax. However, \verb|coolprop| does not have separate methods for saturation properties. Instead, they are evaluated by pairing vapor quality and temperature as inputs.
\lstinputlisting[language=Python,firstline=16,firstnumber=16]{\rootfolder/code/example-coolprop-pvTx.py}
The acceptable input parameter pairs and their corresponding input keys for the \verb|update|, that can be formed from the state variables $p$, $v$, $\rho$, $T$, and $x$, are listed in Table~\ref{tab:coolprop-pvTx-input-pairs}. The full documentation is available at \href{https://coolprop.org}{https://coolprop.org}.
\begin{table}[!htb]
    \centering
    \begin{threeparttable}
    \caption{Acceptable coolprop input parameter pairs from the variables $p$, $v$, $\rho$, $T$, and $x$ and the required key. Applicability domain of the pairs.}
    \begin{tabular}{llcc}
    \toprule 
    Variable Pair & Input Key & Single-phase & Two-phase\\
    \midrule 
    $(p, T)$ & \texttt{PT\_INPUTS} & Yes & No \\
    $(\rho, T)$ & \texttt{DmassT\_INPUTS} & Yes & Yes \\
    $(\rho, p)$ & \texttt{DmassP\_INPUTS} & Yes & Yes \\
    $(x, T)$ & \texttt{QT\_INPUTS} & No & Yes \\
    $(p, x)$ & \texttt{PQ\_INPUTS} & No & Yes \\
    $(\rho, x)$ & \texttt{DmassQ\_INPUTS} & No & Yes \\
    \bottomrule
    \end{tabular}
    \end{threeparttable}
    \label{tab:coolprop-pvTx-input-pairs}
\end{table}

Unlike \verb|pyromat|, \verb|coolprop| does not support \verb|numpy| arrays. Hence all inputs and outputs are always scalar-valued. On the other hand, the \verb|coolprop| provides a large collection of different substances while \verb|pyromat| presently has only a few different substances implemented. Both libraries have their uses, and are best learned by practicing.

\section{Thermal expansion and compressibility coefficients}
The thermodynamic framework outlined in Fig.~\ref{fig:roadmap-thermo} contains not only $p$, $v$, and $T$, but also energy-related quantities such as internal energy $u$, enthalpy $h$, and entropy $s$ that are needed in practical calculations involving power plants and heat pumps. Unlike $p$, $v$, and $T$, the energy-related quantities cannot be obtained from direct measurements. Instead, their experimental determination relies on observing how the measurable variables $p$, $v$, and $T$ change relative to each other when a substance is heated or compressed. At a detailed level, the relative changes are quantified by derivatives.

In the single-phase regions, pressure, specific volume, and temperature relate to each other via the surface equation $f(p,v,T)=0$ and, at least locally, one may define the functions $p(v,T)$, $v(T,p)$, and $T(p,v)$. Calculating total differentials of these functions, one finds the following relations between the differential changes in pressure, specific volume, and temperature
\begin{align}
    dp&=\left(\frac{\partial p}{\partial v}\right)_Tdv+\left(\frac{\partial p}{\partial T}\right)_v dT,\\
    dv&=\left(\frac{\partial v}{\partial T}\right)_pdT+\left(\frac{\partial v}{\partial p}\right)_T dp,\\
    dT&=\left(\frac{\partial T}{\partial p}\right)_vdp+\left(\frac{\partial T}{\partial v}\right)_p dv.
\end{align}
As the differentials $dp$, $dv$, and $dT$ are the same on both sides of the equations, one may algebraically rearrange the equations and find relations between the partial derivatives. For example, one may manipulate the differential resulting from the function $v(T,p)$ into
\begin{align}
    dv&=\left(\frac{\partial v}{\partial T}\right)_pdT+\left(\frac{\partial v}{\partial p}\right)_T dp & &\implies & &\left\{\begin{aligned}
    dp&=\frac{1}{\left(\frac{\partial v}{\partial p}\right)_T}dv-\frac{\left(\frac{\partial v}{\partial T}\right)_p}{\left(\frac{\partial v}{\partial p}\right)_T}dT,\\
    dT&=\frac{1}{\left(\frac{\partial v}{\partial T}\right)_p}dv-\frac{\left(\frac{\partial v}{\partial p}\right)_T }{\left(\frac{\partial v}{\partial T}\right)_p}dp.\end{aligned}\right.
\end{align}
Comparing the manipulated expressions with the differentials resulting directly from differentiation of the functions $p(v,T)$ and $T(p,v)$, one identifies relations such as
\begin{align}
    \left(\frac{\partial p}{\partial T}\right)_v&=-\frac{\left(\frac{\partial v}{\partial T}\right)_p}{\left(\frac{\partial v}{\partial p}\right)_T}.
\end{align}
The mutual derivatives between $p$, $v$, and $T$ are therefore related to each other.

In the single-phase region, the $p$-$v$-$T$ surface has two independent variables, and all mutual derivatives can be expressed in terms of two independent quantities, namely the isobaric thermal expansion coefficient $\beta_p$ and the isothermal compressibility coefficient $\kappa_T$. The coefficients measure the relative change in substance volume (or density) when either temperature or pressure changes and the other is fixed, and are defined by
\begin{align}
    \beta_p&=\frac{1}{v}\left(\frac{\partial v}{\partial T}\right)_p=-\frac{1}{\rho}\left(\frac{\partial \rho}{\partial T}\right)_p,\\
    \kappa_T&=-\frac{1}{v}\left(\frac{\partial v}{\partial p}\right)_T=\frac{1}{\rho}\left(\frac{\partial \rho}{\partial p}\right)_T.
\end{align}
In the two-phase mixture region, pressure and temperature on the $p$-$v$-$T$ surface are not independent but satisfy the saturation relation $g(T,p)=0$. Consequently, the usual definitions of $\beta_p$ and $\kappa_T$ as independent single-phase response coefficients do not apply. The slope of the coexistence curve
\begin{align}
    \left(\frac{dp}{dT}\right)_{\text{sat}}
\end{align}
nevertheless characterises how pressure and temperature vary along the two-phase equilibrium states. In single-phase regions, the coefficient $\beta_p$ can be determined by measuring the volume of a substance at several temperatures while keeping the pressure constant, and the coefficient $\kappa_T$ by measuring volume at several pressures while temperature is fixed. They can therefore be determined from existing $p$-$v$ and $v$-$T$ diagrams by examining the slopes of the isotherms and isobars. For most substances $\beta_p>0$ and $\kappa_T>0$, meaning that volume typically increases with temperature at constant pressure and decreases with pressure at constant temperature.
\begin{example}
    For water at pressure of $225\,\text{bar}$ and temperature of $376^\circ\text{C}$, one determines the slopes of the corresponding isotherm and isobar from the $p$-$v$ and $v$-$T$ diagrams as below
    \begin{center}
    \begin{tikzpicture}
    \node[above right,inner sep = 0] (image) at (0,0) {
        \includegraphics[width=0.85\linewidth]{\rootfolder/figs/example-beta-kappa-pv.pdf}
    };
    \begin{scope}[relative to node=image]
        % \foreach \ratio in {-1,-0.9,...,1}{
        %    \draw[help lines,purple] (-1,\ratio) -- (1,\ratio);
        %    \draw[help lines,purple] (\ratio,-1) -- (\ratio,1);
        % }
        % \draw (-1,0) -- (1,0);
        % \draw (0,-1) -- (0,1);
        
        \node[anchor=center,blue] at (0,0.18) {$\Delta v\approx 1.9\,\ell/\text{kg}$};

        \node[anchor=center,blue,rotate=90] at (-0.40,0.4) {$\Delta p\approx-4.0\,\text{bar}$};
      
        \coordinate (pvT) at (0.067,0.42);
        \draw[latex-, thick,shorten <=3pt,color=blue] (pvT) to[bend left=25] ++(0.1,0.2) node[anchor=west] {$\left\{\begin{aligned}
            p&=225\,\text{bar}\\
            T&=376^\circ\text{C}\\
            v&\approx 4.0\,\ell/\text{kg}
        \end{aligned}\right.$};
        
    \end{scope}
    \end{tikzpicture}    
    \end{center}

    \begin{center}
    \begin{tikzpicture}
    \node[above right,inner sep = 0] (image) at (0,0) {
        \includegraphics[width=0.85\linewidth]{\rootfolder/figs/example-beta-kappa-vT.pdf}
    };
    \begin{scope}[relative to node=image]
        % \foreach \ratio in {-1,-0.9,...,1}{
        %    \draw[help lines,purple] (-1,\ratio) -- (1,\ratio);
        %    \draw[help lines,purple] (\ratio,-1) -- (\ratio,1);
        % }
        % \draw (-1,0) -- (1,0);
        % \draw (0,-1) -- (0,1);
        
        \node[anchor=center,blue] at (0.25,-0.47) {$\Delta T\approx 2.0^\circ\text{C}$};

        \node[anchor=center,blue,rotate=90] at (0.5,0.0) {$\Delta v\approx2.1\,\ell/\text{kg}$};
      
        \coordinate (pvT) at (0.26,0.03);
        \draw[latex-, thick,shorten <=3pt,color=blue] (pvT) to[bend right=25] ++(-0.2,0.2) node[anchor=east] {$\left.\begin{aligned}
            p&=225\,\text{bar}\\
            T&=376^\circ\text{C}\\
            v&\approx 4.0\,\ell/\text{kg}
        \end{aligned}\right\}$};
        
        
    \end{scope}
    \end{tikzpicture}    
    \end{center}

    The isobaric thermal expansion coefficient and the isothermal compressibility coefficient at the specified state are then calculated to be
    \begin{align*}
        \beta_p\approx\frac{1}{v}\frac{\Delta v}{\Delta T}=0.263\,\text{K}^{-1}\\
        \kappa_T\approx-\frac{1}{v}\frac{\Delta v}{\Delta p}=0.119\,\text{bar}^{-1}
    \end{align*}
\end{example}
\begin{example}
    The isobaric thermal expansion coefficient $\beta_p$ and the isothermal compressibility coefficient $\kappa_T$ can be computed also numerically using the thermodynamic data libraries.         
    Using \verb|pyromat|, one can approximate derivatives with a finite difference rule $df/dx\approx (f(x+\delta x)-f(x))/\delta x$ using a small enough step size:
    \lstinputlisting[language=Python,firstline=1,lastline=28,firstnumber=1]{\rootfolder/code/example-beta-kappa-numerical.py}
    Running the script results in
    \begin{Verbatim}
beta_p = 0.261266 (1/K)
kappa_T = 0.116053 (1/bar)
    \end{Verbatim}
    Using \verb|coolprop|, the necessary derivatives can be accessed directly with method calls:
    \lstinputlisting[language=Python,firstline=30,firstnumber=1]{\rootfolder/code/example-beta-kappa-numerical.py}
    Running the script results in 
    \begin{Verbatim}
beta_p = 0.262058 (1/K)
kappa_T = 0.11624 (1/bar)
    \end{Verbatim}
    Both libraries return values close to each other and close to the values obtained visually from the $p$-$v$ and $v$-$T$ diagrams.
\end{example}

\section{Compressibility factor and reduced variables}
While numerical libraries provide high accuracy, the ideal gas model remains a convenient tool for rapid engineering estimates and analytical derivations due to its mathematical simplicity. To determine when the ideal gas model is ``good enough'', the so-called compressibility factor 
\begin{align}
Z = \frac{pv}{R_sT},
\end{align}
can be used as a dimensionless metric to quantify deviation from ideal behaviour. For an ideal gas, $Z=1$ exactly. As the name suggests, the factor compares the volume $v$ the substance occupies in relation to the volume $R_sT/p$ an ideal gas at the same temperature and pressure would occupy. If the value deviates from unity by more than a tolerable threshold (typically 1--5\%), then ideal gas model likely should be abandoned in favor of the more complex equations of state or numerical libraries. 

Diagrams involving the compressibility factor are usually expressed in terms of the so-called reduced variables that are dimensionless and normalized to the critical state values
\begin{align}
    T_r&=\frac{T}{T_\text{crit}},\\
    p_r&=\frac{p}{p_\text{crit}},\\
    v_r&=\frac{v}{v_\text{crit}}.
\end{align}
When illustrated for different substance, as in Fig.~\ref{fig:compressibility-factor} for methane $\text{CH}_4$, nitrogen $\text{N}_2$, and carbon dioxide $\text{CO}_2$, the compressibility factor reveals approximately similar behaviour for many substances, even with significantly different critical states. Even the phase transitions across the liquid-vapor two-phase state, visible on the diagram as strictly vertical segments, seem similar for different substances. This observation is referred to as the principle of corresponding states and it tells that many fluids have approximately similar compressibility factors and deviate from ideal to roughly same degree when compared at the same reduced temperature and pressure. It is for this reason that the Van der Waals, Dieterici, and Redlich-Kwong equations of state can be applied to many substances with relatively good accuracy.
\begin{example}
    In terms of the reduced temperature and pressure, the Van der Waals equation of state becomes
    \begin{align*}
        p_r=\frac{8T_r}{3v_r-1}-\frac{3}{v_r^2},
    \end{align*}    
    where the reduced volume variable
    \begin{align*}
        v_r=\frac{v}{v(T_\text{crit},p_\text{crit})}=\frac{8p_\text{crit}v}{3R_sT_\text{crit}},
    \end{align*}
    is the specific volume normalized to the model-predicted critical-state specific volume. In terms of reduced variables, the Van der Waals equation is seen to contain no substance-dependent parameters.
\end{example}
\begin{figure}[!htb]
    \centering
    \includegraphics[width=0.8\linewidth]{\rootfolder/figs/fig-compressibility-factor.pdf}
    \caption{Selected reduced temperature ($T_r$) isotherms of common gases illustrating the similarity of $Z$ of the gases.}
    \label{fig:compressibility-factor}
\end{figure}


\section{Summary}
Pure substances, such as oxygen, carbon dioxide, and distilled water,
exhibit qualitatively similar phase behaviour despite having widely
different critical temperatures and pressures. Their thermodynamic states
can be mapped experimentally in $p$-$v$-$T$ space and represented in
projections such as $p$-$v$ and $v$-$T$ diagrams. Simple two-parameter
equations of state, such as the Van der Waals model, can reproduce the
main features of this behaviour, while multi-parameter equations of
state implemented in numerical libraries such as \verb|pyromat| and
\verb|coolprop| can provide much more accurate property data.

For two-parameter equations of state, the model parameters can be related
to the critical state by requiring the isotherm to have a stationary
inflection point at the critical point,
\begin{align}
    \left(\frac{\partial p}{\partial v}\right)_{T=T_c,v=v(p_c,T_c)} &= 0,
    &
    \left(\frac{\partial^2 p}{\partial v^2}\right)_{T=T_c,v=v(p_c,T_c)} &= 0.
\end{align}
Measured critical-state data can then be used to determine the model
parameters, or the parameters can be estimated from measurements away
from the critical point. Together with the Maxwell construction, the
resulting model can also be used to estimate saturation states. The
approximately universal behaviour of different substances can be
illustrated using reduced variables and the compressibility factor
\begin{align}
    Z = \frac{v}{v_{\mathrm{ideal}}}=\frac{pv}{R_sT},
\end{align}
which reveals similar behaviour for different fluids at corresponding
reduced temperatures and pressures.

The thermodynamic state can also be described locally through differential
response coefficients. The isobaric thermal expansion coefficient
$\beta_p$ and the isothermal compressibility $\kappa_T$ characterize the
response of specific volume to changes in temperature and pressure,
respectively. Both can be estimated from local slopes in $p$-$v$ and
$v$-$T$ diagrams or calculated directly from property data. These
coefficients will later provide the connection between experimentally
measured thermodynamic response and the differential relations needed to
construct the thermodynamic framework.

\chapter{Mechanics of fluids}
The equations of motion describing the evolution of simple fluids, as those often encountered in engineering applications, are derived here. The derivation process involves first introducing the Leibniz rule for differentiating integrals, and then applying it to the conservation laws of mass, momentum, and energy. The process is mathematically more involved than what is typically encountered in thermodynamics courses, and is not expected to be digested fully during first reading. Instead, the purpose of this chapter is to build a detailed and rigorous enough of a foundation that can be later recalled and leveraged together with thermodynamic relations to analyze common engineering applications. 

\section{Learning goals}
\begin{itemize}
    \item[\elementary] The student conceptually understands the idea of a flow map.
    \item[\elementary] The student recalls how differential line-, surface-, and volume elements can be constructed in general coordinate systems.
    \item[\elementary] The student classifies whether a differential volume grows or shrinks when it is dragged along the flow of a given vector field.
    \item[\elementary] The student provides physical interpretation for the different terms in the Leibniz integral differentiation rule.
    \item[\elementary] The student associates a physical principle to each of the fluid's equations of motion.
    \item[\elementary] The student provides physical interpretation for the different terms in the fluid's equations of motion.
    \item[\intermediate] The student uses python to numerically integrate vector-valued ordinary differential equations of the form $d\bm{x}(t)/dt=\bm{v}(\bm{x}(t),t)$ to find out how points move along vector fields.
    \item[\intermediate] The student uses rank-2 tensors to represent linear relations between vectors and manipulates their basic tensor operations.
    \item[\intermediate] Given a velocity field, the student categorizes flow into isotropic deformation, shearing, and rotation.
    \item[\intermediate] The student recognizes the fluid equations of motion to be incomplete without a closing expression for internal energy and an equation of state.
    \item[\advanced] The student manipulates equations expressed over a volume and its surface to find local equations.
    \item[\advanced] The student confirms that the conservation of momentum implies that the traction vector is odd and linear with respect to the surface normal unit vector, leading to the concept of stress tensor.
    \item[\advanced] The student confirms that the conservation of angular momentum implies the stress tensor to be symmetric.
\end{itemize}


\section{The flow map of a vector field}
Assume the vector field $\bm{v}(\bm{x},t)$ to describe which direction fluid is moving at point $\bm{x}$ at time $t$, and at what speed. Drop a small grain of pollen into the fluid at $\bm{y}_0$ at some chosen time $t_0$. From that moment onwards, the position $\bm{y}(t)$ of the pollen grain will follow the fluid flow and obey the vector-valued ordinary differential equation
\begin{align}\label{eq:point-ODE}
    \frac{d\bm{y}(t)}{dt}=\bm{v}(\bm{y}(t),t), \qquad \bm{y}(t_0)=\bm{y}_0.
\end{align}
Over time, the position $\bm{y}(t)$ draws a curve into space with the tangent vectors of the curve denoting the direction of the local and instantaneous velocity of the grain.
\begin{center}
    \input{\rootfolder/code/fig-point-movement}
\end{center}

Solution to the differential equation~\eqref{eq:point-ODE} describing the motion of the pollen grain can be formally expressed with a time-dependent flow map $\bm{\phi}_t(\bm{x})$ associated with the velocity field $\bm{v}(\bm{x},t)$. Starting at time $t_0$, the map moves any point $\bm{x}$, including the $\bm{y}_0$ where the pollen grain was initially dropped, into a new point $\bm{\phi}_t(\bm{x})$ by integrating along the flow of the time-dependent velocity field $\bm{v}(\bm{x},t)$ until time $t$ is reached. The derivative of the flow map with respect to time $t$ therefore corresponds to the velocity field evaluated at the inspected point and time
\begin{align}
    \frac{d}{dt}\bm{\phi}_t(\bm{x})=\bm{v}(\bm{\phi}_t(\bm{x}),t),
\end{align}
and, at time $t=t_0$, the map reduces to identity $\bm{\phi}_{t_0}(\bm{x})=\bm{x}$. Under the action of the flow map $\bm{\phi}_t$, any set of points $\{\bm{x}|\bm{x}\in\Omega(t_0)\}$ is mapped into an image set $\Omega(t)=\text{Im}_{\bm{\phi}_t}(\Omega(t_0))=\{\bm{\phi}_t(\bm{x})|\bm{x}\in\Omega(t_0)\}$ of the original one as illustrated in Fig.~\ref{fig:flowmap}. The map therefore moves not only points but lines, surfaces, and volumes too along the flow.
\begin{figure}[!htb]
    \centering
    \input{\rootfolder/code/fig-flowmap}
    \caption{Movement of a general point $\bm{x}$ and the image $\text{Im}_{\bm{\phi}_t}(\Omega(t_0))$ of a domain $\Omega(t_0)$ under the flow map $\bm{\phi}_t(\bm{x})$ of a vector field $\bm{v}(\bm{x},t)$. The pollen grain is denoted with~$\bm{y}$.}
    \label{fig:flowmap}
\end{figure}

\begin{example}
    The vector field $\bm{v}(\bm{x},t)=\omega(y\bm{\hat{x}}-x\bm{\hat{y}})+v_z\bm{\hat{z}}$, with $\omega$ and $v_z$ constants, has a flow map
    \begin{align*}
        \bm{\phi}_t(\bm{x})=\cos(\omega(t-t_0))(x\bm{\hat{x}}+y\bm{\hat{y}})+\sin(\omega(t-t_0))(y\bm{\hat{x}}-x\bm{\hat{y}})+(z+v_z(t-t_0))\bm{\hat{z}}
    \end{align*}
    The map describes a helical trajectory with constant angular velocity $\omega$ and translative velocity $v_z$.
\end{example}

\begin{example}
    The flow map of a vector field can be calculated also numerically, solving the corresponding ordinary differential equation for a selected set of points. This can be performed with script as below for the vector field of the previous example
    \lstinputlisting[language=Python,firstline=1,lastline=22,firstnumber=1]{\rootfolder/code/example-odeint.py}
    When the points \verb|xs| are illustrated, one observes a helical trajectory expanding into $\bm{\hat{z}}$ direction starting from the red dot, agreeing with the analytical solution.
    \begin{center}
        \includegraphics[width=0.6\linewidth]{\rootfolder/figs/example-odeint.pdf}
    \end{center}
\end{example}

If an infinitesimally short line, small patch, or volume of pollen is dropped into the flow instead of a single point-like grain, how exactly will they develop over time? Points are known to move according to the differential equation \eqref{eq:point-ODE} but infinitesimally small lines, surface patches, and volumes differ from points as mathematical objects. What equations will their evolution follow and what are the rates at which these objects change in time? Answering these questions enables one to derive the Leibniz integral differentiation rule and to apply the principles of mass, momentum, and energy conservation to ultimately derive equations that describe the evolution of the fluid's velocity field $\bm{v}(\bm{x},t)$, density $\rho(\bm{x},t)$, pressure $p(\bm{x},t)$, and temperature $T(\bm{x},t)$. Before setting the differential objects into motion, one has to understand how they are constructed in the first place.

\section{Differential line-, surface, and volume elements}
Consider a coordinate triplet $(u_1,u_2,u_3)$ that defines a point $\bm{x}$ in space. In other words, a location in space is a vector-valued function $\bm{x}(u_1,u_2,u_3)$ of the coordinates. A differential, vector-valued line element $d\bm{\ell}_{\bm{x},\alpha}$, into the direction of the coordinate $u_\alpha$ and rooted at $\bm{x}$, is generated by the change in the position $\bm{x}$ with respect to a differential change $du_\alpha$ in the coordinate $u_\alpha$ via the definition
\begin{align}
    d\bm{\ell}_{\bm{x},\alpha}=\lim_{\epsilon\to0}\frac{\bm{x}(u_\alpha+\epsilon du_\alpha,u_\beta,u_\gamma)-\bm{x}(u_\alpha,u_\beta,u_\gamma)}{\epsilon}=\frac{\partial\bm{x}}{\partial u_\alpha}du_\alpha.
\end{align}
At point $\bm{x}$, the differential line element $d\bm{\ell}_{\bm{x},\alpha}$ has a differential scalar-valued length 
\begin{align}
    d\ell_{\bm{x},\alpha}=\left|d\bm{\ell}_{\bm{x},\alpha}\right|=\left|\frac{\partial\bm{x}}{\partial u_\alpha}\right|du_\alpha,
\end{align}
and a non-differential unit vector pointing its direction
\begin{align}
    \bm{\hat{\ell}}_{\bm{x},\alpha}=\frac{\partial \bm{x}/\partial u_\alpha}{\left|\partial\bm{x}/\partial u_\alpha\right|}.
\end{align}
In terms of its length and the unit direction vector, the differential line element can be written equally as
\begin{align}
    d\bm{\ell}_{\bm{x},\alpha}=\bm{\hat{\ell}}_{\bm{x},\alpha}d\ell_{\bm{x},\alpha}.
\end{align}

A differential, vector-valued, oriented surface element $d\bm{A}_{\bm{x}}=\bm{\hat{n}}_{\bm{x}}dA_{\bm{x}}$, rooted at $\bm{x}$ and having a surface unit normal vector $\bm{\hat{n}}_{\bm{x}}$, can be constructed by considering the oriented area two different differential line elements span. For example, the vector-valued line elements $d\bm{\ell}_{\bm{x},\alpha}$ and $d\bm{\ell}_{\bm{x},\beta}$ create a vector valued oriented surface element $d\bm{A}_{\bm{x},\gamma}$ via the cross-product
\begin{align}
    d\bm{A}_{\bm{x},\gamma}&=d\bm{\ell}_{\bm{x},\alpha}\times d\bm{\ell}_{\bm{x},\beta}=\frac{\partial \bm{x}}{\partial u_\alpha}\times\frac{\partial \bm{x}}{\partial u_\beta}du_\alpha du_\beta.
\end{align}
At point $\bm{x}$, such element has a non-negative, scalar-valued but differential area
\begin{align}
    dA_{\bm{x},\gamma}=\left|d\bm{A}_{\bm{x},\gamma}\right|=\left|\frac{\partial \bm{x}}{\partial u_\alpha}\times\frac{\partial \bm{x}}{\partial u_\beta}\right|du_\alpha du_\beta,
\end{align}
and a non-differential surface unit normal vector
\begin{align}
    \bm{\hat{n}}_{\bm{x},\gamma}=\frac{\frac{\partial \bm{x}}{\partial u_\alpha}\times\frac{\partial \bm{x}}{\partial u_\beta}}{\left|\frac{\partial \bm{x}}{\partial u_\alpha}\times\frac{\partial \bm{x}}{\partial u_\beta}\right|}.
\end{align}

Finally, a differential volume element $dV_{\bm{x}}$, rooted at $\bm{x}$, can be constructed by considering the volume all three differential line elements span
\begin{align}
    dV_{\bm{x}}&=\left|d\bm{\ell}_{\bm{x},\alpha}\times d\bm{\ell}_{\bm{x},\beta}\cdot d\bm{\ell}_{\bm{x},\gamma}\right|=|d\bm{A}_{\bm{x},\gamma}\cdot d\bm{\ell}_{\bm{x},\gamma}|=\left|\frac{\partial \bm{x}}{\partial u_\alpha}\times\frac{\partial \bm{x}}{\partial u_\beta}\cdot\frac{\partial \bm{x}}{\partial u_\gamma}\right|du_\alpha du_\beta du_\gamma.
\end{align}
The absolute value of the scalar triple product of the three directional vectors define the Jacobian $J$ of the coordinate system
\begin{align}
    J=\left|\frac{\partial \bm{x}}{\partial u_\alpha}\times\frac{\partial \bm{x}}{\partial u_\beta}\cdot\frac{\partial \bm{x}}{\partial u_\gamma}\right|=\left|\text{det}\begin{pmatrix}
        \frac{\partial x_1}{\partial u_1} & \frac{\partial x_1}{\partial u_2} & \frac{\partial x_1}{\partial u_3} \\
        \frac{\partial x_2}{\partial u_1} & \frac{\partial x_2}{\partial u_2} & \frac{\partial x_2}{\partial u_3} \\
        \frac{\partial x_3}{\partial u_1} & \frac{\partial x_3}{\partial u_2} & \frac{\partial x_3}{\partial u_3}
    \end{pmatrix}\right|,
\end{align}
that on mathematics courses is often introduced as the determinant of the matrix containing the partial derivatives of the coordinates, instead of the geometric view presented here. Both definitions give the same result.
Example coordinate-triplet isolines trough a point $\bm{x}$ and the possible vectors $\partial\bm{x}/\partial u_\alpha$ they generate at the point $\bm{x}$ are illustrated in Fig.~\ref{fig:differential-elements}. From the figure, one may observe the vectors to span finite areas and a finite volume. The areas and the volume become differential once the lenght along the vectors are scaled with the associated differentials $du_\alpha$.
\begin{figure}[!htb]
    \centering
    \input{\rootfolder/code/fig-coordinates}
    \caption{Coordinate lines and vectors that define differential line, surface, and volume elements.}
    \label{fig:differential-elements}
\end{figure}
\begin{example}
    A point $\bm{x}$ in 3-D space can be described by the Cartesian coordinate triplet $(x,y,z)$ as
    \begin{align*}
        \bm{x}=x\bm{\hat{x}}+y\bm{\hat{y}}+z\bm{\hat{z}}.
    \end{align*}
    The same point can be described by the cylindrical coordinate triplet $(r,\vartheta,z)$ as 
    \begin{align*}
        \bm{x}=r\cos\vartheta\bm{\hat{x}}+r\sin\vartheta\bm{\hat{y}}+z\bm{\hat{z}},
    \end{align*}
    with the coordinates $x$ and $y$ being determined from the cylindrical coordinates according to $x=r\cos\vartheta$ and $y=r\sin\vartheta$.
\end{example}
\begin{example}
    In cartesian coordinates $(x,y,z)$, the differential line elements at point $\bm{x}$ are
    \begin{align*}
        d\bm{\ell}_{\bm{x},x}=\bm{\hat{x}}dx, \quad d\bm{\ell}_{\bm{x},y}=\bm{\hat{y}}dy, \quad d\bm{\ell}_{\bm{x},z}=\bm{\hat{z}}dz.
    \end{align*}
    At the same point, they generate differential oriented surface elements such as
    \begin{align*}
        d\bm{A}_{\bm{x},x}=d\bm{\ell}_{\bm{x},y}\times d\bm{\ell}_{\bm{x},z}=\bm{\hat{y}}\times\bm{\hat{z}}dydz=\bm{\hat{x}}dydz,
    \end{align*}
    and a differential volume element
    \begin{align*}
        dV_{\bm{x}}=|\bm{\hat{y}}\times\bm{\hat{z}}\cdot\bm{\hat{x}}|dydzdx=dxdydz.
    \end{align*}
    In cylindrical coordinates $(r,\vartheta,z)$, the differential line elements at point $\bm{x}$ are computed to be
    \begin{align*}
        d\bm{\ell}_{\bm{x},r}&=(\cos\vartheta\bm{\hat{x}}+\sin\vartheta\bm{\hat{y}})dr,\\
        d\bm{\ell}_{\bm{x},\vartheta}&=(-r\sin\vartheta\bm{\hat{x}}+r\cos\vartheta\bm{\hat{y}})d\vartheta,\\
        d\bm{\ell}_{\bm{x},z}&=\bm{\hat{z}}dz.
    \end{align*}
    They generate differential oriented surface elements such as
    \begin{align*}
        d\bm{A}_{\bm{x},z}=d\bm{\ell}_{\bm{x},r}\times d\bm{\ell}_{\bm{x},\vartheta}=(\cos\vartheta\bm{\hat{x}}+\sin\vartheta\bm{\hat{y}})\times (-r\sin\vartheta\bm{\hat{x}}+r\cos\vartheta\bm{\hat{y}})drd\vartheta=\bm{\hat{z}}rdrd\vartheta,
    \end{align*}
    and a differential volume element 
    \begin{align*}
        dV_{\bm{x}}=|d\bm{A}_{\bm{x},z}\cdot d\bm{\ell}_{\bm{x},z}|=rdrd\vartheta dz.
    \end{align*}
\end{example}

\section{Rank-2 tensors}
Similarly as the initial position of a pollen grain dropped into a flowing fluid is mapped into a new position over time, also the vectors $\partial\bm{x}/\partial u_\alpha$ that generate differential line, surface, and volume elements evolve if pushed forward in time by the flow map $\bm{\phi}_t$, resulting in changes in the vector orientation and length as illustrated in Fig.~\ref{fig:differential-elements-movement}. This change in the vectors will then, in turn, compound into changes in the differential objects the vectors generate. In calculating the time rate of change of the vectors and the related objects, one encounters often for the first time an expression that can be interpreted as a tensor. As tensors are encountered also in converting physical conservation principles into local differential equations for the fluid, but rarely covered in elementary calculus classes, it is meaningful to address them at least somewhat before moving forward. 
\begin{figure}[!htb]
    \centering
    \input{\rootfolder/code/fig-coordinates-movement}
    \caption{Movement of the vectors generating differential line elements.}
    \label{fig:differential-elements-movement}
\end{figure}

The rate of change of the vectors generating the differential elements can be computed along the flow using the chain rule for differentiation and the definition of the flow map itself. Direct computation then provides 
\begin{align}
    \frac{d}{dt}\frac{\partial\bm{\phi}_t(\bm{x})}{\partial u_\alpha}=\frac{\partial}{\partial u_\alpha}\frac{d\bm{\phi}_t(\bm{x})}{dt}=\frac{\partial \bm{v}(\bm{\phi}_t(\bm{x}),t)}{\partial u_\alpha}=\frac{\partial\bm{\phi}_t(\bm{x})}{\partial u_\alpha}\cdot\nabla\bm{v}\Big|_{(\bm{\phi}_t(\bm{x}),t)}
\end{align}
The expression on the right can be interpreted in two different ways. First, the dot product of the vector $\partial\bm{\phi}_t(\bm{x})/\partial u_\alpha$ and the vector-valued differential operator $\nabla$ can be viewed as a scalar valued differential operator that acts on the vector-valued function $\bm{v}$ at point the point $\bm{\phi}_t(\bm{x})$ at time $t$, resulting in a new vector. The alternative is to interpret the expression as a left dot-product of a rank-2 tensor $\nabla\bm{v}$, evaluated at $\bm{\phi}_t(\bm{x})$ at time $t$, with the vector $\partial \bm{\phi}_t(\bm{x})/\partial u_\alpha$.
In Cartesian coordinates, the velocity vector at point $\bm{x}$ at time $t$ is $\bm{v}(\bm{x},t)=\sum_{i=1}v_i(\bm{x},t)\bm{\hat{x}}_i$ and the tensor $\nabla\bm{v}$ at the same point and time has a Cartesian representation
\begin{align}
    \nabla\bm{v}\Big|_{(\bm{x},t)}=\sum_{ij}\bm{\hat{x}}_i\otimes\bm{\hat{x}}_j\frac{\partial v_j(\bm{x},t)}{\partial x_i},
\end{align}
with $\otimes$ denoting the tensor product that comes with specific rules (listed later). In Cartesian coordinates, the vector on the left is $\partial\bm{\phi}_t(\bm{x})/\partial u_\alpha=\sum_{k=1}^3(\partial\phi^k_t(\bm{x})/\partial u_\alpha)\bm{\hat{x}}_k$, and the dot product with the tensor $\nabla\bm{v}$ at $(\bm{\phi}_t(\bm{x}),t)$ from the left then provides
\begin{align}
    \frac{\partial\bm{\phi}_t(\bm{x})}{\partial u_\alpha}\cdot\nabla\bm{v}\Big|_{(\bm{\phi}_t(\bm{x}),t)}&=\left(\sum_{k=1}^3\frac{\partial \phi^k_t(\bm{x})}{\partial u_\alpha}\bm{\hat{x}}_k\right)\cdot\left(\sum_{ij}\bm{\hat{x}}_i\otimes\bm{\hat{x}}_j\frac{\partial v_j}{\partial x_i}\Big|_{(\bm{\phi}_t(\bm{x}),t)}\right)\nonumber\\
    &=\sum_{i,j,k=1}^3\frac{\partial\phi^k_t(\bm{x})}{\partial u_\alpha}\frac{\partial v_j}{\partial x_i}\Big|_{(\bm{\phi}_t(\bm{x}),t)} \underbrace{\bm{\hat{x}}_k\cdot\bm{\hat{x}}_i}_{=\delta_{ik}}\otimes\,\bm{\hat{x}}_j\nonumber\\
    &=\sum_{i,j=1}^{3}\frac{\partial \phi^i_t(\bm{x})}{\partial u_\alpha}\frac{\partial v_j}{\partial x_i}\Big|_{(\bm{\phi}_t(\bm{x}),t)}\bm{\hat{x}}_j.
\end{align}
The result is a vector and could be obtained also by applying the scalar valued differential operator $\sum_{i=1}^3(\partial\phi^i_t(\bm{x})/\partial u_\alpha)(\partial/\partial x_i)$ on the velocity vector $\bm{v}$ at the point $\bm{\phi}_t(\bm{x})$. Identifying and handling $\nabla\bm{v}$ as a tensor object, however, reveals something fundamental. The expression 
\begin{align}
    \frac{d}{dt}\frac{\partial\bm{\phi}_t(\bm{x})}{\partial u_\alpha}=\frac{\partial\bm{\phi}_t(\bm{x})}{\partial u_\alpha}\cdot\nabla\bm{v}\Big|_{(\bm{\phi}_t(\bm{x}),t)}
\end{align}
tells that the time rate of change of the vector $\partial\bm{\phi}_t(\bm{x})/\partial u_\alpha$ is linearly related to the vector itself, with the tensor $\nabla \bm{v}$ at the point $\bm{\phi}_t(\bm{x})$ acting as a kind of non-constant, non-scalar proportionality factor, generating the temporal change. It will be demonstrated later that the tensor $\nabla\bm{v}$ in fact encodes different types of fluid motion.

A general rank-2 tensor $\bm{\sigma}$ in Cartesian coordinates has a representation
\begin{align}
    \bm{\sigma}=\sum_{i,j=1}^{3}\bm{\hat{x}}_i\otimes\bm{\hat{x}}_j\sigma_{ij}
\end{align}
with $\sigma_{ij}$ being a two-dimensional list of real numbers representing the coefficient values of the tensor in the specified coordinate system, in this case the Cartesian coordinates. The list of values can be accessed by dotting the tensor with the associated coordinate basis vector from left and right
\begin{align}
    \sigma_{ij}=\bm{\hat{x}}_i\cdot\bm{\sigma}\cdot\bm{\hat{x}}_j.
\end{align}
While the two-dimensional list $\{\sigma_{ij}\}_{i,j=1}^3$ can be seen as a matrix, the tensor $\bm{\sigma}$ contains not only the coefficients but also the basis vectors and their tensor product. A tensor therefore should never be misunderstood as just a matrix. If some other coordinate basis is used instead of the Cartesian one, the coefficient values change accordingly, due to changes in the basis vectors. A rank-2 tensor of particular importance is the identity tensor
\begin{align}
    \bm{I}=\sum_{ij}\bm{\hat{x}}_i\otimes\bm{\hat{x}}_j\delta_{ij}.
\end{align}
Below, the most common rules involving dot product and tensor product are listed. Here $\bm{a}$, $\bm{b}$, and $\bm{c}$ are vectors while $\bm{\sigma}$ and $\bm{\pi}$ are rank-2 tensors, and $\bm{I}$ is the identity tensor.
\begin{align}
(\bm{a} \otimes \bm{b}) \cdot \bm{c} &= \bm{a}(\bm{b} \cdot \bm{c}) && \text{Right dot product} \\
\bm{c} \cdot (\bm{a} \otimes \bm{b}) &= (\bm{c} \cdot \bm{a})\bm{b} && \text{Left dot product} \\
\bm{\sigma} : \bm{\pi} &= \sum_{i,j} \sigma_{ij} \pi_{ij} && \text{Double dot product} \\
\bm{I} : \bm{\sigma}&= \sum_{i}\sigma_{ii} && \text{Trace operation} \\
\bm{\sigma}\cdot\bm{a}&= \bm{a}\cdot\bm{\sigma}^T  && \text{Tensor transpose}
\end{align}
\begin{example}
    Consider a point $\bm{x}(\theta)=\cos\theta\bm{\hat{x}}+\sin\theta\bm{\hat{y}}$ on a unit circle in $x$-$y$ plane, parametrized by the angle $\theta$. The motion of the point $\bm{x}(\theta)$ along the circle is described by the vector-valued ordinary differential equation
    \begin{align}
        \frac{d\bm{x}}{d\theta}=-\sin\theta\bm{\hat{x}}+\cos\theta\bm{\hat{y}}
    \end{align}
    In terms of an antisymmetric rank-2 tensor 
    \begin{align}
        \bm{\Omega}=-\bm{\hat{x}}\otimes\bm{\hat{y}}+\bm{\hat{y}}\otimes\bm{\hat{x}}=\sum_{i,j=1}^2\bm{\hat{x}}_i\otimes\bm{\hat{x}}_j\Omega_{ij}
    \end{align}
    which, in the Cartesian coordinates, has the matrix components
    \begin{align}
        \Omega_{ij}=\begin{pmatrix}
            0 & -1\\ 1 & 0
        \end{pmatrix}
    \end{align}
    the differential equation can be written also as
    \begin{align}
        \frac{d\bm{x}}{d\theta}=\bm{\Omega}\cdot\bm{x}=\bm{x}\cdot\bm{\Omega}^T.
    \end{align}
    Here the tensor $\bm{\Omega}$ acts as the generator of rotation along the unit circle.
\end{example}

In addition to serving as generators of motion, tensors can be used also to represent linear relations in higher dimensions in a coordinate independent manner. In a single dimension, a linear function is written as $f(x)=kx$. In multiple dimensions a linear map can be written as $\bm{f}(\bm{x})=\bm{\sigma}\cdot\bm{x}$, with $\bm{\sigma}$ a tensor. In both fluid and solid mechanics, this linear representation is used to present contact-force density on a surface with respect to the surface normal vector.
\begin{example}
    \label{ex:tensor_linear_relation}
    In both fluid and solid mechanics, one will study how forces acting on a continuous body affect the body. A contact force is distributed on a surface $S$ and represented as a surface force density. The differential contact force $\delta\bm{F}$ on a differential surface area $dA$ is expressed as
    \begin{align}
        \delta\bm{F}=\bm{t}dA
    \end{align}
    where $\bm{t}$ is the surface force density, called traction. The total force on the surface is then the integral
    \begin{align*}
        \bm{F}=\int_{S}\bm{t}dA.
    \end{align*}
    The surface $S$ may cover the entire body under inspection, meaning $S=\partial V$ is the boundary of the inspected body volume $V$. If the surface force density can be written as a linear relation 
    \begin{align*}
        \bm{t}=\bm{\sigma}\cdot\bm{\hat{n}}
    \end{align*}
    with respect to the surface unit normal $\bm{\hat{n}}$ and some rank-2 tensor $\bm{\sigma}$, the integral over the entire body surface can be converted into a volume integral of tensor divergence
    \begin{align*}
        \int_{\partial V}\bm{t}dA=\int_{\partial V}\bm{\sigma}\cdot\bm{\hat{n}}dA=\int_{V}\nabla\cdot\bm{\sigma}dV.
    \end{align*}
    In mechanics, inspection of Newton's second law will provide this condition.
\end{example}

\section{Evolution of differential elements}
Now return to the expression 
\begin{align}
    \frac{d}{dt}\frac{\partial\bm{\phi}_t(\bm{x})}{\partial u_\alpha}=\frac{\partial\bm{\phi}_t(\bm{x})}{\partial u_\alpha}\cdot\nabla\bm{v}\Big|_{(\bm{\phi}_t(\bm{x}),t)},
\end{align}
describing the time rate of change of the vector generating the differential elements along the fluid flow. Using this result and regular calculus rules, the rate of change of the differential line, surface, and volume elements at point $\bm{\phi}_t(\bm{x})$ at instance $t$ along the flow can now be calculated. For the differential line element, one finds by direct substitution
\begin{align}\label{eq:dl_dt}
    \frac{d}{dt}d\bm{\ell}_{\bm{\phi}_t(\bm{x}),\alpha}=\frac{\partial\bm{\phi}_t(\bm{x})}{\partial u_\alpha}\cdot\nabla\bm{v}\Big|_{(\bm{\phi}_t(\bm{x}),t)}du_\alpha=(\nabla \bm{v})^T\Big|_{(\bm{\phi}_t(\bm{x}),t)}\cdot d\bm{\ell}_{\bm{\phi}_t(\bm{x}),\alpha}
\end{align}
For the differential surface element one uses the product rule of differentiation and calculates
\begin{align}\label{eq:dA_dt}
    \frac{d}{dt}d\bm{A}_{\bm{\phi}_t(\bm{x}),\gamma}&=\frac{d}{dt}\left(d\bm{\ell}_{\bm{\phi}_t(\bm{x}),\alpha}\times d\bm{\ell}_{\bm{\phi}_t(\bm{x}),\beta}\right)\nonumber\\
    &=\frac{d}{dt}d\bm{\ell}_{\bm{\phi}_t(\bm{x}),\alpha}\times d\bm{\ell}_{\bm{\phi}_t(\bm{x}),\beta}+d\bm{\ell}_{\bm{\phi}_t(\bm{x}),\alpha}\times \frac{d}{dt}d\bm{\ell}_{\bm{\phi}_t(\bm{x}),\beta}\nonumber\\
    &=\left((\nabla \bm{v})^T\Big|_{(\bm{\phi}_t(\bm{x}),t)}\cdot d\bm{\ell}_{\bm{\phi}_t(\bm{x}),\alpha}\right)\times d\bm{\ell}_{\bm{\phi}_t(\bm{x}),\beta}+d\bm{\ell}_{\bm{\phi}_t(\bm{x}),\alpha}\times\left((\nabla \bm{v})^T\Big|_{(\bm{\phi}_t(\bm{x}),t)}\cdot d\bm{\ell}_{\bm{\phi}_t(\bm{x}),\beta}\right)\nonumber\\
    &=\left((\nabla\cdot\bm{v})\bm{I}-\nabla\bm{v}\right)\Big|_{(\bm{\phi}_t(\bm{x}),t)}\cdot d\bm{A}_{\bm{\phi}_t(\bm{x}),\gamma}
\end{align}
For the differential volume element, one similarly uses the product rule and calculates
\begin{align}\label{eq:dV_dt}
    \frac{d}{dt}dV_{\bm{\phi}_t(\bm{x})}&=\frac{d}{dt}\left(d\bm{A}_{\bm{\phi}_t(\bm{x}),\gamma}\cdot d\bm{\ell}_{\bm{\phi}_t(\bm{x}),\gamma}\right)\nonumber\\
    &=\frac{d}{dt}d\bm{A}_{\bm{\phi}_t(\bm{x}),\gamma}\cdot d\bm{\ell}_{\bm{\phi}_t(\bm{x}),\gamma}+d\bm{A}_{\bm{\phi}_t(\bm{x}),\gamma}\cdot\frac{d}{dt}d\bm{\ell}_{\bm{\phi}_t(\bm{x}),\gamma}\nonumber\\
    &=d\bm{\ell}_{\bm{\phi}_t(\bm{x}),\gamma}\cdot \left((\nabla\cdot\bm{v})\bm{I}-\nabla\bm{v}\right)\Big|_{(\bm{\phi}_t(\bm{x}),t)}\cdot d\bm{A}_{\bm{\phi}_t(\bm{x}),\gamma}+d\bm{A}_{\bm{\phi}_t(\bm{x}),\gamma}\cdot(\nabla\bm{v})^T\Big|_{(\bm{\phi}_t(\bm{x}),t)}\cdot d\bm{\ell}_{\bm{\phi}_t(\bm{x}),\gamma}\nonumber\\
    &=(\nabla\cdot\bm{v})\Big|_{(\bm{\phi}_t(\bm{x}),t)}dV_{\bm{\phi}_t(\bm{x})}
\end{align}
From the rate of change of a differential volume, one identifies that a negative velocity divergence corresponds locally to compression of the material volume, while a positive divergence corresponds to volumetric expansion.

Since $\nabla\cdot\bm{v}=\nabla\bm{v}:\bm{I}$, the time evolution of the differential elements along the flow of the vector field $\bm{v}$ is seen to be determined by the velocity gradient tensor $\nabla\bm{v}$ that encodes the local details of the velocity field. The tensor can be decomposed into three parts
\begin{align}
    \nabla\bm{v}=\bm{D}+\frac{1}{3}(\nabla\cdot\bm{v})\bm{I}+\bm{\Omega},
\end{align}
where the symmetric and traceless tensor $\bm{D}$ and the antisymmetric tensor $\bm{\Omega}$ are 
\begin{align}
    \bm{D}&=\frac{1}{2}\left(\nabla\bm{v}+(\nabla\bm{v})^T\right)-\frac{1}{3}(\nabla\cdot\bm{v})\bm{I},\\
    \bm{\Omega}&=\frac{1}{2}\left(\nabla\bm{v}-(\nabla\bm{v})^T\right).
\end{align}
In the velocity gradient tensor, $\tfrac{1}{3}(\nabla\cdot\bm{v})\bm{I}$ describes isotropic deformation, $\bm{D}$ describes deformation that is non-isotropic but volume-preserving, and $\bm{\Omega}$ describes rotation without deformation. Each part therefore categorizes fluid motion into physically distinct mechanisms that can be visualised by tracking movement of points $\bm{x}(t)$ in the flow of $\bm{v}(\bm{x},t)$.
\begin{example}
    Consider three different velocity fields
    \begin{align*}
        \bm{v}_D&=\frac{v_0}{L_0}(x\bm{\hat{x}}-y\bm{\hat{y}}),
        &\bm{v}_\Omega&=\frac{v_0}{L_0}(-y\bm{\hat{x}}+x\bm{\hat{y}}),
        &\bm{v}_\text{div}&=\frac{v_0}{L_0}(x\bm{\hat{x}}+y\bm{\hat{y}}),
    \end{align*}
    For each field, only one contribution to the velocity gradient tensor is non-zero
    \begin{align*}
        \bm{D}(\bm{v}_D)&=\frac{v_0}{L_0}(\bm{\hat{x}}\bm{\hat{x}}-\bm{\hat{y}}\bm{\hat{y}}), & \bm{\Omega}(\bm{v}_D)&=\bm{0}, & \nabla\cdot\bm{v}_D&=0,\\
        \bm{D}(\bm{v}_\Omega)&=\bm{0}, & \bm{\Omega}(\bm{v}_\Omega)&=\frac{v_0}{L_0}(\bm{\hat{x}}\bm{\hat{y}}-\bm{\hat{y}}\bm{\hat{x}}), & \nabla\cdot\bm{v}_\Omega&=0, \\
        \bm{D}(\bm{v}_\text{div})&=\bm{0}, & \bm{\Omega}(\bm{v}_\text{div})&=\bm{0}, & \nabla\cdot\bm{v}_\text{div}&=\frac{2v_0}{L_0}.
    \end{align*}
    Integrate then points in space according to the differential equation
    \begin{align*}
        \frac{d\bm{x}(t)}{dt}=\bm{v}(\bm{x}(t)),
    \end{align*}
    until a time of $t=v_0/(2L_0)$. A square of points, initially standing on it's corner, changes in the following manner under different velocity fields
    \begin{center}
        \includegraphics[width=0.99\linewidth]{\rootfolder/figs/example-D-Omega-div-individually.pdf}
    \end{center}
    The field $\bm{v}_D$ has stretched the square in the $x$-direction and squeezed it in the $y$-direction, corresponding to non-isotropic volume-preserving deformation. The field $\bm{v}_\Omega$ has rotated the square without deforming it. The field $\bm{v}_\text{div}$ has isotropically scaled the square volume.
    
    More complicated combined effects are observed for velocity fields that have more complicated spatial dependencies. For example, the velocity field
    \begin{align*}
        \bm{v}=v_0\frac{-y\bm{\hat{x}}+x\bm{\hat{y}}}{\sqrt{x^2+y^2}}+\frac{v_0}{L_0^2}(x^2\bm{\hat{x}}+y^2\bm{\hat{y}}),
    \end{align*}
    results in non-uniform rotation and deformation.
    \begin{center}
        \includegraphics[width=0.37\linewidth]{\rootfolder/figs/example-complicated-flow.pdf}
    \end{center}
\end{example}

\section{Leibniz rule for differentiating integrals}
Deriving equations for the evolution of fluid quantities requires differentiating volume integrals over domains that flow with the fluid velocity $\bm{v}(\bm{x},t)$. In other situations, such as when analyzing flow through mechanical devices with moving boundaries, one instead consider domains that move with an arbitrary velocity field $\bm{u}(\bm{x},t)$. In both cases, the differentiation of time-dependent volume integrals is carried out using the Leibniz integral differentiation rule. Its derivation relies on a change of coordinates with the help of the flow map associated with the flow, and on the time derivative of a volume element under the flow that waas derived in the previous section.

Consider the integral of a scalar-, vector-, or tensor-valued field $\bm{f}(\bm{x},t)$ over time-dependent volume $V(t)$ evolving according to a vector field $\bm{u}(\bm{x},t)$ and having closed surface $\partial V(t)$
\begin{align}
    \int_{V(t)}\bm{f}(\bm{x},t)dV_{\bm{x}}.
\end{align}
Assume the flow map associated with the velocity field $\bm{u}(\bm{x},t)$ from some fixed moment $t_0$ onwards to be $\bm{\psi}_t$. The volume $V(t)$ can now be expressed as an image set 
\begin{align}
    V(t)=\text{Im}_{\bm{\psi}_t}(V(t_0))=\{\bm{\psi}_t(\bm{x}_0)\,|\,\bm{x}_0\in V(t_0)\},
\end{align}
of a fixed volume $V(t_0)$ and to be generated by the flow map $\psi_t$ moving initial points $\bm{x}_0$ in to the new points $\bm{x}_0$. Introduce a change of coordinates on the integration labels according to $\bm{x}=\bm{\psi}_t(\bm{x}_0)$. Then
\begin{align}
    \int_{V(t)}\bm{f}(\bm{x},t)dV_{\bm{x}}=\int_{\text{Im}_{\bm{\psi}_t}(V(t_0))}\bm{f}(\bm{x},t)dV_{\bm{x}}=\int_{V(t_0)}\bm{f}(\bm{\psi}_t(\bm{x}_0),t)dV_{\bm{\psi}_t(\bm{x}_0)},
\end{align}
where the integration on right is now to be performed over the labelled points $\bm{x}_0\in V(t_0)$. As the domain $V(t_0)$ of the label points $\bm{x}_0$ is independent of the time $t$, the time derivative can be taken inside the integral and 
\begin{align}
    \frac{d}{dt}\int_{V(t)}\bm{f}(\bm{x},t)dV_{\bm{x}}=\int_{V(t_0)}\frac{d}{dt}\left(\bm{f}(\bm{\psi}_t(\bm{x}_0),t)dV_{\bm{\psi}_t(\bm{x}_0)}\right).
\end{align}
The product and composition rule of differentiation, and the expression for the time derivative of a differential volume element under the flow, then provide
\begin{align}
    &\frac{d}{dt}\left(\bm{f}(\bm{\psi}_t(\bm{x}_0),t)dV_{\bm{\psi}_t(\bm{x}_0)}\right)\nonumber\\
    &=\frac{d}{dt}\bm{f}(\bm{\psi}_t(\bm{x}_0),t) dV_{\bm{\psi}_t(\bm{x}_0)}+\bm{f}(\bm{\psi}_t(\bm{x}_0),t)\frac{d}{dt}dV_{\bm{\psi}_t(\bm{x}_0)}\nonumber\\
    &=(\partial_t\bm{f}+\bm{u}\cdot\nabla\bm{f})\Big|_{(\bm{\psi}_t(\bm{x}_0),t)}dV_{\bm{\psi}_t(\bm{x}_0)}+(\bm{f}\nabla\cdot\bm{u})\Big|_{(\bm{\psi}_t(\bm{x}_0),t)}dV_{\bm{\psi}_t(\bm{x}_0)},
\end{align}
and the time derivative of the integral can be written
\begin{align}
    \frac{d}{dt}\int_{V(t)}\bm{f}(\bm{x},t)dV_{\bm{x}}=\int_{V(t_0)}\left(\partial_t\bm{f}+\bm{u}\cdot\nabla\bm{f}+\bm{f}\nabla\cdot\bm{u}\right)\Big|_{(\bm{\psi}_t(\bm{x}_0),t)}dV_{\bm{\psi}_t(\bm{x}_0)}.
\end{align}
Performing then a change of coordinates back to the original ones with $\bm{x}=\bm{\psi}_t(\bm{x}_0)$, one is able to return the integral to be evaluated over the original volume $V(t)$ at time $t$
\begin{align*}
    \int_{V(t_0)}\left(\partial_t\bm{f}+\bm{u}\cdot\nabla\bm{f}+\bm{f}\nabla\cdot\bm{u}\right)\Big|_{(\bm{\psi}_t(\bm{x}_0),t)}dV_{\bm{\psi}_t(\bm{x}_0)}=\int_{V(t)}\left(\partial_t\bm{f}+\bm{u}\cdot\nabla\bm{f}+\bm{f}\nabla\cdot\bm{u}\right)\Big|_{(\bm{x},t)}dV_{\bm{x}}.
\end{align*}
In applications, the explicit labeling of the integration domain with points $\bm{x}$ is usually dropped, resulting in the compact formula for the Leibniz integral differentiation rule
\begin{align}\label{eq:leibniz-rule}
    \frac{d}{dt}\int_{V(t)}\bm{f}dV&=\int_{V(t)}\left(\partial_t\bm{f}+\bm{u}\cdot\nabla\bm{f}+\bm{f}\nabla\cdot\bm{u}\right)dV\nonumber\\
    &=\int_{V(t)}\left(\partial_t\bm{f}+\nabla\cdot(\bm{u}\otimes\bm{f})\right)dV,
\end{align}
The divergence theorem can then be used to provide a meaningful physical interpretation for the Leibniz integral differentiation rule
\begin{align}
    \underbrace{\frac{d}{dt}\int_{V(t)}\bm{f}dV}_{\substack{\text{Rate of change}\\\text{of a volume integral}\\\text{over a moving volume}}}&=\underbrace{\int_{V(t)}\partial_t \bm{f}dV}_{\substack{\text{Volume integral}\\\text{over the}\\\text{local change}}} +\underbrace{\int_{\partial V(t)}\bm{f}\underbrace{\bm{u}\cdot\bm{\hat{n}} dA}_{\substack{\text{Local rate of}\\\text{change of}\\\text{the volume}}}}_{\substack{\text{Net flow of $\bm{f}$ into $V$}\\\text{due to movement of $\partial V$}}}.
\end{align}
\begin{example}
    The three-dimensional derivation above combines several unfamiliar mathematical steps. The same construction is illustrated below in one dimension, where the individual steps can be followed using the ordinary chain and product rules. The calculation also recovers the familiar one-dimensional Leibniz rule.

    To differentiate the time-dependent integral
    \begin{align*}
        \int_{a(t)}^{b(t)} f(x,t) dx,
    \end{align*}
    with $a(t)$ and $b(t)$ some time-dependent, well-behaving functions describing the positions of the integral boundaries in time $t$, assume any point $x\in[a(t),b(t)]$ to move with the speed $u(x,t)$. Let $\psi(x,t)$ be the flow map of $u(x,t)$ so that a point $x_0\in[a(t_0),b(t_0)]$ becomes a point $\psi(x_0,t)=x$ and $\partial_t\psi(x,t)=u(\psi(x,t),t)$. One may now write
    \begin{align*}
        \int_{a(t)}^{b(t)} f(x,t) dx&=\int_{\psi(a(t_0),t)}^{\psi(b(t_0),t)}f(x,t)dx\\
        &=\int_{a(t_0)}^{b(t_0)}f(\psi(x_0,t),t)d\psi(x_0,t)\\
        &=\int_{a(t_0)}^{b(t_0)}f(\psi(x_0,t),t)\frac{\partial\psi}{\partial x}\Big|_{(x_0,t)}dx_0
    \end{align*}
    To proceed, apply the chainrule to the derivative of the flow map
    \begin{align*}
        \frac{d}{dt}\frac{\partial\psi}{\partial x}\Big|_{(x_0,t)}=\frac{\partial^2\psi(x,t)}{\partial x\partial t}\Big|_{(x_0,t)}=\frac{\partial u(\psi(x,t),t)}{\partial x}\Big|_{(x_0,t)}=\frac{\partial u}{\partial x}\Big|_{(\psi(x_0,t),t)}\frac{\partial\psi}{\partial x}\Big|_{(x_0,t)}.
    \end{align*}
    One may now identify that 
    \begin{align*}
        \frac{d}{dt}d\psi(x_0,t)=\frac{d}{dt}\frac{\partial\psi}{\partial x}\Big|_{(x_0,t)}dx_0=\frac{\partial u}{\partial x}\Big|_{(\psi(x_0,t),t)}\frac{\partial\psi}{\partial x}\Big|_{(x_0,t)}dx_0=\frac{\partial u}{\partial x}\Big|_{(\psi(x_0,t),t)}d\psi(x_0,t),
    \end{align*}
    is analogous to the time derivative of a differential volume element with $\partial u/\partial x$ serving the role of velocity divergence. Considering that the time derivative of the composed function $f(\psi(x_0,t),t)$ can be written 
    \begin{align*}
        \frac{d}{dt}f(\psi(x_0,t),t)=\frac{\partial f}{\partial t}\Big|_{(\psi(x_0,t),t)}+\frac{\partial f}{\partial x}\Big|_{(\psi(x_0,t),t)}\frac{\partial \psi}{\partial  t}\Big|_{(x_0,t)}=\left(\frac{\partial f}{\partial t}+u\frac{\partial f}{\partial x}\right)\Big|_{(\psi(x_0,t),t)},
    \end{align*}
    the time derivative of the integral becomes
    \begin{align*}
        \frac{d}{dt}\int_{a(t)}^{b(t)} f(x,t) dx&=\int_{a(t_0)}^{b(t_0)}\frac{d}{dt}\left(f(\psi(x_0,t),t)\frac{\partial\psi}{\partial x}\Big|_{(x_0,t)}\right)dx_0\\
        &=\int_{a(t_0)}^{b(t_0)}\left(\frac{\partial f}{\partial t}+u\frac{\partial f}{\partial x}+f\frac{\partial u}{\partial x}\right)\Big|_{(\psi(x_0,t),t)}\frac{\partial\psi}{\partial x}\Big|_{(x_0,t)}dx_0\\
        &=\int_{a(t)}^{b(t)}\left(\frac{\partial f}{\partial t}+\frac{\partial(uf)}{\partial x}\right)\Big|_{(x,t)}dx\\
        &=\int_{a(t)}^{b(t)}\frac{\partial f(x,t)}{\partial t}dx+\Big[u(x,t)f(x,t)\Big]_{x=a(t)}^{x=b(t)}
    \end{align*}
    The result is the one-dimensional version of the Leibniz rule. 
    Realizing that the speed function $u(x,t)$ at the boundary points satisfies $u(a(t),t)=da/dt$ and $u(b(t),t)=db/dt$ enables writing
    \begin{align*}
        \frac{d}{dt}\int_{a(t)}^{b(t)}f(x,t)dx=\int_{a(t)}^{b(t)}\frac{\partial f(x,t)}{\partial t}dx+f(b(t),t)\frac{db(t)}{dt}-f(a(t),t)\frac{da(t)}{dt},
    \end{align*}
    which recovers the familiar and well-known rule for differentiating one-dimensional integrals with changing boundaries.
\end{example}


\section{Conservation of mass}
Consider a volume $V(t)$ of fluid moving along the fluid's velocity field $\bm{v}(\bm{x},t)$ as in Fig.~\ref{fig:mass-volume}. The mass inside the volume $V(t)$ is calculated by integrating the density of the fluid over the volume. The time rate of change of the mass contained in the volume is then obtained with the Leibniz rule, providing
\begin{align}
    \underbrace{\frac{d}{dt}\int_{V(t)}\rho dV}_{\substack{\text{Rate of change of}\\\text{fluid mass in $V(t)$}}}=\underbrace{\int_{V(t)}\partial_t\rho \,dV}_{\substack{\text{Volume integral}\\\text{of the local rate}\\\text{of change of density}}}+\underbrace{\int_{\partial V(t)}\rho\underbrace{\bm{v}\cdot\bm{\hat{n}}dA}_{\substack{\text{Rate of}\\\text{change of}\\\text{the volume}}}}_{\substack{\text{Net flow of mass into $V$}\\\text{due to movement of $\partial V$}}}=\int_{V(t)}\left(\partial_t\rho+\nabla\cdot(\bm{v}\rho)\right)dV.
\end{align}
The volume is tied to the flow, and therefore the mass of fluid inside the volume $V(t)$ does not change in time. One therefore has
\begin{align}
    \int_{V(t)}\left(\partial_t\rho+\nabla\cdot(\bm{v}\rho)\right)dV=0.
\end{align}
The condition has to hold regardless of the initial choice for $V(t)$. Therefore the fluid density has to satisfy the following partial differential equation, often referred to as the fluid's continuity or density equation
\begin{align}\label{eq:density}
    \partial_t\rho+\nabla\cdot(\bm{v}\rho)=0.
\end{align}
\begin{figure}[!htb]
    \centering
    \input{\rootfolder/code/fig-mass-volume}
    \caption{Evolution of the volume $V(t)$ tied to the flow of the fluid velocity field $\bm{v}(\bm{x},t)$.}
    \label{fig:mass-volume}
\end{figure}
\begin{example}
    Consider a velocity field of the form
    \begin{align*}
        \bm{v}(\bm{x},t)=v(z)\bm{\hat{z}}
    \end{align*}
    describing one-dimensional time-independent flow into $\bm{\hat{z}}$-direction with speed $v(z)$ that may vary with respect to the coordinate $z$. Assuming that also the density is steady in time and changes only with respect to the coordinate $z$, the fluid's density equation reduces to
    \begin{align*}
        \partial_t\rho+\nabla\cdot(\bm{v}\rho) \quad \implies \quad \frac{d}{dz}(v(z)\rho(z))=0
    \end{align*}
    This means the scalar velocity $v(z)$ times the density $\rho(z)$ has to be a constant, and that density is given by
    \begin{align*}
        \rho(z)=\frac{\rho_0v_0}{v(z)}
    \end{align*}
    with $\rho_0$ and $v_0$ some constants. If the fluid at some location $z$ flows with slower speed, the density at that location must be higher to satisfy the conservation of mass. Conversely, if the fluid locally has a higher speed, the density must be lower.
\end{example}

The physical law of mass conservation has a convenient mathematical consequence. Over a volume $V(t)$ moving with the fluid's velocity field $\bm{v}(\bm{x},t)$, the density-weighted rate of change of any quantity $\bm{f}(\bm{x},t)$, be it scalar, vector, or tensor, becomes expressible
\begin{align}\label{eq:rho-identity}
    \frac{d}{dt}\int_{V(t)}\rho \bm{f} dV
    %&=\int_{V(t)}\left(\partial_t(\rho\bm{f})+\nabla\cdot(\bm{v}\rho\bm{f})\right)dV\nonumber\\
    &=\int_{V(t)}\rho\left(\partial_t\bm{f}+\bm{v}\cdot\nabla \bm{f}\right)dV
    %\nonumber\\&
    =\int_{V(t)}\rho\frac{D\bm{f}}{Dt}dV,
\end{align}
On the right, the scalar-valued differential operator 
\begin{align}\label{eq:material-derivative}
    \frac{D}{Dt}=\frac{\partial}{\partial t}+\bm{v}\cdot\nabla,
\end{align}
is typically referred to as the material derivative. The term and the operator are used extensively in fluid and solid mechanics texts. In terms of it, the equation~\eqref{eq:density} for density can be written also as
\begin{align}
    \frac{1}{\rho}\frac{D\rho}{Dt}=-\nabla\cdot\bm{v}.
\end{align}
Further, the scalar differential operator $D/Dt$ illustrates the connection between the so-called Lagrangian and Eulerian views of motion. 
\begin{example}
    The acceleration of a point $\bm{x}(t)$ following the velocity field $\bm{v}(\bm{x},t)$ equals the material derivative of the velocity field evaluated along the trajectory at the point $\bm{x}(t)$
    \begin{align*}
        \frac{d^2\bm{x}(t)}{dt^2}=\frac{d\bm{v}(\bm{x}(t),t)}{dt}=(\partial_t\bm{v}+\bm{v}\cdot\nabla\bm{v})\Big|_{(\bm{x}(t),t)}=\frac{D\bm{v}}{Dt}\Big|_{(\bm{x}(t),t)}.
    \end{align*}
\end{example}


\section{Conservation of momentum}
According to Newton's second law, the rate of change of linear momentum of a body equals the net forces acting on it. The forces acting on a volume of fluid include gravity that is distributed throughout the fluid volume, and mechanical contact force that is distributed on the surface of the body. Since the total force must be finite, yet distributed over the body volume and surface, both the body and surface force need to be treated differentially while the total force is an integral of the differential force. 

Denote with $\delta \bm{F}_b$ the differential body force from gravity acting on a differential volume $dV$ of the body. The differential mass the volume $dV$ carries equals $\rho dV$, and the differential body force from gravity acting on this volume therefore is given by
\begin{align}
    \delta\bm{F}_b=\rho\bm{g}dV,
\end{align}
where $\bm{g}$ is the gravitational acceleration vector. Denote then with $\delta \mathbf{F}_s$ the differential surface force from mechanical contact acting on a differential surface area $dA$ oriented with respect to a surface unit normal vector $\mathbf{\hat{n}}$. Let the surface force density per unit area be a traction vector $\mathbf{t}(\mathbf{\hat{n}})$ that is dependent on the unit normal of the surface inspected. While the underlying mechanical state at a point within the material is independent of how we choose to observe it, the traction vector changes with the orientation of the imaginary surface on which it is evaluated. This reflects the nature of a contact force and Newton's third law. For every action, there is a counter action. In case of a contact force, the counteraction happens at the same point and must depend on the orientation chosen for inspection. The differential force acting on a surface is therefore expressible as
\begin{align}
    \delta\bm{F}_s=\bm{t}(\bm{\hat{n}})dA.
\end{align}
Together, the differential forces and the related volume and surface elements are illustrated in Fig.~\ref{fig:forces-velocity}. The velocity vector $\bm{v}$ in the Figure denotes the potential movement direction of the body in a specific point on its surface.
\begin{figure}[!htb]
    \centering
    \input{\rootfolder/code/fig-forces-velocity}
    \caption{The differential surface and body forces $\delta\bm{F}_s$ and $\delta\bm{F}_b$ acting on a fluid volume.}
    \label{fig:forces-velocity}
\end{figure}

Assuming the fluid moves with velocity $\bm{v}$, the differential amount of momentum the fluid carries in a differential volume $dV$ is $\rho\bm{v}dV$. For a macroscopic volume $V(t)$ moving along the fluid's velocity field $\bm{v}$, Newton's law then reads
\begin{align}\label{eq:newton}
    \underbrace{\frac{d}{dt}\int_{V(t)}\rho\bm{v}dV}_{\substack{\text{Rate of change of}\\\text{momentum within}\\\text{volume $V(t)$}}} = \underbrace{\int_{V(t)}\delta\bm{F}_b}_{\substack{\text{Force acting}\\\text{on fluid inside}\\\text{the volume}}} + \underbrace{\int_{\partial V(t)}\delta\bm{F}_s}_{\substack{\text{Force from}\\\text{mechanical}\\\text{surface contact}}}.
\end{align}
On the left of Eq.~\eqref{eq:newton}, the material derivative identity~\eqref{eq:rho-identity} allows one to move the time derivative inside the integral
\begin{align}
    \frac{d}{dt}\int_{V(t)}\rho\bm{v}dV = \int_{V(t)}\rho\frac{D\bm{v}}{Dt}dV.
\end{align}
On the right of Eq.~\eqref{eq:newton}, the total force involves both a volume and a surface integral
\begin{align}
    \int_{V(t)}\delta \bm{F}_b+\int_{\partial V(t)}\delta\bm{F}_s=\int_{V(t)}\rho\bm{g}dV+\int_{\partial V(t)}\bm{t}(\bm{\hat{n}})dA,
\end{align}
which, for now, prevents finding a local differential equation to describe the time derivative of the fluid velocity. The surface integral on the right must first be expressed as a volume integral. This is achieved by requiring Newton's second law to hold over arbitrary volumes and investigating two specific volumes to derive a linear relation for the traction with respect to the surface unit normal vector
\begin{align}
    \bm{t}(\bm{\hat{n}})=\bm{\sigma}\cdot\bm{\hat{n}},
\end{align} 
where the rank-2 tensor $\bm{\sigma}$ is referred to as stress. The process is detailed below.
\begin{example}\label{ex:traction}
    For a continuous vector valued function $\bm{f}:\Omega\mapsto \mathbb{R}^n$ defined over a connected set $\Omega\subset\mathbb{R}^m$, the Mean Value Theorem tells that the average value of the integral of $\bm{f}$ over the set $\Omega$ lies in the convex hull of the image of $\Omega$ under $\bm{f}$, namely 
    \begin{align*}
        \frac{1}{vol(\Omega)}\int_\Omega\bm{f}d^m\bm{x} \in \text{conv}(\{\bm{f}(\bm{x})|\bm{x}\in\Omega\}).
    \end{align*}
    Using this result, the traction vector can be shown to satisfy the identity
    \begin{align*}
        \bm{t}(\bm{x},\bm{\hat{n}}(\bm{x}))=\bm{\sigma}(\bm{x})\cdot\bm{\hat{n}}(\bm{x}).
    \end{align*}
    
    Choose an $\epsilon$-parametrized volume $V(t,\epsilon)$ moving with the flow, and which at the instant $t$ is a cylinder centered at the point $\bm{x}$ with the symmetry axis aligned with the unit normal vector $\bm{\hat{n}}(\bm{x})$ of the surface on which one inspects the traction vector. See the figure for clarification.
    \begin{center}
        \input{\rootfolder/code/fig-cylinder} 
    \end{center}
    Denote the top, bottom, and shell surfaces of the cylinder as the sets of points $S_n(t,\epsilon)$, $S_{-n}(t,\epsilon)$ and $S_r(t,\epsilon)$ with corresponding areas $A_n(\epsilon)=\pi R(\epsilon)^2$ for the top and bottom, and $A_r(\epsilon)=2\pi R(\epsilon)H(\epsilon)$ for the shell. The volume of the cylinder is $V(t,\epsilon)=A_n(\epsilon)H(\epsilon)$. According to the mean value theorem, one now has
    \begin{align*}
        \frac{1}{V(t,\epsilon)}\int_{V(t,\epsilon)}\left(\rho\frac{D\bm{v}}{Dt}-\rho\bm{g}\right)dV&=\bm{F}(\epsilon)\in \text{conv}\left(\left\{\left(\rho\frac{D\bm{v}}{Dt}-\rho\bm{g}\right)_{\bm{y}}|\bm{y}\in V(t,\epsilon)\right\}\right),
        \\
        \frac{1}{A_n(t,\epsilon)}\int_{S_n(t,\epsilon)}\bm{t}(\bm{\hat{n}})dA&= \bm{t}_n(\epsilon)\in\text{conv}(\{\bm{t}(\bm{y},\bm{\hat{n}}(\bm{x}))|\bm{y}\in S_n(t,\epsilon)\}),
        \\
        \frac{1}{A_n(t,\epsilon)}\int_{S_{-n}(t,\epsilon)}\bm{t}(\bm{\hat{n}})dA&=\bm{t}_{-n}(\epsilon)\in \text{conv}(\{\bm{t}(\bm{y},-\bm{\hat{n}}(\bm{x}))|\bm{y}\in S_{-n}(t,\epsilon)\}),\\
        \frac{1}{A_r(t,\epsilon)}\int_{S_r(t,\epsilon)}\bm{t}(\bm{\hat{n}})dA&=\bm{t}_r(\epsilon)\in \text{conv}(\{\bm{t}(\bm{y},\bm{\hat{r}}(\bm{y}))|\bm{y}\in S_r(t,\epsilon)\}).
    \end{align*}
    At instance $t$, Newton's law~\eqref{eq:newton} over the volume $V(t,\epsilon)$ then becomes
    \begin{align*}
        \bm{F}(\epsilon) V(t,\epsilon)&=(\bm{t}_n(\epsilon)+\bm{t}_{-n}(\epsilon))A_n(t,\epsilon)+\bm{t}_r(\epsilon)A_r(t,\epsilon).
    \end{align*} 
    Divide both sides with the area $A_n(t,\epsilon)$ and choose the parametrization $R(\epsilon)=\epsilon R_0$, $H(\epsilon)=\epsilon^2H_0$. At the limit $\epsilon\to 0$, the ratios $V(t,\epsilon)/A_n(t,\epsilon)\to 0$ and $A_r(t,\epsilon)/A_n(t,\epsilon)\to 0$ vanish, the sets $S_n(t,\epsilon)\to \bm{x}$ and $S_{-n}(t,\epsilon)\to\bm{x}$ collapse to the point $\bm{x}$, while for the vectors one has $\bm{t}_n(\epsilon)\to\bm{t}(\bm{x},\bm{\hat{n}}(\bm{x}))$ and $\bm{t}_{-n}(\epsilon)\to\bm{t}(\bm{x},-\bm{\hat{n}}(\bm{x}))$. At the limit, Newton's law for the chosen cylindrical volume therefore provides a relation
    \begin{align*}
        \bm{t}(\bm{x},\bm{\hat{n}}(\bm{x}))=-\bm{t}(\bm{x},-\bm{\hat{n}}(\bm{x})).
    \end{align*}
    
    Mathematically the result implies that the traction vector is odd with respect to the dependency on the unit normal of the surface it acts on. Physically the result reflects Newton's third law. A surface has outward side ($\bm{\hat{n}}$) and an inward side ($-\bm{\hat{n}}$). On any action on the outward side, there is always a corresponding reaction on the inward side, and of exactly the same magnitude but into opposite direction, as Newton's third law says. This effect is not seen locally for the body forces because the reaction to gravity does not happen inside the fluid but at a distant point outside the fluid where Earth, the cause for gravity, resides.

    Next consider a different $\epsilon$-parametrized volume $V(t,\epsilon)$, which at instance $t$ is a trirectangular tetrahedron attached to the point $\bm{x}$ from one corner. Set the direction of the face with the largest surface area to align with the normal $\bm{\hat{n}}(\bm{x})$ as in the figure below.
    \begin{center}
        \input{\rootfolder/code/fig-tetrahedron}
    \end{center}
    Define the sets $S_i(t,\epsilon)$ to correspond to the faces of the tetrahedron aligned with the local orthogonal unit normal vectors $\bm{\hat{e}}_i(\bm{x})$ and the set $S_n(t,\epsilon)$ to correspond to the face aligned with $\bm{\hat{n}}(\bm{x})$. The corresponding areas are $A_i(t,\epsilon)$ and $A_n(t,\epsilon)$ and satisfy
    \begin{align*}
        A_i(t,\epsilon)=\bm{\hat{e}}_i(\bm{x})\cdot\bm{\hat{n}}(\bm{x})A_n(t,\epsilon).
    \end{align*}
    Now the mean value theorem provides
     \begin{align*}
        \frac{1}{V(t,\epsilon)}\int_{V(t,\epsilon)}\left(\rho\frac{D\bm{v}}{Dt}-\rho\bm{g}\right)dV&=\bm{F}(\epsilon)\in \text{conv}\left(\left\{\left(\rho\frac{D\bm{v}}{Dt}-\rho\bm{g}\right)_{\bm{y}}|\bm{y}\in V(t,\epsilon)\right\}\right),
        \\
        \frac{1}{A_n(t,\epsilon)}\int_{S_n(t,\epsilon)}\bm{t}(\bm{\hat{n}})dA&= \bm{t}_n(\epsilon)\in \text{conv}(\{\bm{t}(\bm{y},\bm{\hat{n}}(\bm{x}))|\bm{y}\in S_n(t,\epsilon)\}),
        \\
        \frac{1}{A_i(t,\epsilon)}\int_{S_i(t,\epsilon)}\bm{t}(\bm{\hat{n}})dA&=\bm{t}_i(\epsilon)\in \text{conv}(\{\bm{t}(\bm{y},-\bm{\hat{e}}_i(\bm{x}))|\bm{y}\in S_i(t,\epsilon)\}).
    \end{align*}
    At instance $t$, Newton's law~\eqref{eq:newton} over the volume $V(t,\epsilon)$ then becomes
    \begin{align*}
        \bm{F}(\epsilon) V(t,\epsilon)&=\bm{t}_n(\epsilon)A_n(t,\epsilon)+\sum_{i=1}^3\bm{t}_i(\epsilon)A_i(t,\epsilon).
    \end{align*} 
    Divide both sides with $A_n(t,\epsilon)$ and choose the parametrization $L_i=\epsilon L_{i,0}$. At the limit $\epsilon\to 0$, the ratio $V(t,\epsilon)/A_n(t,\epsilon)\to 0$ vanishes, the sets $S_i(t,\epsilon)\to\bm{x}$ and $S_n(t,\epsilon)\to\bm{x}$ collapse to point $\bm{x}$, while for the vectors one has $\bm{t}_i(\epsilon)\to \bm{t}(\bm{x},-\bm{\hat{e}}_i(\bm{x}))$ and $\bm{t}_n(\epsilon)\to\bm{t}(\bm{x},\bm{\hat{n}}(\bm{x}))$. At the limit, Newton's law for the chosen tetrahedron volume therefore provides the relation
    \begin{align*}
        \bm{t}(\bm{x},\bm{\hat{n}}(\bm{x}))=-\sum_{i=1}^3\bm{\hat{e}}_i(\bm{x})\cdot\bm{\hat{n}}(\bm{x}) \,\bm{t}(\bm{x},-\bm{\hat{e}}_i(\bm{x})).
    \end{align*}
    
    The result established with the cylindrical volume can be applied with respect to any surface. In particular, one could choose to study traction vectors on surfaces aligned with respect to the $\bm{\hat{e}}_i(\bm{x})$ to obtain $\bm{t}(\bm{x},\bm{\hat{e}}_i(\bm{x}))=-\bm{t}(\bm{x},-\bm{\hat{e}}_i(\bm{x}))$. Now one can write
    \begin{align*}
        \bm{t}(\bm{x},\bm{\hat{n}}(\bm{x}))=\sum_{i=1}^3\bm{\hat{e}}_i(\bm{x})\cdot\bm{\hat{n}}(\bm{x}) \,\bm{t}(\bm{x},\bm{\hat{e}}_i(\bm{x})).
    \end{align*}
    With respect to the basis vectors $\bm{\hat{e}}_i(\bm{x})$, the surface normal can be written $\bm{\hat{n}}(\bm{x})=\sum_{i=1}^3n_i(\bm{x})\,\bm{\hat{e}}_i(\bm{x})$ where the coefficients are $n_i(\bm{x})=\bm{\hat{n}}(\bm{x})\cdot\bm{\hat{e}}_i(\bm{x})$. Therefore
    one finds
    \begin{align*}
        \bm{t}\left(\bm{x},\sum_{i=1}^3n_i(\bm{x})\,\bm{\hat{e}}_i(\bm{x})\right)=\sum_{i=1}^3n_i(\bm{x})\,\bm{t}(\bm{x},\bm{\hat{e}}_i(\bm{x}))
    \end{align*}
    This condition on $\bm{t}$ defines it to be a linear relation with respect to the surface normal $\bm{\hat{n}}$. Therefore
    \begin{align*}
        \bm{t}(\bm{x},\bm{\hat{n}}(\bm{x}))=\bm{\sigma}(\bm{x})\cdot\bm{\hat{n}}(\bm{x}),
    \end{align*}
    where $\bm{\sigma}(\bm{x})$ is a tensor, referred to as stress.
\end{example}

The linear dependency of traction on the surface unit normal then enables the use of divergence theorem to convert the surface integral into a volume integral
\begin{align}
    \int_{\partial V(t)}\bm{t}(\bm{\hat{n}})dA=\int_{\partial V(t)}\bm{\sigma}\cdot\bm{\hat{n}}dA=\int_{V(t)}\nabla\cdot\bm{\sigma} dV,
\end{align}
and the Newton's second law over the volume $V(t)$ can be seen to equal the equation
\begin{align}
    \int_{V(t)}\rho\frac{D\bm{v}}{Dt}dV=\int_{V(t)}(\rho \bm{g}+\nabla\cdot \bm{\sigma}) dV.
\end{align}
Since the initial choice of the volume $V(t)$ is arbitrary, the condition must hold true for all volumes within the fluid. Consequently, the fluid velocity has to satisfy the following local partial differential equation
\begin{align}\label{eq:cauchy-momentum}
    \rho\frac{D\bm{v}}{Dt}=\rho\bm{g}+\nabla\cdot\bm{\sigma}.
\end{align}
This equation is commonly referred to as the Cauchy's momentum equation. 
\begin{example}
    Recall that the material derivative of the velocity field at a point $\bm{x}(t)$ moving along the fluid flow represented the local acceleration of the point
    \begin{align*}
        \frac{d^2\bm{x}(t)}{dt^2}=\frac{D\bm{v}}{Dt}\Big|_{(\bm{x}(t),t)}.
    \end{align*}
    Cauchy's momentum equation then tells that the local acceleration of the point $\bm{x}(t)$ is caused by gravity and divergence of the stress tensor
    \begin{align*}
        \frac{d^2\bm{x}(t)}{dt^2}=\left(\bm{g}+\rho^{-1}\nabla\cdot\bm{\sigma}\right)\Big|_{(\bm{x}(t),t)}.
    \end{align*}
\end{example}
\begin{example}
    Consider the fluid is at rest and not moving at all, meaning the velocity field vanishes $\bm{v}(\bm{x},t)=\bm{0}$. The Cauchy's momentum equation still holds true, and reads
    \begin{align*}
        \rho\bm{g}+\nabla\cdot\bm{\sigma}&=\bm{0}, && \text{Fluid at rest}
    \end{align*}
    From experiments one knows that the gravity in column of static water or air is balanced by pressure difference at the column bottom and top. For fluid at rest, the divergence of the stress tensor must therefore reduce to the negative of a the pressure gradient. The only form of stress then is pressure and
    \begin{align*}
        \bm{\sigma}&=-p\bm{I}, && \text{Fluid at rest}
    \end{align*}
\end{example}

\section{Conservation of angular momentum}
For point masses, the definition $d\bm{x}/dt=\bm{v}$ for velocity together with Newton's second law $m\,d\bm{v}/dt=\bm{F}$ imply a vector-valued equation $d\bm{L}/dt=\bm{T}$ between the angular momentum vector $\bm{L}=\bm{x}\times m\bm{v}$ and the torque vector $\bm{T}=\bm{x}\times\bm{F}$. Similarly for continuous matter, the time rate of change of angular momentum in a volume $V(t)$ bound to the matter must equal the total torque on it. This requirement is observed to result in an additional restriction on the stress tensor.

The differential amount of angular momentum a differential volume $dV$ of the fluid carries is $\delta\bm{L}=\bm{x}\times \rho\bm{v}dV$. With $\delta\bm{F}_b=\rho\bm{g}dV$ the differential body force acting on a differential volume $dV$ and $\delta\bm{F}_s=\bm{\sigma}\cdot\bm{\hat{n}}dA$ the differential surface force acting on a differential surface area $dA$, the differential torques resulting from body and surface forces are $\delta\bm{T}_b=\bm{x}\times\delta\bm{F}_b=\bm{x}\times\rho\bm{g}dV$ and $\delta\bm{T}_s=\bm{x}\times\delta \bm{F}_s=\bm{x}\times\bm{\sigma}\cdot\bm{\hat{n}}dA$ respectively. For the time rate of change of angular momentum of the fluid within a macroscopic material volume $V(t)$ to equal the total torque on the fluid volume, the following condition must then hold
\begin{align}
    \underbrace{\frac{d}{dt}\int_{V(t)}\delta\bm{L}}_{\substack{\text{Rate of change of}\\\text{angular momentum}\\\text{within $V(t)$}}}=\underbrace{\int_{V(t)}\delta\bm{T}_b}_{\substack{\text{Torque on}\\\text{volume $V(t)$}\\\text{due to body}\\\text{forces}}}+\underbrace{\int_{\partial V(t)}\delta\bm{T}_s}_{\substack{\text{Torque on}\\\text{volume $V(t)$}\\\text{due to surface}\\\text{forces}}}.
\end{align}

Since $V(t)$ is a material volume and $\delta\bm{L}=\bm{x}\times\rho\bm{v}dV$ is proportional to the density, the expression on left can be simplified in terms of the material derivative using the identity~\eqref{eq:rho-identity}. Recalling also that the material derivative is a scalar differential operator, one may apply the usual product rule for differentiation and obtain
\begin{align}
    \frac{d}{dt}\int_{V(t)}\bm{x}\times\rho\bm{v}dV\underset{\eqref{eq:rho-identity}}{=}\int_{V(t)}\rho\frac{D(\bm{x}\times\bm{v})}{Dt}dV=\int_{V(t)}\left(\frac{D\bm{x}}{Dt}\times\rho\bm{v}+\bm{x}\times\rho\frac{D\bm{v}}{Dt}\right)dV
\end{align}
Because $D\bm{x}/Dt=\bm{v}$, the first cross product on the right vanishes identically. Substituting Cauchy's momentum equation~\eqref{eq:cauchy-momentum} into the second cross product, one finds
\begin{align}
    \frac{d}{dt}\int_{V(t)}\delta\bm{L}=\int_{V(t)}\bm{x}\times\rho\frac{D\bm{v}}{Dt}dV=\int_{V(t)}\left(\bm{x}\times\rho\bm{g}+\bm{x}\times\nabla\cdot\bm{\sigma}\right)dV.
\end{align}

On the right of the angular momentum equation, the total torque is
\begin{align}
    \int_{V(t)}\delta\bm{T}_b+\int_{\partial V(t)}\delta\bm{T}_s&=\int_{V(t)}\bm{x}\times\rho\bm{g}dV+\int_{\partial V(t)}\bm{x}\times\bm{\sigma}\cdot\bm{\hat{n}}dA\nonumber\\
    &=\int_{V(t)}(\bm{x}\times\rho\bm{g}+\nabla\cdot(\bm{x}\times\bm{\sigma}))dV.
\end{align}
Equating this expression with the rate of change of angular momentum obtained above, the body-force contributions cancel and the angular momentum condition reduces to
\begin{align}
    \int_{V(t)}\bm{x}\times\nabla\cdot\bm{\sigma}dV=\int_{V(t)}\nabla\cdot(\bm{x}\times\bm{\sigma})dV.
\end{align}
Since the material volume $V(t)$ is arbitrary, the integrands must be equal and the stress tensor therefore has to satisfy a local relation
\begin{align}
    \bm{x}\times\nabla\cdot\bm{\sigma}=\nabla\cdot(\bm{x}\times\bm{\sigma}).
\end{align}
An expansion of this condition in Cartesian coordinates shows that it is satisfied if and only if the stress tensor is symmetric
\begin{align}
    \bm{\sigma}^T=\bm{\sigma}.
\end{align}

\section{Conservation of energy}
A differential volume $dV$ of fluid contains a differential mass $\rho dV$ and therefore a differential amount of kinetic energy $\tfrac{1}{2}\rho|\bm{v}|^2dV$. Using the identity~\eqref{eq:rho-identity}, the rate of change of fluid's kinetic energy in a volume $V(t)$ moving with the fluid becomes
\begin{align}
    \frac{d}{dt}\int_{V(t)}\frac{1}{2}\rho|\bm{v}|^2dV=\int_{V(t)}\rho\frac{D}{Dt}\frac{1}{2}|\bm{v}|^2dV=\int_{V(t)}\rho\bm{v}\cdot\frac{D\bm{v}}{Dt}dV.
\end{align}
Using the equation~\eqref{eq:cauchy-momentum} for velocity, and the vector identity $\bm{v}\cdot\nabla\cdot\bm{\sigma}=\nabla\cdot(\bm{\sigma}\cdot\bm{v})-\bm{\sigma}:\nabla\bm{v}$, one finds
\begin{align}
    \int_{V(t)}\rho\bm{v}\cdot\frac{D\bm{v}}{Dt}dV&=\int_{V(t)}\left(\bm{v}\cdot\rho\bm{g}+\nabla\cdot(\bm{\sigma}\cdot\bm{v})-\bm{\sigma}:\nabla\bm{v}\right) dV.
\end{align}
As the initial choice for volume $V(t)$ is arbitraty, a local equation for the kinetic energy density is found
\begin{align}
    \rho\frac{D}{Dt}\frac{1}{2}|\bm{v}|^2=\bm{v}\cdot\rho\bm{g}+\nabla\cdot(\bm{\sigma}\cdot\bm{v})-\bm{\sigma}:\nabla\bm{v}=\rho \bm{v}\cdot\frac{D\bm{v}}{Dt}.
\end{align}
Then using the divergence theorem to convert the volume integral of $\nabla\cdot(\bm{\sigma}\cdot\bm{v})$ into a surface integral, one identifies the driving mechanisms for the rate of change of fluid's kinetic energy
\begin{align}\label{eq:material-rate-of-change-of-kinetic-energy}
    \underbrace{\frac{d}{dt}\int_{V(t)}\frac{1}{2}\rho|\bm{v}|^2dV}_{\substack{\text{Rate of change}\\\text{of kinetic energy}}}&=\int_{V(t)}\bm{v}\cdot\rho\bm{g}dV+\int_{\partial V(t)}\bm{v}\cdot\bm{\sigma}\cdot\bm{\hat{n}}dA-\int_{V(t)}\bm{\sigma}:\nabla\bm{v}dV\nonumber\\
    &=\underbrace{\int_{V(t)}\bm{v}\cdot\delta\bm{F}_b}_{\substack{\text{Rate of mechanical}\\\text{work performed}\\\text{by body forces}}}+\underbrace{\int_{\partial V(t)}\bm{v}\cdot\delta\bm{F}_s}_{\substack{\text{Rate of mechanical}\\\text{work performed}\\\text{by surface forces}}}-\underbrace{\int_{V(t)}\bm{\sigma}:\nabla\bm{v}dV}_{\substack{\text{Rate of mechanical}\\\text{energy change from}\\\text{local deformation}}},
\end{align}
The rate of change of the fluid's kinetic energy therefore does not equal simply the rate of work done by the body and surface forces but includes also an additional volumetric term originating from the rate of deformation of the fluid's local velocity field.

Because of the microscopic internal structure, fluid has internal energy $u(\bm{x},t)$ per unit mass, and a differential volume $dV$ contains a differential amount $\rho u dV$ of internal energy. The total energy the fluid contains is the sum of kinetic and internal energy. The rate at which the total energy changes is driven by the rate at which the fluid receives heat or yields it to its surroundings though it's boundary via conduction or internally by absorbing or emitting radiation, and by the rate at which the body and surface forces perform mechanical work on the fluid
\begin{align}\label{eq:total-energy-balance-material-volume}
    \underbrace{\frac{d}{dt}\int_{V(t)}(\rho(u+\frac{1}{2}|\bm{v}|^2)dV}_{\substack{\text{Rate of change of kinetic}\\\text{and internal energy}}}
    &=
    \underbrace{\int_{\partial V(t)}\bm{q}\cdot(-\bm{\hat{n}})dA}_{\substack{\text{Net rate of heat received}\\\text{through the boundary}}}+\underbrace{\int_{V(t)}q_VdV}_{\substack{\text{Net rate of heat}\\\text{absorbed}\\\text{volumetrically}}}\nonumber\\
    &+\underbrace{\int_{V(t)}\bm{v}\cdot\delta\bm{F}_b}_{\substack{\text{Rate of mechanical}\\\text{work performed}\\\text{by body forces}}}+\underbrace{\int_{\partial V(t)}\bm{v}\cdot\delta\bm{F}_s}_{\substack{\text{Rate of mechanical}\\\text{work performed}\\\text{by surface forces}}}
\end{align}
Here the vector $\bm{q}(\bm{x},t)$ describes the local conductive heat flux and the scalar $q_V(\bm{x},t)$ describes the local net rate per unit volume at which the fluid absorbs heat from radiation passing through it and emits thermal radiation. The negative sign in the heat rate over the surface accounts for the sufrace normal $\bm{\hat{n}}$ pointing outwards from the volume $V(t)$. Subtracting the rate of change of kinetic energy~\eqref{eq:material-rate-of-change-of-kinetic-energy}, the rate of change of fluids internal energy becomes 
\begin{align}\label{eq:energy-balance-material-volume}
    \underbrace{\frac{d}{dt}\int_{V(t)}\rho udV}_{\substack{\text{Rate of change}\\\text{of internal energy}}}=\underbrace{\int_{\partial V(t)}\bm{q}\cdot(-\bm{\hat{n}})dA}_{\substack{\text{Net rate of heat received}\\\text{through the boundary}}}+\underbrace{\int_{V(t)}q_VdV}_{\substack{\text{Net rate of heat}\\\text{absorbed}\\\text{volumetrically}}}+\underbrace{\int_{V(t)}\bm{\sigma}:\nabla\bm{v}dV}_{\substack{\text{Rate of mechanical}\\\text{energy change from}\\\text{local deformation}}}.
\end{align}
Using the identity~\eqref{eq:rho-identity} for the time derivative on the left
\begin{align}
    \frac{d}{dt}\int_{V(t)}\rho udV=\int_{V(t)}\rho\frac{Du}{Dt}dV,
\end{align}
the divergence theorem for the heat flux term on the right
\begin{align}
    \int_{\partial V(t)}\bm{q}\cdot(-\bm{\hat{n}})dA=\int_{V(t)}(-\nabla\cdot\bm{q})dV,
\end{align}
and requiring energy balance~\eqref{eq:energy-balance-material-volume} to hold regardless of the initial choice for the volume $V(t)$, one finds the local partial differential equation for the specific internal energy
\begin{align}\label{eq:internal-energy-equation}
    \rho\frac{Du}{Dt}=-\nabla\cdot\bm{q}+q_V+\bm{\sigma}:\nabla\bm{v}.
\end{align}
\begin{example}
    The local kinetic-energy equation obtained above shows that changes in
    kinetic energy are caused by mechanical work performed by body and
    surface forces,
    \begin{align*}
        \rho\frac{D}{Dt}\frac{1}{2}|\bm{v}|^2
        =
        \bm{v}\cdot(\rho\bm{g}+\nabla\cdot\bm{\sigma}).
    \end{align*}
    In a still fluid, $\bm{v}=0$ and hence $D/Dt=\partial_t$ and
    $\nabla\bm{v}=0$. In such case, the kinetic energy therefore cannot change, while the equation for internal energy reduces to
    \begin{align*}
        \rho\partial_t u=-\nabla\cdot\bm{q}+q_V.
    \end{align*}
    Changes in internal energy can thus occur through heat transfer even
    when the fluid has no macroscopic motion. The coupling between
    macroscopic mechanical and microscopic internal energy occurs through
    the deformation of the fluid, encoded in the term $\bm{\sigma}:\nabla\bm{v}$.
\end{example}

\section{Constitutive relations for stress and heat flux}
The conservation principles of mass, momentum, and energy discussed above lead to local differential equations from which the local time rate of change of fluid's density, velocity, and internal energy ought to be computed. To evaluate these expressions, additional relations are needed to describe how stress $\bm{\sigma}$ and heat flux $\bm{q}$ depend on the local state and motion of the fluid. Such relations are called constitutive. The form of a constitutive relation is restricted by physical principles such as invariance with respect to the choice of reference frame and the principle of isotropy in a still fluid, while the remaining constitutive behavior must typically be determined experimentally.

For non-relativistic motion obeying Newton's laws, the principle of Galilean invariance requires the laws of motion to be invariant in all inertial frames. A uniform translational $\bm{v} \to \bm{v} + \bm{v}_0$ between two inertial frames moving with relative velocity $\bm{v}_0$ with respect to each other must therefore leave the momentum equation~\eqref{eq:cauchy-momentum} invariant. As the material derivative of a constant velocity vector is zero, $D\bm{v}_0/Dt=\bm{0}$, the stress tensor must also remain invariant under such translation. Otherwise additional terms would be introduced into the Cauchy's momentum equation describing non-physical forces. This requirement with respect to all constant velocities $\bm{v}_0$ forces the stress to be independent of the absolute velocity $\bm{v}$. For a simple fluid, whose stress is assumed to be determined locally by the thermodynamic state and the local velocity gradient, its constitutive dependence may consequently be written as
\begin{align}
    \bm{\sigma} = \bm{\sigma}(\rho, p, T, \nabla\bm{v}).
\end{align}
In principle, the stress could depend also on higher derivatives of the velocity field, and on the spatial derivatives of the thermodynamic state variables, but experimental observations of simple fluids can be explained to good accuracy with stress assumed to depend on the thermodynamic state and the velocity gradient only. 

Consider fluid confined in a cylindrical annulus. Rotate both the inner and outer walls of the annulus with the same and constant angular velocity in the same direction as in Fig.~\ref{fig:cylinder-flow-rotation}. The fluid within is observed to rotate with exactly the same angular velocity at every point. In other words, the fluid undergoes rigid rotation. If one were to observe the fluid from a frame rotating at the same angular velocity, the fluid would appear perfectly still, regardless of the magnitude of the angular velocity. Resulting centrifugal force would introduce a pressure gradient in the radial direction the same way gravity normally induces hydrostatic pressure, but no other forces would appear measurable on the cylindrical walls. The stress must therefore be independent of the rigid rotation of the observer's frame, implying that the antisymmetric part $\bm{\Omega}$ of the velocity gradient tensor cannot contribute to the stress. Consequently, the dependency of stress on the velocity gradient can appear only through the symmetric traceless tensor $\bm{D}$ describing the rate of shear in non-isotropic deformation and the isotropic part $\tfrac{1}{3}(\nabla\cdot\bm{v})\bm{I}$ describing the rate of volumetric deformation. In an undeforming fluid, the stress is isotropic and represented by the thermodynamic pressure. The stress can consequently be decomposed as
\begin{align}\label{eq:stress}
    \bm{\sigma} = -p\bm{I} + \bm{\tau}(\rho, p, T, \bm{D}, \nabla\cdot\bm{v}),
\end{align}
with $\bm{\tau}$ referred to as the deviatoric stress tensor. The minus sign for pressure accounts for the outward unit normal of the surface on which the traction must counter the pressure in undeforming fluid.
\begin{figure}[!htb]
    \centering
    \input{\rootfolder/code/fig-rotation.tex}
    \caption{Fluid flow in cylindrical annulus with walls co-rotating with the same angular velocity. Only rigid rotation characterized by $\bm{\Omega}$ is present, and no torque on the walls is required to maintain the fluid flow.}
    \label{fig:cylinder-flow-rotation}
\end{figure}

If only the outer wall of the cylindrical annulus is spun as in Fig.~\ref{fig:cylinder-flow-shear}, or both walls spun but in opposite directions, shear is observed to develop into the flow, and a constant torque at the annulus walls proportional to the shear rate is required to maintain the flow. This indicates that, for sufficiently small deformation rates, the deviatoric stress depends approximately linearly on $\bm{D}$. Studying then the travel of sound waves in a pipe as illustrated in Fig.~\ref{fig:pipe-wave}, measurements of the wave amplitude attenuation over distance can be associated with a traction force caused by volumetric deformation. In this case as well, the force deduced to have a linear dependency on the volumetric deformation rate characterized by $\nabla\cdot\bm{v}$. Together, the measurements inducate the following expression for the deviatoric stress tensor
\begin{align}\label{eq:deviatoric-stress}
    \bm{\tau} = 2\mu \bm{D} + \zeta(\nabla \cdot \bm{v})\bm{I},
\end{align}
which is symmetric and therefore compatible with the requirement of angular momentum conservation. Fluids obeying this linear relationship are classified as Newtonian fluids, with the viscous proportionality coefficients $\mu$ and $\zeta$ being properties of the fluid and generally functions of the thermodynamic state $(\rho, p, T)$. In contrast, non-Newtonian fluids exhibit a more complex, non-linear dependency on the rates of deformation.
\begin{figure}[!htb]
    \centering
    \input{\rootfolder/code/fig-shear.tex}
    \caption{Fluid flow in a cylindrical annulus with fixed inner and rotating outer cylindrical walls. Constant torque is required on the walls to maintain the flow, proportional to the shear rate characterized by the tensor $\bm{D}$.}
    \label{fig:cylinder-flow-shear}
\end{figure}
\begin{figure}[!htb]
    \centering
    \input{\rootfolder/code/fig-wave.tex}
    \caption{A traveling wave in a fluid. The flow is characterized by $\nabla\cdot\bm{v}$, related to the volumetric deformation. The deformation attenuates in a long pipe, decreasing the wave amplitude.}
    \label{fig:pipe-wave}
\end{figure}

The heat flux $\bm{q}$ describes the local transfer of energy mediated by microscopic motion of particles relative to the local motion of the fluid, i.e., the transfer of heat. Expressing it in terms of the fluid's variables requires a constitutive relation of its own that is determined experimentally. If a fluid is confined between two surfaces that are maintained at different temperatures, heat is observed to flow from the warmer surface toward the colder surface. Reversing the temperatures on the plates reverses the direction of the heat transfer. If the plates are oriented so that temperature-induced density changes do not couple with gravity to drive buoyancy-related flow, then the fluid remains still and the heat transfer occurs by conduction, as in solids. For sufficiently small temperature gradients, the heat flux is observed to be approximately proportional to the temperature gradient, providing the constitutive relation
\begin{align}
    \bm{q}=-\lambda\nabla T,
    \label{eq:fourier-law}
\end{align}
where $\lambda$ is the thermal conductivity of the fluid with a value that depends on the thermodynamic state of the fluid. The negative sign indicates that heat is transported from high temperature toward low temperature. The thermal conductivity, together with the viscosity coefficients $\mu$ and $\zeta$, is referred to as a transport coefficient. These coefficients depend on the fluid and its thermodynamic state and are generally determined experimentally.

\section{Summary}
The motion of a fluid can be described by following individual points along the flow and studying how differential line, surface, and volume elements evolve. The velocity gradient tensor $\nabla\bm{v}$ describes the local evolution of these differential elements and can be decomposed into isotropic deformation, volume-preserving deformation, and rotation. The Leibniz integral differentiation rule then provides the connection
between the local evolution of the fluid and conservation laws applied over finite moving volumes.

Conservation of mass, momentum, and energy leads to the local equations
\begin{align}
    \label{eq:density-collection}
    \frac{1}{\rho}\frac{D\rho}{Dt}&= -\nabla\cdot\bm{v},\\
    \label{eq:momentum-collection}
    \rho\frac{D\bm{v}}{Dt}&= \rho\bm{g}+\nabla\cdot\bm{\sigma},\\
    \label{eq:internal-energy-collection}
    \rho\frac{Du}{Dt}&= -\nabla\cdot\bm{q}+q_V+\bm{\sigma}:\nabla\bm{v},
\end{align}
where the scalar-valued differential operator
\begin{align}
    \frac{D}{Dt}
    &= \frac{\partial}{\partial t}+\bm{v}\cdot\nabla
\end{align}
is the material derivative. Conservation of angular momentum further requires the stress tensor to be symmetric $\bm{\sigma}^T=\bm{\sigma}$. For a simple, Newtonian fluid, the stress and heat flux are described by the constitutive relations
\begin{align}
    \label{eq:sigma-collection}
    \bm{\sigma}&=-p\bm{I}+\bm{\tau},\\
    \label{eq:tau-collection}
    \bm{\tau}&=2\mu\bm{D}+\zeta(\nabla\cdot\bm{v})\bm{I},\\
    \label{eq:heat-flux-collection}
    \bm{q}&=-\lambda\nabla T,
\end{align}
where the tensor
\begin{align}
    \label{eq:D-tensor-collection}
    \bm{D}=\frac{1}{2}\left(\nabla\bm{v}+(\nabla\bm{v})^{\mathsf T}\right)-\frac{1}{3}(\nabla\cdot\bm{v})\bm{I},
\end{align}
describes the non-isotropic rate of deformation. The coefficients $\mu$ and $\zeta$ are referred to as shear and bulk viscosity, respectively, while $\lambda$ is the heat conductivity.

The differential equations describe the evolution of density, velocity, and specific internal energy once the fluid state and its spatial gradients are known. Assuming an equation of state of the form $p(\rho,T)$ to exist, the equations provide a calculation method
\begin{align}
    \rho,\, T,\, u,\, \bm{v} \implies \partial_t\rho,\, \partial_t\bm{v},\, \partial_tu.
\end{align}
However, the equations are not closed. Four independent variables exists while time evolution is provided only for three of them, with the time evolution for temperature $\partial_tT$ missing. A thermodynamic relation $u(\rho,T)$ connecting the internal energy, pressure, temperature, and density is therefore needed. Establishing this thermodynamic closure is the subject of the next chapter.

\chapter{Thermodynamic closure}
In this chapter, the fluid motion is restricted to very specific type of flows which enables associating the derivatives of the internal energy with respect to the state variables to experimental data measured during heating the substance. An observation of the existence of inaccessible states without addition or removal of heat then together with mathematics establish the existence of entropy as a state function, revealing the thermodynamic framework. Once the connection of state variables to internal energy and entropy is revealed, thermodynamic potential functions enable calculation of all thermodynamic data and form the foundation for thermodynamic data libraries such as \verb|pyromat| and \verb|coolprop|. Simultaneously the framework closes the fluid equations of motion into a set of well-defined partial differential equations.

\section{Learning goals}
\begin{itemize}
    \item[\elementary] The student lists conditions assumed in quasistatic processes.
    \item[\elementary] The student evaluates specific heats, internal energy, enthalpy, and entropy values for a given state point using \verb|pyromat| or \verb|coolprop|.
    \item[\elementary] The student calculates internal energy, enthalpy, and entropy changes for an ideal gas and compares the changes to values from \verb|pyromat| or \verb|coolprop|.
    \item[\intermediate] The student manipulates differential expressions to express heat and expansion work differentials in a desired set of state variables.
    \item[\intermediate] The student numerically integrates expansion work and heat along given paths and verifies their path dependency.
    \item[\intermediate] The student uses differential relations to calculate changes in thermodynamic state functions for a non-ideal fluid.
    \item[\advanced] The student infers why specific internal energy and specific entropy are state functions from the path dependence of heat and work and the existence of inaccessible states, respectively.
    \item[\advanced] The student calculates thermodynamic properties from Helmholtz or Gibbs free energy functions and their derivatives.
    \item[\advanced] The student programs a python script utilizing \verb|pyromat| or \verb|coolprop| to create a $p$-$h$ or a $T$-$s$ diagram.
\end{itemize}

% \section{Rearrangement of the energy balance}
% As the first step in constructing a link between fluid mechanics and thermodynamics, define the experimentally-measurable volume-averaged density, temperature, and pressure as integrals over the material volume $V(t)$
% \begin{align}
%     \bar{\rho}(t)&=\frac{1}{V(t)}\int_{V(t)}\rho dV=\frac{m}{V(t)},\\
%     \bar{T}(t)&=\frac{1}{V(t)}\int_{V(t)}TdV,\\
%     \bar{p}(t)&=\frac{1}{V(t)}\int_{V(t)}pdV,
% \end{align}
% where $m=\int_{V(t)}\rho dV$ is the mass of the fluid the volume contains and constant in time. In terms of the Neumann Green's function $G_N$, one then has the identities
% \begin{align}
%     \rho(\bm{x},t)-\bar{\rho}(t)&=\int_{V'(t)}\nabla'G_N(\bm{x},\bm{x}',t)\cdot\nabla'\rho(\bm{x}',t) dV',\\
%     T(\bm{x},t)-\bar{T}(t)&=\int_{V'(t)}\nabla'G_N(\bm{x},\bm{x}',t)\cdot\nabla'T(\bm{x}',t) dV',\\
%     p(\bm{x},t)-\bar{p}(t)&=\int_{V'(t)}\nabla'G_N(\bm{x},\bm{x}',t)\cdot\nabla'p(\bm{x}',t) dV',
% \end{align}
% highlighting that the local deviation from the volume averaged value relates to the field gradient in the entire volume, and is therefore bounded by the maximum value of the gradient within the volume. 
% \begin{example}
%     Define the Neumann Green's function $G_N(\bm{x},\bm{x}',t)$ as the solution to the time-dependent boundary value problem
%     \begin{align*}
%         \nabla^2G_N(\bm{x},\bm{x}',t)&=\frac{1}{V(t)}-\delta^3(\bm{x}-\bm{x}'), & &\text{for} \quad \bm{x},\bm{x}'\in V(t)\\
%         \nabla G_N(\bm{x},\bm{x}',t)\cdot\bm{\hat{n}}(\bm{x})&=0, & &\text{for}\quad \bm{x}\in\partial V(t).
%     \end{align*}
%     Use the properties of the Dirac delta-function, the divergence theorem, and the boundary condition of $G_N$ to express the local deviation of any smooth scalar field $f(\bm{x},t)$ from its volume average in terms of the field gradient
%     \begin{align*}
%         f(\bm{x},t)-\frac{1}{V(t)}\int_{V(t)}f(\bm{x}',t)dV'&=\int_{V(t)}f(\bm{x}',t)\left(\delta^3(\bm{x}'-\bm{x})-\frac{1}{V(t)}\right)dV'\\
%         &=-\int_{V(t)}f(\bm{x}',t)\nabla'^2G_N(\bm{x}',\bm{x},t)dV'\\
%         &=-\int_{\partial V(t)}f(\bm{x}',t)\nabla'G_N(\bm{x}',\bm{x},t)\cdot\bm{\hat{n}}(\bm{x}',t)dA\\
%         &\qquad+\int_{V(t)}\nabla'G_N(\bm{x}',\bm{x},t)\cdot\nabla'f(\bm{x}',t)dV'\\
%         &=\int_{V(t)}\nabla'G_N(\bm{x}',\bm{x},t)\cdot\nabla'f(\bm{x}',t)dV'.
%     \end{align*}
% \end{example}

% Now turn the attention to the terms in the fluid's energy balance~\eqref{eq:energy-balance-material-volume}. For the term involving stress, use $\bm{I}:\nabla\bm{v}=\nabla\cdot\bm{v}$, $\bm{I}:\bm{I}=3$, and $\bm{D}:\bm{I}=0$ to express
% \begin{align}
%     \bm{\sigma}:\nabla\bm{v}=-p\nabla\cdot\bm{v}+2\mu\bm{D}:\bm{D}+\zeta(\nabla\cdot\bm{v})^2.
% \end{align}
% In the first term, the divergence of the velocity is related to the time derivative of the volume-averaged density $\bar{\rho}$. According to the Leibniz integral differentiation rule, one finds
% \begin{align}
%     \int_{V(t)}\nabla\cdot\bm{v}\,dV=\frac{dV(t)}{dt}=m\frac{d\bar{v}(t)}{dt}=-\frac{m}{\bar{\rho}^2}\frac{d\bar{\rho}(t)}{dt},
% \end{align}
% with $\bar{v}(t)=1/\bar{\rho}(t)=V(t)/m$ the volume-averaged specific volume. Formally splitting $p=\bar{p}+(p-\bar{p})$, one may then divide the pressure contribution to the stress integral in the energy balance into averaged quantities and gradients-driven term
% \begin{align}
%     \int_{V(t)}p\nabla\cdot\bm{v}\,dV=m\bar{p}\frac{d\bar{v}}{dt}+\int_{V(t)}\int_{V'(t)}(\nabla\cdot\bm{v})\nabla'G_{N}\cdot\nabla'pdV'dV.
% \end{align}
% Denote the mechanical rate of work per unit mass done by the fluid's average pressure with
% \begin{align}\label{eq:dot-w_p-expressions}
%     \dot{\bar{w}}_p=\bar{p}\frac{d\bar{v}}{dt}=-\frac{\bar{p}}{\bar{\rho}^2}\frac{d\bar{\rho}}{dt},
% \end{align}
% and the mechanical rate of work done by pressure deviations and viscous forces within the fluid with the quadratically gradient-dependent functionals
% \begin{align}
%     \dot{\mathcal{W}}_{\delta p}[\nabla p,\nabla\bm{v}]&=\int_{V(t)}\int_{V'(t)}(\nabla\cdot\bm{v})\nabla'G_{N}\cdot\nabla'pdV'dV,\\
%     \dot{\mathcal{W}}_{\mu,\zeta}[\rho,T,p,\nabla\bm{v}]&=-\int_{V(t)}\left(2\mu\bm{D}:\bm{D}+\zeta(\nabla\cdot\bm{v})^2\right)dV.
% \end{align}
% Note that the viscous work-rate functional generally depends on the fluid state variables through the viscosity coefficients. While the pressure-deviation related work rate $\dot{\mathcal{W}}_{\delta p}$ can be positive or negative, the viscous work rate $\dot{\mathcal{W}}_{\mu,\zeta}$ is always non-positive because $\bm{D}:\bm{D}=\sum_{ij}D_{ij}^2\geq 0$ and $(\nabla\cdot\bm{v})^2\geq 0$, and the experimentally determined viscous coefficients $\mu$ and $\zeta$ are non-negative. The stress-related term in the energy balance can now be expressed
% \begin{align}
%     \int_{V(t)}\bm{\sigma}:\nabla\bm{v}dV=-m\dot{\bar{w}}_p-\dot{\mathcal{W}}_{\delta p}-\dot{\mathcal{W}}_{\mu,\zeta}.
% \end{align}

% As the final step in reformatting the energy balance, define the density-weighted volume-averaged specific internal energy over the volume $V(t)$ as
% \begin{align}
%     \bar{u}(t)=\frac{1}{m}\int_{V(t)}\rho u dV,
% \end{align}
% and denote the rate of heat received by the fluid with the heat rate functional
% \begin{align}
%     \dot{\mathcal{Q}}[\rho,T,p,\nabla T]&=\int_{\partial V(t)}\lambda \nabla T\cdot\bm{\hat{n}}dA+\int_{V(t)}q_VdV.
% \end{align}
% In addition to the temperature gradient, the heat rate functional generally depends also on the fluid's state variables through the heat conductivity coefficient $\lambda$ and the local heat absorbtion rate $q_V$. The energy balance~\eqref{eq:energy-balance-material-volume} over the volume $V(t)$ can now be written as
% \begin{align}
%     \underbrace{m\frac{d\bar{u}}{dt}}_{\substack{\text{Rate of}\\\text{internal}\\\text{energy}\\\text{change}}}=\underbrace{\dot{\mathcal{Q}}}_{\substack{\text{Net rate}\\\text{of heat}\\\text{into the}\\\text{fluid}}}-\underbrace{m\dot{\bar{w}}_p}_{\substack{\text{Rate of}\\\text{work by}\\\text{average}\\\text{pressure}}}-\underbrace{\dot{\mathcal{W}}_{\delta p}}_{\substack{\text{Rate of}\\\text{work by}\\\text{pressure}\\\text{variations}}}-\underbrace{\dot{\mathcal{W}}_{\mu,\zeta}}_{\substack{\text{Rate of}\\\text{work by}\\\text{viscous}\\\text{forces}}}.
% \end{align}

\section{Rearrangement of the energy balance}
The conservation laws investigated over an arbitrary material volume provide local equations for the time rate of change of density, velocity, and specific internal energy. These equations involve four local quantities, ($\rho,T,u,\bm v$), but provide evolution equations for only three of them ($\rho,u,\bm v$). An additional relation connecting the internal energy to the thermodynamic state is therefore required to close the equations. The purpose of this section is to rearrange the energy balance to enable construction of such a relation from experimentally measurable data.

The energy balance over a material volume $V(t)$ links heat rate and mechanical deformation work rate to the rate at which internal energy of the volume changes
\begin{align}
    \underbrace{\frac{d}{dt}\int_{V(t)}\rho udV}_{\substack{\text{Rate of change}\\\text{of internal energy}}}=\underbrace{\int_{\partial V(t)}\bm{q}\cdot(-\bm{\hat{n}})dA}_{\substack{\text{Net rate of heat received}\\\text{through the boundary}}}+\underbrace{\int_{V(t)}q_VdV}_{\substack{\text{Net rate of heat}\\\text{absorbed}\\\text{volumetrically}}}+\underbrace{\int_{V(t)}\bm{\sigma}:\nabla\bm{v}dV}_{\substack{\text{Rate of mechanical}\\\text{energy change from}\\\text{local deformation}}}.
\end{align}
As the mass $m=\int_{V(t)}\rho dV$ of the material volume $V(t)$ is constant in time, the time rate of change of the total internal energy on the left side of the energy balance becomes
\begin{align}
    \frac{d}{dt}\int_{V(t)}\rho udV=m\frac{d\bar{u}}{dt},
\end{align}
where $\bar{u}$ is the density-weighted volume average of the local internal energy field, defined
\begin{align}
    \bar{u}=\frac{1}{m}\int_{V(t)}\rho u dV.
\end{align}
Associating the heat and mechanical deformation work rates on the right of the energy balance with the values and time rates change of volume-averaged state variables and excess terms, provides a link between internal energy and experimentally measurable state variables. Further inspection of the energy balance at a limit where the excess terms can be considered negligible, and assuming the reduced energy balance to hold regardless of the material volume chosen, provides a relation for internal energy in terms of state variables, which can then be interpreted as a local thermodynamic relation.

This process begins with associating local field values to volume averages. Over a time-dependent material volume $V(t)$, the Neumann Green's function $G_N(\bm{x},\bm{x}',t)$ is defined as the solution to the time-dependent boundary-value problem
\begin{align}
    \nabla'^2G_N(\bm{x}',\bm{x},t)
    &=
    \frac{1}{V(t)}-\delta^3(\bm{x}'-\bm{x}),
    &&
    \bm{x},\bm{x}'\in V(t),\\
    \nabla'G_N(\bm{x}',\bm{x},t)\cdot\bm{\hat n}(\bm{x}',t)
    &=0,
    &&
    \bm{x}'\in\partial V(t).
\end{align}
For a sufficiently smooth scalar field $f(\bm{x},t)$, the function $G_N$ provides the following identity between the function value and it's volume average $\bar{f}(t)=(1/V(t))\int_{V(t)}f(\bm{x},t)dV$
\begin{align}
    f(\bm{x},t)-\bar{f}(t)
    =
    \int_{V(t)}
    \nabla'G_N(\bm{x}',\bm{x},t)
    \cdot
    \nabla'f(\bm{x}',t)\,dV'.
    \label{eq:field-average-deviation}
\end{align}
According to the identity, the local deviation of a scalar field from its volume averaged value relates to the field gradient in the entire volume, and is bounded by the maximum value of the gradient within the volume and a geometric factor
\begin{align}
    |f(\bm{x},t)-\bar{f}(t)|\leq \int_{V(t)}|\nabla G_N|dV\,\max_{V(t)}|\nabla f|.
\end{align}
\begin{example}
    Verifying the identity~\eqref{eq:field-average-deviation} proceeds by multiplying the defining equation for $G_N$ and by $f$ and integrating it over the volume:
    \begin{align*}
        f(\bm{x},t)-\frac{1}{V(t)}
        \int_{V(t)}f(\bm{x}',t)\,dV'
        &=
        \int_{V(t)}
        f(\bm{x}',t)
        \left[
            \delta^3(\bm{x}'-\bm{x})-\frac{1}{V(t)}
        \right]dV'\\
        &=
        -\int_{V(t)}
        f(\bm{x}',t)
        \nabla'^2G_N(\bm{x}',\bm{x},t)\,dV'\\
        &=
        -\int_{\partial V(t)}
        f(\bm{x}',t)
        \nabla'G_N(\bm{x}',\bm{x},t)
        \cdot\bm{\hat n}(\bm{x}',t)\,dA\\
        &\qquad
        +\int_{V(t)}
        \nabla'G_N(\bm{x}',\bm{x},t)
        \cdot\nabla'f(\bm{x}',t)\,dV'\\
        &=
        \int_{V(t)}
        \nabla'G_N(\bm{x}',\bm{x},t)
        \cdot\nabla'f(\bm{x}',t)\,dV'.
    \end{align*}
    On the second to last line, the boundary term vanishes because of the normal derivative of $G_N$ vanishes on the boundary.
\end{example}

With the help of the Neumann Green's function, one is able to separate the mechanical deformation work rate in the energy balance into a term expressible in terms of volume-averaged state variables and their rates, and excess terms. For this to work, use first the identities $\bm{I}:\nabla\bm{v}=\nabla\cdot\bm{v}$, $\bm{I}:\bm{I}=3$, and $\bm{D}:\bm{I}=0$ to express the term $\bm{\sigma}:\nabla\bm{v}$ as
\begin{align}
    \bm{\sigma}:\nabla\bm{v}
    =
    -p\nabla\cdot\bm{v}
    +2\mu\bm{D}:\bm{D}
    +\zeta(\nabla\cdot\bm{v})^2,
    \label{eq:stress-work-decomposition}
\end{align}
so that the mechanical deformation work rate in the energy balance is split into pressure and viscosity related contributions
\begin{align}
    \underbrace{\int_{V(t)}\bm{\sigma}:\nabla\bm{v}dV}_{\substack{\text{Mechanical deformation}\\\text{work rate}}}=-\underbrace{\int_{V(t)}p\nabla\cdot\bm{v}dV}_{\substack{\text{Pressure related}\\\text{deformation work rate}}}+\underbrace{\int_{V(t)}\left(2\mu\bm{D}:\bm{D}
    +\zeta(\nabla\cdot\bm{v})^2\right)dV}_{\substack{\text{Viscosity related}\\\text{deformation work rate}}}
\end{align}
Next use the Leibniz rule to realize that the divergence of the velocity is related to the rate of change of volume-averaged specific volume or density via
\begin{align}
    \int_{V(t)}\nabla\cdot\bm{v}\,dV
    =
    \frac{dV(t)}{dt}
    =
    m\frac{d\bar v(t)}{dt}
    =
    -\frac{m}{\bar{\rho}^2}
    \frac{d\bar{\rho}(t)}{dt},
\end{align}
where $\bar v(t)=V(t)/m$ is the volume-averaged specific volume and $\bar{\rho}(t)=1/\bar{v}(t)$ the corresponding volume-averaged density. If the pressure were constant over the volume $V(t)$, it could be directly pulled out of the integral of $p\nabla\cdot\bm{v}$, resulting in a term expressible in volume averages only. Generally this is not true, though, and a trick involving the Neumann Green's function is needed. For the trick, define the volume-averaged pressure
\begin{align}
    \bar{p}(t)=\frac{1}{V(t)}\int_{V(t)}p\,dV,
\end{align}
with the deviation in the local value from the volume average $\bar{p}$ being expressible in terms of the Neumann Green's function
\begin{align}
    p(\bm{x},t)-\bar{p}(t)=\int_{V(t)}\nabla'G_N(\bm{x}',\bm{x},t)\cdot\nabla'p(\bm{x}',t)\,dV'.
\end{align}
Then formally decompose the local pressure into average and a pressure deviation according to $p=\bar p+(p-\bar p)$. The average part can now be pulled out of the integral while the deviation provides a gradient-related correction term, resulting in
\begin{align}
    \underbrace{\int_{V(t)}p\nabla\cdot\bm v\,dV}_{\substack{\text{Pressure related}\\\text{deformation work rate}}}
    &=
    \bar p\int_{V(t)}\nabla\cdot\bm v\,dV
    +
    \int_{V(t)}(p-\bar p)\nabla\cdot\bm v\,dV\\
    &=
    \underbrace{m\bar p\frac{d\bar v}{dt}}_{\substack{\text{Average}\\\text{pressure}\\\text{related}\\\text{deformation}\\\text{work rate}}}
    +
    \underbrace{\int_{V(t)}
    \int_{V'(t)}
    (\nabla\cdot\bm v)
    \nabla'G_N\cdot\nabla'p\,
    dV'\,dV}_{\substack{\text{Local pressure deviation related}\\\text{deformation work rate}}}.
\end{align}

Denote the mass-specific average pressure deformation work rate, the pressure variation related deformation work rate, and the viscous contribution to the deformation work rate with
\begin{align}
    \label{eq:dot-w_p-expressions}
    \dot{\bar w}_p
    &=
    \bar p\frac{d\bar v}{dt}
    =
    -\frac{\bar p}{\bar{\rho}^2}
    \frac{d\bar{\rho}}{dt},\\
    \dot{\mathcal W}_{\delta p}
    &=
    \int_{V(t)}
    \int_{V'(t)}
    (\nabla\cdot\bm v)
    \nabla'G_N\cdot\nabla'p\,
    dV'\,dV,\\
    \dot{\mathcal W}_{\mu,\zeta}
    &=
    -\int_{V(t)}
    \left(
        2\mu\bm D:\bm D
        +
        \zeta(\nabla\cdot\bm v)^2
    \right)dV,
\end{align}
and the total deformation work rate contributing to the energy balance is then expressed
\begin{align}
    \int_{V(t)}
    \bm{\sigma}:\nabla\bm v\,dV
    =
    -m\dot{\bar w}_p
    -\dot{\mathcal W}_{\delta p}
    -\dot{\mathcal W}_{\mu,\zeta},
    \label{eq:stress-work-averaged}
\end{align}
The deformation work rate is now separated into a contribution expressed entirely in terms of volume-averaged state variables and two correction terms associated with spatial non-uniformity. The pressure-deviation term involves a product of the pressure and velocity gradients, while the viscous contribution is quadratic in the velocity gradient. Furthermore, since $\bm{D}:\bm{D}=\sum_{ij}D_{ij}^2\geq 0$, $(\nabla\cdot\bm{v})^2\geq 0$, and the viscous coefficients $\mu$ and $\zeta$ are non-negative, the viscous deformation work rate $\dot{\mathcal{W}}_{\mu,\zeta}$ is always non-positive. The pressure-deviation contribution $\dot{\mathcal{W}}_{\delta p}$, in contrast, may be positive or negative.
For the heat rate terms on the right of the energy balance, denote the total rate of heat received by the fluid as
\begin{align}
    \dot{\mathcal Q}
    &=
    \int_{\partial V(t)}
    \bm{q}\cdot(-\bm{\hat{n}})dA
    +
    \int_{V(t)}q_V\,dV=
    \int_{\partial V(t)}
    \lambda\nabla T\cdot\bm{\hat n}\,dA
    +
    \int_{V(t)}q_V\,dV.
\end{align}
The dependence of this functional on the state variables enters both through the thermal conductivity $\lambda$ and through the local volumetric energy-transfer rate $q_V$. Substituting the above expressions into the material-volume energy balance finally provides
\begin{align}
    \underbrace{m\frac{d\bar u}{dt}}
    _{\substack{\text{Rate of}\\\text{internal}\\\text{energy}\\\text{change}}}
    =
    \underbrace{\dot{\mathcal Q}}
    _{\substack{\text{Net rate}\\\text{of heat}\\\text{into the}\\\text{fluid}}}
    -
    \underbrace{m\dot{\bar w}_p}
    _{\substack{\text{Rate of}\\\text{work by}\\\text{average}\\\text{pressure}}}
    -
    \underbrace{\dot{\mathcal W}_{\delta p}}
    _{\substack{\text{Rate of}\\\text{work by}\\\text{pressure}\\\text{variations}}}
    -
    \underbrace{\dot{\mathcal W}_{\mu,\zeta}}
    _{\substack{\text{Rate of}\\\text{work by}\\\text{viscous}\\\text{forces}}}.
    \label{eq:energy-balance-macroscopic}
\end{align}
The important feature of this decomposition is that the first work term depends only on the macroscopic state variables, whereas the remaining terms are explicitly controlled by spatial gradients that can be assessed in an appropriate limit. As the next step, one this limit is considered to connect the heat rate $\dot{\mathcal{Q}}$ into rates of change of the volume-averaged state variables.


\section{Quasistatic processes}
The decomposition obtained in the previous section shows that the departure from the volume-average scalar description is contained in terms associated with spatial gradients. Quasistatic processes are defined as processes in which these spatial variations relax on time scales much shorter than the time scale on which the volume-averaged state evolves, while thin layers on the boundary of the volume may remain with significant gradients to enable heat transfer into the volume from its surroundings. In such processes, the volume-averaged state approaches the experimentally measured equilibrium equation of state
\begin{align}
    \bar{p}(t)\to p(\bar{\rho}(t),\bar{T}(t)),
\end{align}
and the mechanical work rates related to pressure deviation and viscous effects are small compared to the heat rate, the average pressure work rate, or both
\begin{align}
    \frac{|\dot{\mathcal{W}}_{\delta p}|}{\max(|\dot{\mathcal{Q}}|,m|\dot{\bar{w}}_p|)}\to 0,\qquad \frac{|\dot{\mathcal{W}}_{\mu,\zeta}|}{\max(|\dot{\mathcal{Q}}|,m|\dot{\bar{w}}_p|)}\to 0.
\end{align}
Here the denominator provides the macroscopic energy-transfer scale against which the corrections identified in the previous section are compared.

In experiments involving quasistatic processes, a relation is sought associating the rate of heat the fluid receives to the measured macroscopic state variables and their rates of change
\begin{align}
    \dot{\mathcal{Q}}\to m \dot{\bar{q}}(\bar{\rho},\bar{T},\bar{p},\dot{\bar{\rho}},\dot{\bar{T}},\dot{\bar{p}}).
\end{align}
The heat rate per unit mass, $\dot{\bar{q}}$, is typically observed to be proportional to the rates of change of the macroscopic volume-averaged state variables with proportionality coefficients that depend on the fluid's state. Since the equation of state binds the state variables $\bar{\rho}$ (or $\bar{v}$), $\bar{T}$, and $\bar{p}$ together and only two of them are needed to define the state. Therefore, generally three different representations for the specific heat rate are available
\begin{align}
    \dot{\bar{q}}=\left\{\begin{aligned}
        &\bar{q}_{T,\rho}(\bar{T},\bar{\rho})\frac{d\bar{T}}{dt}+\bar{q}_{\rho,T}(\bar{T},\bar{\rho})\frac{d\bar{\rho}}{dt},\\
        &\bar{q}_{T,p}(\bar{T},\bar{p})\frac{d\bar{T}}{dt}+\bar{q}_{p,T}(\bar{T},\bar{p})\frac{d\bar{p}}{dt},\\
        &\bar{q}_{\rho,p}(\bar{\rho},\bar{p})\frac{d\bar{\rho}}{dt}+\bar{q}_{p,\rho}(\bar{\rho},\bar{p})\frac{d\bar{p}}{dt}.
    \end{aligned}\right.
\end{align}
This experimentally observed linear dependence on the rates will allow the temporal energy balance to be converted into a differential relation on the thermodynamic state space. Considering the differential representation for the equation of state
\begin{align}\label{eq:differential-equation-of-state}
    \frac{1}{\bar{\rho}}\frac{d\bar{\rho}}{dt}=-\beta_p\frac{d\bar{T}}{dt}+\kappa_T\frac{d\bar{p}}{dt},
\end{align}
the proportionality coefficients for $\dot{\bar{q}}$ are observed to relate to each other via
\begin{align}\label{eq:q-Tp-relations}
    \bar{q}_{T,p}&=\bar{q}_{T,\rho}-\bar{\rho}\beta_p\bar{q}_{\rho,T}=-\bar{\rho}\beta_p\bar{q}_{\rho,p},\\
    \label{eq:q-pT-relations}
    \bar{q}_{p,T}&=\bar{\rho}\kappa_T\bar{q}_{\rho,T}=\bar{q}_{p,\rho}+\bar{\rho}\kappa_T\bar{q}_{\rho,p}.
\end{align}

As an example, the coefficient $\bar{q}_{T,\rho}$ represents the heat per unit mass associated with temperature changes in constant-density processes, while $\bar{q}_{T,p}$ represents the heat per unit mass associated with with temperature changes in constant-pressure processes. In terms of the experimental response rate $\dot{\bar{q}}$, the quasistatic limit of the exact material-volume energy balance now reduces to the following equation
\begin{align}\label{eq:quasistatic-energy-balance}
    \frac{d\bar{u}}{dt}=\dot{\bar{q}}-\dot{\bar{w}}_p.
\end{align}
This equation is now expressed in terms of macroscopic volume-averaged state variables and via experimentally measurable response rates. The following sections will discuss exclusively quasistatic processes. With $\bar{\rho}, \bar{T}, \bar{p},$ and $\bar{u}$ understood to characterize such processes, the explicit overbar notation will be dropped for clarity, and all variables will be assumed to refer to the macroscopic average state properties until further notice. 

\begin{example} 
    For a class of gravity-free, isotropic flows with no heat transfer, the fluid equations reduce from the spatio-temporal partial differential equations to time-de\-pen\-dent ordinary differential equations for spatially homogeneous fluid variables. The ratio of viscous work rate to the average-pressure work rate is then observed be proportional to the ratio of viscous relaxation time and the macroscopic time scale of the fluid motion.
    
    Assume the fluid velocity field to be isotropic and given by
    \begin{align*}
        \bm{v}(\bm{x},t)=\frac{v(t)}{L}\bm{x},
    \end{align*}
    with $v(t)$ a time-dependent function. Assume further density $\rho(\bm{x},t)=\bar{\rho}(t)$, temperature $T(\bm{x},t)=\bar{T}(t)$, pressure $p(\bm{x},t)=\bar{p}(t)$, and specific internal energy $u(\bm{x},t)=\bar{u}(t)$ all to be spatially homogeneous but time-dependent. Since $\nabla\cdot\bm{v}=3(v/L)$, substituting the ansatz to equation~\eqref{eq:density-collection} for density results in
    \begin{align*}
        \frac{1}{\rho}\frac{D\rho}{dt}=-\nabla\cdot\bm{v}\quad \implies \quad \frac{1}{\bar{\rho}}\frac{d\bar{\rho}}{dt}=-3\frac{v}{L}.
    \end{align*}
    Since $\nabla\bm{x}=(\nabla\bm{x})^T=\bm{I}$, one observes $\bm{D}(\bm{v})=\bm{0}$, and the stress tensor becomes isotropic $\bm{\sigma}=(-p+3\zeta (v/L))\bm{I}$ and spatially homogeneous. Under the ansatz, the gravity-free velocity equation~\eqref{eq:momentum-collection} then results in
    \begin{align*}
        \rho\frac{D\bm{v}}{Dt}=\nabla\cdot\bm{\sigma} \quad \implies \quad \left(\frac{1}{L}\frac{dv}{dt}+\frac{v^2}{L^2}\right)\bm{x}=\bm{0} \quad\implies\quad \frac{dv}{dt}+\frac{v^2}{L}=0.
    \end{align*}
    Using the equation for density, the equation for specific internal energy then becomes
    \begin{align*}
        \rho\frac{Du}{Dt}=\bm{\sigma}:\nabla\bm{v} \quad \implies\quad \frac{d\bar{u}}{dt}=\frac{\bar{p}}{\bar{\rho}^2}\frac{d\bar{\rho}}{dt}+9\frac{\zeta}{\bar{\rho}}\frac{v^2}{L^2}
    \end{align*}
    The original system of partial differential equation has now been reduced to a system of ordinary differential equations. Further, the ratio of viscous work rate to the average-pressure work rate now becomes
    \begin{align*}
        \frac{|\dot{\mathcal{W}}_{\mu,\zeta}|}{m|\dot{\bar{w}}_p|}=3\frac{\zeta}{\bar{p}}\frac{|v|}{L}\sim \frac{t_\text{relax}}{t_\text{macro}},
    \end{align*}
    where $t_\text{relax}\sim \zeta/\bar{p}$ is the time scale associated with viscous effect driving relaxation (originates from molecular level) and $t_\text{macro}\sim L/|v|$ is the time scale associated with the macroscopic motion of the fluid. As long as the macroscopic time scale is much longer than the micrscopic time scale driving the relaxation process, the viscous work rate can be neglected, leading to a quastistatic limiting description of the fluid.
\end{example}


\section{Specific heats}
In quasistatic experiments, the heat rate can be related to the state variables and their rates of change. Specific heat capacities are defined from particularly simple experimental processes in which one of the state variables is held constant. Define a general quasistatic process in the space of state variables, represented in terms of the pair $(\rho,T)$, $(T,p)$, or $(\rho,p)$, over a time interval $t\in[t_0,t_1]$ as the curve
\begin{align}
    t\mapsto c(t)=\left\{\begin{aligned}
        &(\rho(t),T(t))\\
        &(T(t),p(t))\\
        &(\rho(t),p(t))
    \end{aligned}\right.
\end{align}
For a quasistatic process $c(t)=(\rho,T(t))$ illustrated in Fig.~\ref{fig:isochoric-heating}, the density $\rho$ (and the specific volume $v$) is constant over time. In such case, the experimentally measurable response of the fluid to the heat it receives reduces to
\begin{align}
    \dot{q}=q_{T,\rho}\frac{dT}{dt}.
\end{align}
Under such processes, the ratio of the measured heat rate per unit mass to the measured rate of temperature change is referred to as isochoric specific heat $c_v(\rho,T)$, meaning $q_{T,\rho}=c_v$. If the pressure $p$ is kept constant instead of density as in Fig.~\ref{fig:isobaric-heating}, one considers a process path $c(t)=(p,T(t))$, and the experimentally measurable response of the fluid to the heat it receives then reduces to
\begin{align}
    \dot{q}=q_{T,p}\frac{dT}{dt}.
\end{align}
In such processes, the ratio of the measured heat rate per unit mass to the measured rate of temperature change is referred to as isobaric specific heat $c_p(\rho,T)$, meaning $q_{T,p}=c_p$. With $q_{T,\rho}$ and $q_{T,p}$ known, equations~\eqref{eq:q-Tp-relations} and~\eqref{eq:q-pT-relations} then provide the rest of the proportionality coefficients, and the general response of the fluid to the heat it receives while remaining in single phase is expressed through the relations
\begin{align}\label{eq:dot-q-expressions}
    \dot{q}=\left\{\begin{aligned}
    &c_v\frac{dT}{dt}+\frac{c_v-c_p}{\rho\beta_p}\frac{d\rho}{dt},\\
    &c_p\frac{dT}{dt}+(c_v-c_p)\frac{\kappa_T}{\beta_p}\frac{dp}{dt},\\
    &-\frac{c_p}{\rho\beta_p}\frac{d\rho}{dt}+c_v\frac{\kappa_T}{\beta_p}\frac{dp}{dt}.
    \end{aligned}\right.
\end{align}
\begin{figure}[!htb]
    \centering
    \input{\rootfolder/code/fig-heating-isochoric.tex}
    \caption{Isochoric quasistatic heating process.}
    \label{fig:isochoric-heating}
\end{figure}
\begin{figure}[!htb]
    \centering
    \input{\rootfolder/code/fig-heating-isobaric.tex}
    \caption{Isobaric quasistatic heating process.}
    \label{fig:isobaric-heating}
\end{figure}

 
The specific heats $c_v$ and $c_p$ for water are illustrated in Fig.~\ref{fig:cv-cp-H2O} and can be computed with the \verb|pyromat| and \verb|coolprop| libraries. The values of $c_v$ and $c_p$ are defined only in the single-phase regions, not in the phase-change regions. In Figure~\ref{fig:cv-cp-H2O}, one recognizes the isocontours of both $c_v$ and $c_p$ for water to be almost perfectly vertical in a region where temperature is high enough and water vapor is dilute enough (low density, high specific volume). This indicates the specific heats in the ideal gas region to approximately depend on temperature only. For gases, this is a commonly used approximation.
\begin{figure}
    \centering
    \includegraphics[width=0.8\linewidth]{\rootfolder/figs/fig-cv-H2O.pdf}
    \includegraphics[width=0.8\linewidth]{\rootfolder/figs/fig-cp-H2O.pdf}
    \caption{Specific heats for water ($c_v$ on top, $c_p$ on bottom).}
    \label{fig:cv-cp-H2O}
\end{figure}

\begin{example}
    The specific heats can be computed with both the pyromat and coolprop libraries. For water:
    
    \lstinputlisting[language=Python,firstline=1,lastline=36,firstnumber=1]{\rootfolder/code/example-cv-cp-evaluation.py}
    Both pyromat and coolprop return $c_v=2.23\,\text{kJ}/\text{kg}/\text{K}$ and $c_p=4.00\,\text{kJ}/\text{kg}/\text{K}$ at the requested state point. The values can be confirmed from Fig.~\ref{fig:cv-cp-H2O}.
\end{example}

\section{The lift-off from the temporal realm}
Experimental observations and numerical calculations confirm that the integrals of heat rate and pressure work rate over a quasistatic process depend on the process path, whereas their difference depends only on the initial and final states of the paths. For any two different quasistatic paths $c_a(t)$ and $c_b(t)$ that begin from a common state $c_a(t_{0,a})=c_b(t_{0,b})$ at initial times $t_{0,a}$ and $t_{0,b}$ and end at the same state $c_a(t_{1,a})=c_b(t_{1,b})$ at later times $t_{1,a}$ and $t_{1,b}$, one observes
\begin{align}
    \int_{c_a(t)}\dot{q}dt&\neq\int_{c_b(t)}\dot{q}dt,\\ \int_{c_a(t)}\dot{w}_pdt&\neq\int_{c_b(t)}\dot{w}_pdt,\\
    \int_{c_a(t)}(\dot{q}-\dot{w}_p)dt&=\int_{c_b(t)}(\dot{q}-\dot{w}_p)dt.
\end{align}
According to the Fundamental Theorem of line integrals, if the integral of a
quantity along a path depends only on the initial and final points of the path,
then the quantity being integrated can be interpreted as the rate of change of
some state function along the path. The observations above therefore show that
there are no state functions whose changes would correspond to the heat transfer
or the pressure work individually: the accumulated heat and pressure work depend
on the path. Their difference, however, is independent of the path. Since the
quasistatic energy balance identifies this difference with the rate of change of
specific internal energy, the specific internal energy can be represented as a
function of the fluid state variables. Because the variables are bound together by the equation of state, only two of the variables $(\rho,T,p)$ are independent. Equivalently, one may use the specific volume $v=1/\rho$ instead of density. Thus the internal energy may be represented using any convenient pair of independent state variables.
\begin{align}
    u(t)=\left\{\begin{aligned}&u(\rho(t),T(t))\\
    &u(T(t),p(t))\\
    &u(\rho(t),p(t))\end{aligned}\right.
\end{align}
\begin{example}\label{ex:path-independence}
    Define two different state points $(v_0,T_0)$ and $(v_1,T_1)$ and two different paths $c_a(t)$ and $c_b(t)$ between the points:
    \lstinputlisting[language=Python,firstline=1,lastline=20,firstnumber=1]{\rootfolder/code/example-path-independence.py}
    \begin{center}
        \includegraphics[width=0.7\textwidth]{\rootfolder/figs/example-path-independence.pdf}
    \end{center}
    Then calculate numerically the integrals of the heat rate $\dot{q}$ and mechanical pressure-related deformation work rate $\dot{w}_p$ along the two paths $c_a(t)$ and $c_b(t)$
    \lstinputlisting[language=Python,firstline=22,lastline=50,firstnumber=22]{\rootfolder/code/example-path-independence.py}
    Running the script prints out the following numbers for the heat, work, and their difference over the different paths. 
    \begin{Verbatim}
     | path a:     | path b:
--------------------------------
q    |  1.3030e+06 |  1.0295e+06 
w    |  1.5102e+06 |  1.2367e+06 
--------------------------------
q-w  | -2.0724e+05 | -2.0724e+05 
    \end{Verbatim}
    While the heat and pressure work both differ along different paths, their difference is independent of the path between the two state points.
\end{example}

Since both the specific heat rate $\dot{q}$ in~\eqref{eq:dot-q-expressions} and the specific pressure work rate $\dot{w}_p$ in~\eqref{eq:dot-w_p-expressions} are linear in the temporal rates of change of the state variables, and the specific internal energy is a function of the state variables, the entire quasistatic energy balance~\eqref{eq:quasistatic-energy-balance} becomes linearly proportional to the temporal rates of change of the state variables. This means the energy balance can be multiplied by the differential time $dt$ and lifted off from the temporal realm.

In quasistatic processes, the absolute time used to traverse the path loses its
meaning, and the energy balance becomes a differential relation on the
thermodynamic state space. This relation is the quasistatic differential form of
the First Law of Thermodynamics
\begin{align}\label{eq:first-law-differential}
    du=\delta q-\delta w_p.
\end{align}
Here $du$ is the total differential of the specific internal energy which can be expressed in terms of any two of the three state variables $(\rho,T,p)$ as
\begin{align}\label{eq:du-presentations}
    du=\left\{\begin{aligned}
        &\left(\frac{\partial u}{\partial T}\right)_{\rho}dT+\left(\frac{\partial u}{\partial \rho}\right)_{T}d\rho,\\
        &\left(\frac{\partial u}{\partial T}\right)_{p}dT+\left(\frac{\partial u}{\partial p}\right)_{T}dp,\\
        &\left(\frac{\partial u}{\partial \rho}\right)_{p}d\rho+\left(\frac{\partial u}{\partial p}\right)_{\rho}dp,
        \end{aligned}\right.
\end{align}
and the expressions $\delta q=\dot{q}dt$ and $\delta w_p=\dot{w}_pdt$ on the right, representing the heat and pressure work rates multiplied by the time differential $dt$, are differential one-forms. They too have multiple different representations depending on the choice of state variables
\begin{align}
    \label{eq:delta-q-representations}
    \delta q&=\left\{\begin{aligned}
    &c_vdT+\frac{c_v-c_p}{\rho\beta_p}d\rho,\\
    &c_pdT+(c_v-c_p)\frac{\kappa_T}{\beta_p}dp,\\
    &-\frac{c_p}{\rho\beta_p}d\rho+c_v\frac{\kappa_T}{\beta_p}dp,\end{aligned}\right.\\
    \label{eq:delta-w_p-representations}
    \delta w_p&=\left\{\begin{aligned}
    &-\frac{p}{\rho^2}d\rho,\\
    &\beta_p\frac{p}{\rho}dT-\kappa_T\frac{p}{\rho}dp.
    \end{aligned}\right.
\end{align}
Consequently, a quasistatic process $c(t)$ can be parametrized with respect to an arbitrary path variable, not just the time~$t$. Here and in what follows, the symbol $\delta$ is used instead of $d$ as a reminder that heat and pressure work
are path-dependent and are not changes of state functions.


\section{Internal energy and enthalpy}
Given the energy balance in the differential form~\eqref{eq:first-law-differential}, and the different representations for the heat and pressure work differentials $\delta q$ and $\delta w_p$ in equations~\eqref{eq:delta-q-representations} and~\eqref{eq:delta-w_p-representations}, different representations for the differential form of internal energy are identified
\begin{align}\label{eq:du-preliminary-representations}
    du=\left\{\begin{aligned}&c_vdT+\left(\frac{c_v-c_p}{\rho\beta_p}+\frac{p}{\rho^2}\right)d\rho,\\
    &\left(c_p-\beta_p\frac{p}{\rho}\right)dT+\left((c_v-c_p)\frac{\kappa_T}{\beta_p}+\kappa_T\frac{p}{\rho}\right)dp,\\
    &\left(\frac{p}{\rho^2}-\frac{c_p}{\rho\beta_p}\right)d\rho+c_v\frac{\kappa_T}{\beta_p}dp,\end{aligned}\right.
\end{align}
specifying the partial derivatives $(\partial u/\partial x)_{y}$ with respect to any two state variables $x,y\in\{\rho,T,p\}$. 

Within a single-phase region where the response coefficients are smooth, the change in internal energy between two state points $(\rho_0,T_0,p_0)$ and
$(\rho_1,T_1,p_1)$ can then be constructed by calculating the path integral
\begin{align}
    u_1-u_0=\int_c du,
\end{align}
over any chosen path $c$ connecting the two states in a single-phase region. The calculation can be performed by parametrizing the path and then computing the integral numerically, as in the example~\ref{ex:path-independence}. After choosing a reference point and tabulating the change in $u$ for many different end points, the data can be included in numerical libraries. For example, both \verb|pyromat| and \verb|coolprop| support evaluating internal energy via function calls. 

While the specific heats introduced above are smooth material properties only in single-phase regions, the specific internal energy can also be assigned to equilibrium two-phase states by calculating the mass-weighted average of the saturated liquid and saturated vapor values using the vapor quality
\begin{align}
    u=(1-x)\,u_\text{liq}+x\,u_\text{vap},\qquad \text{two-phase region}.
\end{align}

For water, the specific internal energy is illustrated in Fig.~\ref{fig:u-H2O}. As visible in the figure, the level sets of internal energy $u$ become vertical for large enough specific volumes and high enough temperatures. This indicates that the specific internal energy is approximately a function of temperature only in the ideal gas region.
\begin{figure}[!htb]
    \centering
    \includegraphics[width=0.8\linewidth]{\rootfolder/figs/fig-u-H2O.pdf}
    \caption{Specific internal energy of water.}
    \label{fig:u-H2O}
\end{figure}
\begin{example}
    The specific internal energy can be computed with both the pyromat and coolprop libraries. For water:
    \lstinputlisting[language=Python,firstline=1,lastline=35,firstnumber=1]{\rootfolder/code/example-internal-energy-evaluation.py}
    For the specified state point of $p=225\,\text{bar}$ and $T=376^\circ\text{C}$ that resides just outside the two-phase region, pyromat returns $u=2152.60\,\text{kJ}/\text{kg}$ while coolprop returns $u=2152.44\,\text{kJ}/\text{kg}$.
\end{example}


In a quasistatic heating process that happens at constant density (or specific volume), the isochoric specific heat according to~\eqref{eq:du-preliminary-representations} is directly associated with the isochoric temperature change of the specific internal energy. Since $(\partial u/\partial T)_\rho=c_v$, the isochoric change is obtained through the integral
\begin{align}
    u(T_1,\rho)-u(T_0,\rho)=\int_{T_0}^{T_1}c_v(T,\rho)dT.
\end{align}
In a quasistatic heating process that happens at constant pressure, the connection of isobaric specific heat to internal energy is less direct due to density change. Since $(\partial u/\partial T)_p=c_p-(p/\rho)\beta_p$, the isobaric change in the specific internal energy is obtained from
\begin{align}
    u(T_1,p)-u(T_0,p)&=\int_{T_0}^{T_1}\left(c_p(T,p)-\frac{p}{\rho}\beta_p\right) dT\nonumber\\
    &=\int_{T_0}^{T_1}c_p(T,p)dT+p\int_{T_0}^{T_1}\frac{1}{\rho^2}\left(\frac{\partial \rho}{\partial T}\right)_pdT\nonumber\\
    &=\int_{T_0}^{T_1}c_p(T,p)dT-\frac{p}{\rho(T_1,p)}+\frac{p}{\rho(T_0,p)},
\end{align}
illustrating the contribution from density change. This observation motivates defining a derived quantity, the so-called specific enthalpy 
\begin{align}
    h=u+\frac{p}{\rho}=u+pv,
\end{align}
for which the isobaric change from the relation above becomes an integral of the isobaric specific heat alone
\begin{align}
    h(T_1,p)-h(T_0,p)=\int_{T_0}^{T_1}c_p(T,p)dT.
\end{align}
Because both $u$ and $p/\rho$ are state functions, $h$ is also a state function. The specific enthalpy is not merely a curiosity, though. It will be an important quantity in the energy balance analysis of real flow processes over non-material volumes. Such processes include flows through turbines, pumps, heat exchangers, valves, and basically all flow-related components needed to construct power plants.

Since both the specific internal energy and the specific volume of an equilibrium
two-phase mixture can be expressed as vapor-quality-weighted averages of the
corresponding saturated liquid and saturated vapor values, so can the specific
enthalpy
\begin{align}
    h=(1-x)\,h_\text{liq}+x\,h_\text{vap}, \quad \text{two-phase region}.
\end{align}
Both \verb|pyromat| and \verb|coolprop| offer methods for calculating the specific enthalpy. For water, the specific enthalpy is illustrated in Fig.~\ref{fig:h-H2O} and, as visible in the figure, the level sets of $h$ become vertical for large enough specific volumes and high enough temperatures. This indicates that the specific enthalpy is approximately a function of temperature only in the ideal gas region, similarly as the specific internal energy. In certain applications, such in analyzing power plants and heat pumps, one often prefers a diagram with pressure and specific enthalpy on the axes for a better view of the two-phase region. Such $p$-$h$ diagram is provided in Fig.~\ref{fig:ph-T-x-H2O}.
\begin{figure}[!htb]
    \centering
    \includegraphics[width=0.8\linewidth]{\rootfolder/figs/fig-h-H2O.pdf}
    \caption{Specific enthalpy for water.}
    \label{fig:h-H2O}
\end{figure}
\begin{figure}[!htb]
    \centering
    \includegraphics[width=0.8\linewidth]{\rootfolder/figs/fig-ph-T-x-H2O.pdf}
    \caption{The pressure--enthalpy diagram for water.}
    \label{fig:ph-T-x-H2O}
\end{figure}

\section{Entropy as a state function}
If there is no heating involved in a quasistatic process, meaning $\delta q=0$, the differential form of the energy balance~\eqref{eq:first-law-differential} provides a differential equation for specific volume (or alternatively for density) in terms of temperature (or vice versa)
\begin{align}\label{eq:adiabatic-differential-equation} 
    \delta q=0
    \quad \implies \quad c_vdT+\frac{c_p-c_v}{v\beta_p}dv=0\quad \implies\quad
    \frac{dv}{dT}=-\frac{v\beta_pc_v}{c_p-c_v}.
\end{align}
In a region where the coefficients are smooth and $c_p-c_v\neq 0$, solutions to this equation define quasistatic adiabatic curves, each corresponding to a different initial state. The quasistatic adiabatic curves populate the entire single-phase domain but never cross each other. Therefore, in the neighbourhood of any initial state $(v_0,T_0)$, other states $(v,T)$ exist that are inaccessible from $(v_0,T_0)$ along adiabatic quasistatic processes alone. Reaching states off the curves requires heat to be added or removed during the process. For water such curves are illustrated in Fig.~\ref{fig:adiabatic-curves-H2O}, demonstrating the smooth alignment of the curves and the corresponding inaccessibility.
\begin{figure}[!htb]
    \centering
    \includegraphics[width=0.8\linewidth]{\rootfolder/figs/fig-adiabatic-curves-H2O.pdf}
    \caption{Selected quasistatic adiabatic curves  for water as solutions to the differential equation~\eqref{eq:adiabatic-differential-equation} resulting from the condition $\delta q=0$.}
    \label{fig:adiabatic-curves-H2O}
\end{figure}

The important observation is not merely that adiabatic curves exist, but that they
do not fill a neighbourhood of the initial state. In every neighbourhood of a state, there are nearby states that cannot be reached by following curves satisfying $\delta q=0$. This observation enables the exploitation of Carath\'eodory's theorem from differential geometry and sets the premise for finding the state function referred to as entropy. Carath\'eodory's theorem states the following. For any smooth differential one-form 
\begin{align}
    \delta \omega=\sum_{i=1}^nf_i(x_1,\ldots,x_n)dx_i,
\end{align}
defined on a domain $\Omega\subset\mathbb{R}^n$ with the functions $f_i:\Omega\to\mathbb{R}$, if for every point $P_0\in\Omega$ there exist arbitrarily close points $P\in\Omega$ inaccessible through the solution curves of $\delta \omega=0$, then the differential $\delta\omega$ admits, at least locally, an integrating factor $\gamma(x_1,\ldots,x_n)$ and a potential function $\xi(x_1,\ldots,x_n)$ such that
\begin{align}
    \delta \omega=\gamma d\xi.
\end{align}
\begin{example}
    In a two-dimensional space with coordinates $x$ and $y$, the differential 
    \begin{align*}
        \delta \omega=xdy-ydx,
    \end{align*}
    cannot be written as a total differential of any function. The family of solution curves of $\delta\omega=0$ are given by
    \begin{align*}
        \delta \omega=0 \quad \implies\quad \frac{dy}{y}=\frac{dx}{x} \quad \implies \quad y=\frac{y_0}{x_0}x.
    \end{align*}
    In the domain $\Omega=\{(x,y)\in\mathbb{R}^2|\,x\neq 0\}$, in every neighbourhood of any point $(x_0,y_0)\in\Omega$, there exist points $(x_1,y_1)\in\Omega$ inaccessible via curves satisfying $\delta\omega=0$. As an example, consider the point $(x_1,y_1)=(x_0,y_0+\epsilon)$ with $\epsilon\neq 0$. Then $(x_1,y_1)$ can be made arbitrarily close to $(x_0,y_0)$ but does not lie on the curve $y=(y_0/x_0)x$ passing through $(x_0,y_0)$. Per Carath\'eodory's theorem, an integrating factor and a potential function must exist locally in $\Omega$. In this case, they can be identified by rewriting
    \begin{align*}
        \delta \omega=xdy-ydx=x^2\left(\frac{dy}{x}-\frac{ydx}{x^2}\right)=x^2 d\left(\frac{y}{x}\right)=\gamma d\xi,
    \end{align*}
    providing $\gamma(x,y)=x^2$ and $\xi(x,y)=y/x$.
\end{example}

For water, and other fluids as well, Figure~\ref{fig:adiabatic-curves-H2O} demonstrates that the solution curves of $\delta q=0$ in the single-phase region comply with the requirements of Carath\'eodory's theorem, and that the heat differential $\delta q$ can therefore be written as
\begin{align}
    \delta q=\gamma(v,T) d\xi(v,T),
\end{align}
where $\gamma$ and $\xi$ are, for now, unknown smooth functions of the fluid's state variables. To find the functions $\gamma$ and $\xi$, one considers a system consisting of two fluids, labelled (1) and (2), with masses $m_1$ and $m_2$ in thermal equilibrium sharing a common temperature $T$ while the specific volumes $v_1$ and $v_2$ of the fluids may differ. 

The two fluids may even be entirely different substances, for example liquid water and gaseous carbon dioxide. The important requirement is that they remain in mutual thermal equilibrium and therefore share the same temperature. Such a system could be constructed by placing the fluids in separate containers, with only a thin heat-conducting wall separating the fluids. The sum of the mass-weighted heat differentials
\begin{align}
    m_1\delta q_1+m_2\delta q_2=m_1\gamma_1(v_1,T)d\xi_1(v_1,T)+m_2\gamma_2(v_2,T)d\xi_2(v_2,T),
\end{align}
then defines a new heat differential specific to the total mass of the two fluids
\begin{align}
    \delta q_{12}&=\frac{m_1}{m_1+m_2}\delta q_1+\frac{m_2}{m_1+m_2}\delta q_2
\end{align}
in a three-dimensional space with respect to the variables $T$, $v_1$, and $v_2$. 
Experimentally, one can again verify that, in every neighbourhood of a state of the combined system, there are nearby states that cannot be reached by adiabatic
quasistatic changes of $v_1$ and $v_2$ while maintaining mutual thermal equilibrium. By Carath\'eodory's theorem of inaccessible states, also the combined system then admits an integrating factor $\gamma_{12}(v_1,v_2,T)$ and a potential $\xi_{12}(v_1,v_2,T)$ such that
\begin{align}
    \delta q_{12}=\gamma_{12}d\xi_{12}.
\end{align}

Next assume $\xi(v,T)$ to be invertible and introduce a change of coordinates from $(v_1,T)$ and $(v_2,T)$ to $(\xi_1,T)$ and $(\xi_2,T)$ such that $v_1=v_1(\xi_1,T)$ and $v_2=v_2(\xi_2,T)$. The combined heat differential can now be written in the coordinates $(T,\xi_1,\xi_2)$ as
\begin{align}\label{eq:delta-q12-xi1-xi2-T}
    \delta q_{12}=\frac{m_1}{m}\gamma_1^\ast(\xi_1,T)d\xi_1+\frac{m_2}{m}\gamma_2^\ast(\xi_2,T)d\xi_2,
\end{align}
with respect to the functions $\gamma^\ast(\xi,T)=\gamma(v(\xi,T),T)$. Then apply another result from differential geometry, namely the Fröbenius theorem. In three-dimensional space the theorem states the following. If a differential form $\delta\omega =\sum_{i=1}^3f_i(x_1,x_2,x_3)dx_i$ admits an integrating factor, then the component functions $f_i$ of the differential form satisfy
\begin{align}\label{eq:frobenius-3d}
    \sum_{i,j,k=1}^{3}\varepsilon_{ijk}f_i\frac{\partial f_k}{\partial x_j}=0,
\end{align}
where $\varepsilon_{ijk}$ is the Levi-Civita tensor. 
\begin{example}
    In a three-dimensional space with coordinates $(x_1,x_2,x_3)$, consider the differential 
    \begin{align*}
        \delta \omega=\sum_{i=1}^3 f_i(x_1,x_2,x_3)dx_i.
    \end{align*} 
    Assume $\delta \omega$ to admit an integrating factor $\gamma(x_1,x_2,x_3)$ and a potential $\xi(x_1,x_2,x_3)$ such that 
    \begin{align*}
        \delta\omega=\gamma d\xi.
    \end{align*}
    Component-wide this equals
    \begin{align*}
         f_i=\gamma\frac{\partial \xi}{\partial x_i}, \quad \forall\, i\in\{1,2,3\}
    \end{align*}
    Construct then the vector $\bm{f}(\bm{x})=\sum_{i=1}^3f_i(\bm{x})\bm{\hat{x}}_i$. In vector notation, the component-wide condition now provides
    \begin{align*}
        \frac{1}{\gamma}\bm{f}=\nabla\xi \quad\implies\quad \nabla\times\left(\frac{1}{\gamma}\bm{f}\right)=\bm{0} \quad\implies\quad \bm{f}\cdot\nabla\times\bm{f}=0,
    \end{align*}
    which can be written in terms of the alternating Levi-Civita pseudo-tensor $\varepsilon_{ijk}$ as
    \begin{align*}
        \sum_{i,j,k=1}^{3}\varepsilon_{ijk}f_i\frac{\partial f_k}{\partial x_j}=0.
    \end{align*}
\end{example}
Applying the Fröbenius theorem~\eqref{eq:frobenius-3d} to the differential~\eqref{eq:delta-q12-xi1-xi2-T} with respect to coordinates $(x_1,x_2,x_3)=(\xi_1,\xi_2,T)$ and the functions $(f_1,f_2,f_3)=(\gamma_1^\ast,\gamma_2^\ast,0)$ results in
\begin{align}\label{eq:frobenius-result}
    \left(\frac{\partial \ln\gamma_1^\ast}{\partial T}\right)_{\xi_1}=\left(\frac{\partial \ln\gamma_2^\ast}{\partial T}\right)_{\xi_2}.
\end{align}
Because the variables $\xi_1$ and $\xi_2$ of systems (1) and (2) can be varied independently of each other (recall $\xi_i=\xi_i(v_i,T)$ and volumes $v_i$ can be changed independently), the equality~\eqref{eq:frobenius-result} can hold if and only if both sides depend solely on the shared variable $T$. This implies that for both systems $\gamma^\ast_i$ must be of the form
\begin{align}
    \gamma^\ast_i(\xi_i,T)=\chi(T)\Xi_i(\xi_i),
\end{align}
where $\chi(T)$ is, for now, an unknown function of temperature that is common to both systems, while the function $\Xi_i(\xi_i)$ may be different for systems (1) and (2). By the Fundamental Theorem of calculus the differential $\Xi_i(\xi_i)d\xi_i$ for any smooth $\Xi_i(\xi_i)$ can always be written as a total differential of some other function. One may therefore define a new potential $s_i^\ast(\xi_i)=s^\ast(\xi_i(v_i,T))=s_i(v_i,T)$ such that $ds_i^\ast(\xi_i)=\Xi_i(\xi_i)d\xi_i$. Therefore the heat differential for each fluid can be written
\begin{align}
    \delta q_i = \gamma_i^\ast(\xi_i,T) d\xi_i=\chi(T)\Xi_i(\xi_i)d\xi_i=\chi(T) ds_i(v_i,T).
\end{align}
This result establishes that, while $\delta q$ cannot always be written directly as a total differential, for all fluids there exists a single universal function $\chi(T)$ of temperature alone that transforms it into the total differential of a state function, $s(v,T)$, referred to as the specific entropy. Further, since $\chi(T)$ is common for all fluids, the mass-weighted sum of heat differentials demonstrates that the specific entropy of the total system consisting of two fluids is a mass-weighted sum of the specific entropies of the individual fluids
\begin{align}
    \delta q_{12}=\chi(T)d\left(\frac{m_1s_1(v_1,T)}{m_1+m_2}+\frac{m_2s_2(v_2,T)}{m_1+m_2}\right)=\chi(T)ds_{12}(v_1,v_2,T).
\end{align}

As the final step, one identifies the functional form of $\chi(T)$. Since $\delta q=\chi ds$ and $\delta q=du+pdv$, one may use $d(\chi s)=\chi ds+s\chi'(T)dT$ to write
\begin{align}
    d(u-\chi s)=-s(v,T)\chi'(T)dT-p(v,T)dv.
\end{align}
The left side is a total differential of the function $u-\chi s$ with respect to coordinates $T$ and $v$. For any smooth function, the second-order cross derivatives must be independent of the order of differentiation. Therefore the following must hold
\begin{align}
    \frac{\partial^2(u-\chi s)}{\partial T\partial v}=\frac{\partial^2(u-\chi s)}{\partial v\partial T} \quad \implies \quad
    \chi'(T)\left(\frac{\partial s}{\partial v}\right)_T=\left(\frac{\partial p}{\partial T}\right)_v=\frac{\beta_p}{\kappa_T}.
\end{align}
On the other hand, the heat differential $\delta q=\chi ds$ can be written in terms of the specific heats, providing
\begin{align}
    \delta q=c_vdT+\frac{c_p-c_v}{v\beta_p}dv \quad\implies \quad \chi\left(\frac{\partial s}{\partial v}\right)_T=\frac{c_p-c_v}{v\beta_p}
\end{align}
Dividing the two expressions, the derivative of specific entropy with respect to specific volume is eliminated, and a differential equation for the unknown function $\chi(T)$ is obtained
\begin{align}
    \frac{\chi'(T)}{\chi(T)}=\frac{v\beta_p^2}{\kappa_T(c_p-c_v)}.
\end{align}

The right side can be experimentally verified to equal $1/T$ for all substances, a result that is commonly known as the generalized Mayer relation
\begin{align}\label{eq:mayer-relation}
    \frac{vT\beta_p^2}{\kappa_T(c_p-c_v)}=1.
\end{align} 
The solution to the remaining equation $T\chi'(T)=\chi(T)$ then is $\chi(T)=a\,T$ with $a$ an arbitrary constant. Choosing to measure specific entropy $s$ in units of energy per unit mass per unit of absolute temperature sets $a=1$. For quasistatic processes, the heat differential can therefore be expressed as
\begin{align}
    \delta q=Tds.
\end{align}
\begin{example}
    The general Mayer relation can be confirmed with \verb|coolprop|:
    \lstinputlisting[language=Python,firstline=1,lastline=26,firstnumber=1]{\rootfolder/code/example-mayer-relation.py}
    Running the script provides
    \begin{Verbatim}
v = 2.00000e-03 (m^3/kg), T = 6.73150e+02 (K), mayer(v,T) = 1.00000e+00
v = 2.00000e-02 (m^3/kg), T = 7.73150e+02 (K), mayer(v,T) = 1.00000e+00
v = 1.00000e+00 (m^3/kg), T = 4.73150e+02 (K), mayer(v,T) = 1.00000e+00
v = 1.00000e+01 (m^3/kg), T = 8.73150e+02 (K), mayer(v,T) = 1.00000e+00    
    \end{Verbatim}
    confirming the relation to hold at the various selected state points.
\end{example}

The relation $\chi'(T)(\partial s/\partial v)_T=\beta_p/\kappa_T$ encountered during the derivation process of $\chi(T)=T$ provides the derivative for entropy with respect to specific volume (and density) at constant temperature. With the heat differential for quasistatic processes now expressible as $\delta q=Tds$, the experimentally obtained expressions~\eqref{eq:delta-q-representations} for $\delta q$ provide the derivative for entropy with respect to temperature at constant specific volume (and density). Derivatives with respect to the other state variable combinations are obtained by transforming the state variable differentials via $d\ln\rho=-\beta_pdT+\kappa_Tdp$, leading to the following representations for the specific entropy differential
\begin{align}
    ds=\left\{\begin{aligned}
    &\frac{c_v}{T}dT-\frac{1}{\rho^2}\frac{\beta_p}{\kappa_T}d\rho,\\
    &\frac{c_p}{T}dT-\frac{\beta_p}{\rho}dp,\\
    &-\frac{c_p}{\rho T\beta_p}d\rho+\frac{c_v}{T}\frac{\kappa_T}{\beta_p}dp.\end{aligned}\right.
\end{align}

The change in specific entropy between two state points can then be constructed by calculating a path integral of the differential $ds$ over any chosen path
\begin{align}
    s_1-s_0=\int_{c}ds.
\end{align}
After fixing a reference point and tabulating the change in $s$ for many different end points, the data can be included in numerical libraries, such as \verb|pyromat| and \verb|coolprop|. While the specific heats introduced earlier are smooth material properties only in single-phase regions, the specific entropy can also be assigned to equilibrium two-phase states by calculating the mass-weighted average of the saturated liquid and saturated vapor values using the vapor quality,
\begin{align}
    s=(1-x)\,s_\text{liq}+x\,s_\text{vap}, \qquad \text{two-phase region}.
\end{align}
For water, the specific entropy is illustrated in Figure~\ref{fig:s-H2O}. As expected, the curves of constant entropy in the singe-phase region align with the solution curves of $\delta q=0$ shown in Fig.~\ref{fig:adiabatic-curves-H2O}. In certain applications, such as in analyzing power plants and heat pumps, one often prefers a diagram with temperature and specific entropy on the axes for a better view of the two-phase region. Such a $T$-$s$ diagram for water is provided in Fig.~\ref{fig:Ts-p-x-H2O}.
\begin{figure}[!htb]
    \centering
    \includegraphics[width=0.8\linewidth]{\rootfolder/figs/fig-s-H2O.pdf}
    \caption{Specific entropy of water.}
    \label{fig:s-H2O}
\end{figure}
\begin{figure}[!htb]
    \centering
    \includegraphics[width=0.8\linewidth]{\rootfolder/figs/fig-Ts-p-x-H2O.pdf}
    \caption{The temperature--entropy diagram for water.}
    \label{fig:Ts-p-x-H2O}
\end{figure}

Finally, the observed Mayer-relation can be used to simplify the differential representations for the specific internal energy 
\begin{align}\label{eq:du-representations}
    du=\left\{\begin{aligned}&c_vdT+\left(\frac{p}{\rho^2}-\frac{\beta_p}{\kappa_T}\frac{T}{\rho^2}\right)d\rho,\\
    &\left(c_p-\beta_p\frac{p}{\rho}\right)dT+\left(\kappa_T\frac{p}{\rho}-\beta_p\frac{T}{\rho}\right)dp,\\
    &\left(\frac{p}{\rho^2}-\frac{c_p}{\rho\beta_p}\right)d\rho+c_v\frac{\kappa_T}{\beta_p}dp,\end{aligned}\right.
\end{align}
and to equivalently define the simplified differentials for specific enthalpy
\begin{align}
    dh=\left\{\begin{aligned}&\left(c_v+\frac{1}{\rho}\frac{\beta_p}{\kappa_T}\right)dT+\frac{1}{\rho^2}\left(-T\frac{\beta_p}{\kappa_T}+\frac{1}{\kappa_T}\right)d\rho,\\
    &c_pdT+\frac{1}{\rho}\left(1-T\beta_p\right)dp,\\
    &-\frac{c_p}{\rho\beta_p}d\rho+\left(\frac{1}{\rho}+c_v\frac{\kappa_T}{\beta_p}\right)dp.\end{aligned}\right.
\end{align}
In the next section, these representations are reorganized into more compact
fundamental thermodynamic relations involving $u$, $h$, $s$, $p$, $v$, and $T$.


\section{Thermodynamic potentials and computation of properties}
Given that the heat differential in a quasistatic process can be expressed in terms of a state function $\delta q=Tds$, and the energy balance written $du=\delta q-\delta w_p$ with $\delta w_p=pdv=-(p/\rho^2)d\rho$, one finds a fundamental relation between the state variables and their changes
\begin{align}\label{eq:fundamental-relation-u}
    du=Tds-pdv.
\end{align}
In terms of specific enthalpy $h=u+pv=u+p/\rho$, the fundamental relation can be written also as
\begin{align}\label{eq:fundamental-relation-h}
    dh=Tds+vdp.
\end{align}
The relations~\eqref{eq:fundamental-relation-u} and~\eqref{eq:fundamental-relation-h} are completely detached from any process details and describe how the state variables, the specific internal energy, specific enthalpy, and the specific entropy relate to each other differentially. 
Moreover, the relations enable a convenient way to calculate thermodynamic data. 

Use $d(Ts)=Tds+sdT$ and define the quantities $f=u-Ts$ and $g=h-Ts$, referred to as the specific Helmholtz free energy and the specific Gibbs free energy, respectively. Because $u$, $h$, $T$, and $s$ are state functions, $f$ and $g$ are
also state functions. The differentials of $f$ and $g$ then satisfy
\begin{align}
    \label{eq:df-helmholtz}
    df&=-sdT-pdv,\\
    \label{eq:dg-gibbs}
    dg&=-sdT+vdp.
\end{align}
The differentials reveal the natural variables of the potentials. Since $df$
contains only $dT$ and $dv$, the specific Helmholtz free energy is naturally
represented as $f(T,v)$. Since $dg$ contains only $dT$ and $dp$, the specific
Gibbs free energy is naturally represented as $g(T,p)$. Similarly the fundamental relations indicate that the natural representation for internal energy and enthalpy are the functions $u(s,v)$ and $h(s,p)$, respectively. These connections are summarized in the table~\ref{tab:natural-variable-pairings}.
\begin{table}[!htb]
    \centering
    \begin{threeparttable}
    \caption{Natural variable pairings.}
    \begin{tabular}{cccc}
        \toprule
        Quantity & Definition & Differential & Natural variables \\
        \midrule
        $u$ & --- & $du=Tds-pdv$ & $(s,v)$ \\
        $h$ & $u+pv$ & $dh=Tds+vdp$ & $(s,p)$ \\
        $f$ & $u-Ts$ & $df=-sdT-pdv$ & $(T,v)$ \\
        $g$ & $h-Ts$ & $dg=-sdT+vdp$ & $(T,p)$ \\
        \bottomrule
    \end{tabular}
    \end{threeparttable}
    \label{tab:natural-variable-pairings}
\end{table}

Since $df$ and $dg$ are total differentials, changes in $f$ and $g$ between two
state points can be computed as path integrals of the corresponding differentials~\eqref{eq:df-helmholtz} and~\eqref{eq:dg-gibbs}, provided the needed functions $s$, $p$, or $v$ are known along the path. After choosing a reference state, one may tabulate values of $f(T,v)$ or $g(T,p)$ and fit them with suitable functions. Once such a representation is available, other thermodynamic quantities can be computed from either free energy and its derivatives. Thermodynamic data libraries such as \verb|pyromat| and \verb|coolprop| are based on this idea: a fitted thermodynamic representation is used to compute other properties through differentiation and algebraic relations.
\begin{example}
    Assume the function $f(T,v)$ is known. Since
    \begin{align*}
    df&=\left(\frac{\partial f}{\partial T}\right)_vdT
    +\left(\frac{\partial f}{\partial v}\right)_Tdv
    =-sdT-pdv,
    \end{align*}
    the specific entropy and pressure as functions of $(T,v)$ are
    \begin{align*}
        s&=-\left(\frac{\partial f}{\partial T}\right)_v, &
        p&=-\left(\frac{\partial f}{\partial v}\right)_T.
    \end{align*}
    From the definition $f=u-Ts$, the specific internal energy is
    \begin{align*} 
        u=f+Ts=f-T\left(\frac{\partial f}{\partial T}\right)_v,
    \end{align*}
    and the isochoric specific heat becomes
    \begin{align*}
        c_v&=\left(\frac{\partial u}{\partial T}\right)_v
        =-T\left(\frac{\partial^2f}{\partial T^2}\right)_v.
    \end{align*}
\end{example}
\begin{example}
    Assume the function $g(T,p)$ is known. Since
    \begin{align*}
        dg=\left(\frac{\partial g}{\partial T}\right)_pdT
        +\left(\frac{\partial g}{\partial p}\right)_Tdp
        =-sdT+vdp,
    \end{align*}
    the specific entropy and specific volume as functions of $(T,p)$ are
    \begin{align*}
        s&=-\left(\frac{\partial g}{\partial T}\right)_p, &
        v&=\left(\frac{\partial g}{\partial p}\right)_T.
    \end{align*}
    From the definition $g=h-Ts$, the specific enthalpy is
    \begin{align*}
        h=g+Ts=g-T\left(\frac{\partial g}{\partial T}\right)_p,
    \end{align*}
    and the isobaric specific heat becomes
    \begin{align*}
        c_p=\left(\frac{\partial h}{\partial T}\right)_p
        =-T\left(\frac{\partial^2 g}{\partial T^2}\right)_p.
    \end{align*}
\end{example}

In terms of the variables $(T,v)$, one may use the definition $f=u-Ts$, the differential $df=-sdT-pdv$, and the differential equation of state $d\ln v=\beta_pdT-\kappa_Tdp$ to find the equation of state for pressure
\begin{align}
    p=-\left(\frac{\partial f}{\partial v}\right)_T,
\end{align}
and the expressions for the state functions specific internal energy, specific enthalpy, and the specific entropy
\begin{align}
    u&=f-T\left(\frac{\partial f}{\partial T}\right)_v,\\
    h&=f-T\left(\frac{\partial f}{\partial T}\right)_v-v\left(\frac{\partial f}{\partial v}\right)_T,\\
    s&=-\left(\frac{\partial f}{\partial T}\right)_v.
\end{align}
Further, the expressions for isobaric thermal expansion coefficient and the isothermal compressibility coefficient can be computed from
\begin{align}
    \beta_p&=-\frac{1}{v}\frac{\frac{\partial^2f}{\partial T\partial v}}{\left(\frac{\partial^2f}{\partial v^2}\right)_T}, &
    \kappa_T&=\frac{1}{v\left(\frac{\partial^2f}{\partial v^2}\right)_T},
\end{align}
and the expressions for isochoric and isobaric specific heats from
\begin{align}
    c_v&=-T\left(\frac{\partial^2f}{\partial T^2}\right)_v, &
    c_p&=-T\left(\frac{\partial^2f}{\partial T^2}\right)_v+T\frac{\left(\frac{\partial^2f}{\partial T\partial v}\right)^2}{\left(\frac{\partial^2f}{\partial v^2}\right)_T}.
\end{align}


\section{Relations for ideal gas}
The ideal-gas model is a commonly used approximation in many engineering applications. For this reason, it is useful to study it in more detail. The \verb|pyromat| library also includes ideal-gas models for various substances

The most important observation regarding ideal gases is that the specific heats are functions of temperature only. This fact is reflected for water in Fig.~\ref{fig:cv-cp-H2O}. At the same time, in both Figs.~\ref{fig:u-H2O} and~\ref{fig:h-H2O} describing the internal energy $u$ and enthalpy $h$ of water in $(T,v)$ coordinates, the level sets are practically vertical at high enough temperatures and specific volumes, i.e., in the ideal gas region. In such case both $u$ and $h$ are independent of volume, functions of temperature, and the differentials for specific internal energy and enthalpy reduce to 
\begin{align}
    du&=c_v(T)dT, & dh&=c_p(T)dT.
\end{align}
The definition of enthalpy $h=u+pv$ and the ideal gas equation of state $pv=R_sT$ then imply that the difference in the ideal gas specific heats is the specific gas constant
\begin{align}
    dh-du=(c_p-c_v)dT=d(pv)=R_sdT \quad \implies \quad c_p-c_v=R_s.
\end{align}
The same result also follows from the generalized Mayer relation~\eqref{eq:mayer-relation} by inserting
$\beta_p=1/T$, $\kappa_T=1/p$, and $v=R_sT/p$. Consequently, an ideal gas model can be described by specifying either one of the specific heats as a function of temperature. Typically, the isobaric specific heat is chosen and fitted with a NASA-polynomial
representation, for example
\begin{align}
    \frac{c_p}{R_s}=\frac{a_1}{T^2}+\frac{a_2}{T}+a_3+a_4T+a_5T^2+a_6T^3+a_7T^4.
\end{align}
The changes is specific internal energy, the specific enthalpy, and the specific entropy can then be calculated with respect to reference point from the following formulas 
\begin{align}
    u-u_0&=\int_{T_0}^{T}c_p(T')dT'-R_s(T-T_0),\\
    h-h_0&=\int_{T_0}^{T}c_p(T')dT',\\
    s-s_0&=\int_{T_0}^{T}\frac{c_p(T')}{T'}dT'-R_s\ln\frac{p}{p_0}.
\end{align}
If the specific heats are approximated as constants over the temperature interval,
these expressions reduce to
\begin{align}
    u-u_0&=c_v(T-T_0),\\
    h-h_0&=c_p(T-T_0),\\
    s-s_0&=c_p\ln\frac{T}{T_0}-R_s\ln\frac{p}{p_0}.
\end{align}
The ideal-gas equation of state alone does not require the specific heats to be
constant. When $c_p$ and $c_v$ are allowed to vary with temperature, the model is
sometimes called a thermally perfect gas. If the specific heats are additionally
approximated as constants, one obtains the calorically perfect gas model. The values of isobaric specific heats for selected ideal gases are illustrated in Figure~\ref{fig:ideal-gas-cp}.
\begin{figure}[!htb]
    \centering
    \includegraphics[width=0.8\linewidth]{\rootfolder/figs/fig-ideal-gas-cp.pdf}
    \caption{Isobaric specific heats for selected ideal gases.}
    \label{fig:ideal-gas-cp}
\end{figure}
\begin{example}
    Evaluation of the ideal gas state functions with pyromat happens identically to evaluation of multi-phase substance models. 
    \lstinputlisting[language=Python,firstline=1,lastline=29,firstnumber=1]{\rootfolder/code/example-ideal-gas-evaluation.py}
\end{example}

\section{Summary}
The thermodynamic closure missing from the fluid equations is a relation connecting specific internal energy to the thermodynamic state. This closure can be established by investigating the fluid in the quasistatic limit. In this limit, spatial pressure variations and viscous effects within the fluid become negligible, and the energy balance reduces to a relation between the change in specific internal energy, heat transfer, and pressure-driven expansion work. Measurements of quasistatic heating processes provide the specific heat capacities needed to relate heat transfer to changes in the macroscopic state variables.

Heat transfer and pressure work at the quasistatic limit are path-dependent, whereas their difference is independent of the path between two thermodynamic states. The difference can therefore be represented as the total differential of a state function, the specific internal energy,
\begin{align}
    du = \delta q - p\,dv.
\end{align}
The specific enthalpy,
\begin{align}
    h = u + pv,
\end{align}
is a derived state function that is particularly useful when describing flow processes in the next chapter.

Further investigation of quasistatic processes reveals that states exist which cannot be reached from a given state without adding or removing heat. Carath\'eodory's theorem then establishes the existence of an integrating factor for the heat differential. Experimental observations identify the corresponding temperature dependence, resulting in the definition of specific entropy through
\begin{align}
    \delta q = T\,ds.
\end{align}
Specific entropy is therefore a state function, while heat and pressure work remain path-dependent differentials.

Combining the two differential relations gives the fundamental thermodynamic relations
\begin{align}
    du &= T\,ds-p\,dv,\\
    dh &= T\,ds+v\,dp.
\end{align}
The Helmholtz and Gibbs specific free energies are then introduced as
\begin{align}
    f &= u-Ts,\\
    g &= h-Ts,
\end{align}
with corresponding differentials
\begin{align}
    df &= -s\,dT-p\,dv,\\
    dg &= -s\,dT+v\,dp.
\end{align}
Thus $f$ is naturally represented as $f(T,v)$, while $g$ is naturally represented as $g(T,p)$.

For ideal gases, the specific internal energy and specific enthalpy become functions of temperature only. The differentials reduce to
\begin{align}
    du&=c_v(T)dT, &
    dh&=c_p(T)dT,
\end{align}
and the ideal-gas equation of state gives
\begin{align}
    c_p-c_v=R_s.
\end{align}
Consequently, changes in $u$, $h$, and $s$ can be calculated from temperature-dependent specific-heat functions.

The thermodynamic properties of a substance can consequently be obtained from either a Helmholtz or Gibbs free-energy representation and its derivatives. Together, these relations provide the thermodynamic closure needed to combine the fluid equations with experimentally determined thermodynamic data, and form the basis of numerical property libraries such as \verb|pyromat| and \verb|coolprop|.

\chapter{Applications}
In this chapter the theory developed so far is used to analyze engineering applications. The partial differential equations describing fluid motion are integrated over control volumes that can represent e.g., jet engines, turbines, compressor, heat exchangers etc., and simplified balance laws for mass, momentum, energy, and entropy are discovered. The energy and entropy balances then enable finding the theoretical maximum for the efficiency at which thermal energy can be converted into mechanical work. Using the equations, the working principles of power plants, heat pumps, and internal combustion engines are explained and analyzed.

\section{Learning goals}
\begin{itemize}
    \item[\elementary] The student provides physical interpretations for the terms appearing in mass, momentum, angular momentum, energy, and entropy balances over open and closed volumes.
    \item[\elementary] The student classifies entropy transfer by heat and mass flow and entropy generation by irreversible processes.
    \item[\elementary] The student calculates heat and mechanical power rates for elementary steady-flow components from their inlet and outlet conditions.
    \item[\elementary] The student calculates the thermal efficiency of a simplified Rankine cycle and the COP of a simplified condensing heat pump from state diagrams.
    \item[\elementary] The student estimates the theoretical maximum efficiency of a heat engine and the theoretical maximum COP of a heat pump from thermodynamic state diagrams.
    \item[\intermediate] The student selects an appropriate control volume and identifies the accumulation, transport, transfer, and generation terms relevant to a given engineering process.
    \item[\intermediate] The student identifies the steady-flow and other modeling assumptions used to simplify the mass, momentum, and energy balances for an engineering component.
    \item[\intermediate] The student calculates forces and torques acting on a steady-flow device from momentum and angular momentum balances.
    \item[\intermediate] The student calculates the thermal efficiency of a power plant and the COP of a heat pump from a cycle description using \verb|pyromat| or \verb|coolprop| for thermodynamic data and component efficiencies.
    \item[\advanced] The student calculates the thermal efficiency of a reciprocating engine from experimentally sampled pressure, volume, and mass-flow-rate data.
    \item[\advanced] The student evaluates entropy generation in the components of an engineering cycle and identifies the dominant sources of irreversibility using energy and entropy balances.
\end{itemize}

\section{Control volumes in applications}
In the previous chapters, the conservation laws were applied to a material volume that moves with the fluid. Often in applications, it is more useful to choose a volume that encloses all of the fluid inside a device and inspect how, e.g., mass and energy are transferred through its boundary. Such a volume is referred to as a control volume.

Consider a time-dependent control volume $\mathcal{V}(t)$ enclosed by a boundary
\begin{align}
    \partial \mathcal{V}(t)=\mathcal{S}(t)\cup\mathcal{O}(t).
\end{align}
The boundary is formed by the piecewise union of solid parts $\mathcal{S}(t)=\bigcup_{i\in\mathcal{I}}\mathcal{S}_{i}(t)$ and open parts $\mathcal{O}(t)=\bigcup_{j\in\mathcal{J}}\mathcal{O}_j(t)$, assumed disjoint in their relative interiors. The fluid may enter or leave the control volume through the open parts of the boundary $\partial\mathcal{V}(t)$, and the boundary may move with an arbitrary velocity $\bm{u}(\bm{x},t)$. For an illustration, see Fig.~\ref{fig:control-volume}.
\begin{figure}[!htb]
    \centering
    \input{\rootfolder/code/fig-control-volume}
    \caption{A time-evolving volume $\mathcal{V}(t)$ with boundary $\partial\mathcal{V}(t)$ consisting of a stationary solid surface $\mathcal{S}_1$, a solid surface $\mathcal{S}_2(t)$ moving with velocity $\bm{u}_2(\bm{x},t)$, and stationary imaginary open surfaces $\mathcal{O}_3$ and $\mathcal{O}_4$ serving as inlet and outlet for the flow.}
    \label{fig:control-volume}
\end{figure}

A solid part of the boundary does not allow fluid to pass through it. Usually, the fluid also does not slip on such surfaces but instead follows the surface movement. Therefore, on a solid boundary of a control volume, the fluid velocity is assumed to equal the boundary velocity
\begin{align}
    \bm{v}(\bm{x},t)=\bm{u}(\bm{x},t),\quad \bm{x}\in \mathcal{S}(t).
\end{align}
An open surface is instead an imaginary surface on which the relative normal velocity $(\bm{v}-\bm{u})\cdot\hat{\bm{n}}$ need not vanish, allowing the fluid to pass through. 

Such volumes are encountered in many engineering applications. In reciprocating engines such as the Otto and Diesel engines, the control volume involves intake and exhaust valve openings that may be open or closed depending on the piston position, while the engine cylinder volume continually decreases and increases during the cycle. In a Stirling engine, the size of the control volume changes continually but no fluid passes through its boundary. In steam power plants, water passes through individual components such as turbines, steam generators, condensers, pumps and heat exchangers, which all can be considered control volumes with open surfaces. Turbines, pumps, and compressors have moving blades or pistons, making the control volume non-stationary in time while heat exchangers typically have no moving parts. Yet combined together, the components of a steam power plant enclose a volume with no open surfaces. In contrast, a control volume enclosing the fluid flowing through a gas power plant necessarily has open surfaces.

Analysis of such systems involves assuming the fluid to remain in local thermodynamic equilibrium at all times, meaning the thermodynamic relations developed in previous chapter hold point-wise within the fluid and can be combined with the local fluid equations of motion. After integration of the local fluid equations over the control volume, balance equations for mass, linear momentum, angular momentum, energy, and entropy are obtained. In general, these balances have the structure
\begin{align}
    \text{Rate of change} + \text{Outflow rate} - \text{Inflow rate} = \text{Net generation and transfer rate}
\end{align}
with the precise terms depending on the physical quantity being considered. For mass there is no generation term, while entropy is generated by temperature and velocity gradients in the flow and transferred by heat and mass flow through the boundary. Energy and momentum can be transferred through the boundary by several different mechanisms.

The balance equations can then be used to understand how the different devices function. For example, they reveal how enthalpy change in steady flow is related to mechanical work in turbines and compressors, and how the performance of such components affects the thermal efficiency of power plants. They also provide the basis for deriving the theoretical maximum efficiency of machines that convert heat into mechanical work.


\section{Mass balance}
The simplest control-volume balance is obtained by applying conservation of mass to the fluid contained in $\mathcal{V}(t)$. Since mass can enter or leave the control volume only through its open surfaces, the change of the mass must be balanced by the net mass flow through those surfaces. 

To make this statement precise, use the Leibniz rule~\eqref{eq:leibniz-rule} to write
\begin{align}
    \frac{d}{dt}\int_{\mathcal{V}(t)}\rho dV
    &=\int_{\mathcal{V}(t)}\partial_t\rho dV
    +\int_{\partial \mathcal{V}(t)}\rho\bm{u}\cdot \bm{\hat{n}}dA.
\end{align}
Using the equation~\eqref{eq:density-collection} for density, the definition~\eqref{eq:material-derivative} of the material derivative, and the divergence theorem, the right side can be written
\begin{align}
    \int_{\mathcal{V}(t)}\partial_t\rho dV+\int_{\partial \mathcal{V}(t)}\rho\bm{u}\cdot\bm{\hat{n}} dA=\int_{\partial \mathcal{V}(t)}\rho(\bm{u}-\bm{v})\cdot\bm{\hat{n}} dA.
\end{align}
Since the velocity on the solid surfaces equals the surface velocity, only the open surfaces where the surface and the flow velocity may differ, contribute to the term on the right. Therefore the rate of change of mass within $\mathcal{V}(t)$ becomes
\begin{align}
    \underbrace{\frac{d}{dt}\int_{\mathcal{V}(t)}\rho dV}_{\substack{\text{Rate of change}\\\text{of mass within}\\\text{volume}\,\mathcal{V}(t)}}+\sum_{j\in\mathcal{J}}\underbrace{\int_{\mathcal{O}_j(t)}\rho(\bm{v}-\bm{u}_{j})\cdot\bm{\hat{n}} dA}_{\substack{\text{Net mass-flow-rate out of}\,\mathcal{V}(t)\\\text{through the open surface}\,\mathcal{O}_{j}}}=0.
\end{align}
Denote the net mass-flow-rate out of the volume $\mathcal{V}(t)$ on the open surface $\mathcal{O}_j(t)$ with
\begin{align}\label{eq:mass-flow-rate-defition}
    \dot{m}_j=\int_{\mathcal{O}_j(t)}\rho(\bm{v}-\bm{u}_j)\cdot\bm{\hat{n}}dA,
\end{align}
and the mass balance over the volume $\mathcal{V}(t)$ can be expressed in a compact form
\begin{align*}
    \frac{dm}{dt}+\sum_{j\in\mathcal{J}}\dot{m}_j=0.
\end{align*}
Note the sign convention of $\dot{m}_j$ to account for the surface normal $\bm{\hat{n}}$ pointing out of the volume $\mathcal{V}(t)$. A positive $\dot{m}_j$ means mass is flowing out of the volume through the surface $\mathcal{O}_j$ while a negative value means the mass is flowing into the volume. Naturally, if the volume contains no open surfaces, then the mass within is always constant in time. 

\begin{example}
    Consider a rigid high-pressure gas tank with a small leakage. When the pressure in the tank is significantly higher than the ambient pressure surrounding the tank, the mass-flow rate $\dot{m}_\ell$ leaking out of the tank is experimentally observed to be proportional to the absolute gas pressure $p$ in the tank
    \begin{align*}
        \dot{m}_\ell=c p,
    \end{align*}
    with $c>0$ an unknown proportionality constant. Assuming the temperature of the gas in the tank to remain constant and the gas inside the tank to behave like an ideal gas, the mass in the tank can be written 
    \begin{align*}
        m(t)=\rho(t) V=\frac{p(t)V}{R_sT}.
    \end{align*}
    The balance law for mass then gives
    \begin{align*}
        \frac{dm}{dt}+\dot{m}_\ell=0 \quad \implies \quad \frac{dp}{dt}\frac{V}{R_sT}+cp=0,
    \end{align*}
    for which the solution is
    \begin{align*}
        p(t)=p(t_0)\exp\left(-\frac{cR_sT}{V}(t-t_0)\right).
    \end{align*}
    The coefficient $c$ can therefore be determined by measuring the pressure curve, and fitting it with an exponential function.
\end{example}

\section{Linear momentum balance}
Along with mass, the fluid contained in a control volume $\mathcal{V}(t)$ also carries linear momentum. The corresponding balance equation is the control-volume analog of Newton's second law. 
The linear momentum
\begin{align}
    \bm{P}(t)=\int_{\mathcal{V}(t)}\rho\bm{v}\,dV,
\end{align}
contained within $\mathcal{V}(t)$ can change because momentum is carried through the open surfaces by mass flow and because forces act on the fluid.

To obtain its rate of change, use first the density equation~\eqref{eq:density-collection} and the definition~\eqref{eq:material-derivative} of the
material derivative to write
\begin{align}
    \partial_t(\rho\bm{v})=\rho\frac{D\bm{v}}{Dt}-\nabla\cdot(\rho\bm{v}\otimes\bm{v}).
\end{align}
Then use the Leibniz rule~\eqref{eq:leibniz-rule} together with the divergence theorem to obtain
\begin{align}
    \frac{d}{dt}\int_{\mathcal{V}(t)}\rho\bm{v}dV&=\int_{\mathcal{V}(t)}\partial_t(\rho\bm{v})dV+\int_{\partial \mathcal{V}(t)}\rho\bm{v}(\bm{u}\cdot\bm{\hat{n}}) dA\nonumber\\
    &=\int_{\mathcal{V}(t)}\rho\frac{D\bm{v}}{Dt}dV+\int_{\partial \mathcal{V}(t)}\rho\bm{v}(\bm{u}-\bm{v})\cdot\bm{\hat{n}} dA.
\end{align}
Using the Cauchy momentum equation~\eqref{eq:momentum-collection} and the divergence theorem, the first integral on the right can be written as
\begin{align}
    \int_{\mathcal{V}(t)}\rho\frac{D\bm{v}}{Dt}dV&=\int_{\mathcal{V}(t)}\rho\bm{g}dV+\int_{\partial \mathcal{V}(t)}(\bm{\sigma}+p_\infty\bm{I})\cdot\bm{\hat{n}} dA.
\end{align}
Here, the constant $p_\infty$ is introduced as an ambient pressure to relate measurements to gauge pressure if desired. This addition does not change the integral, since for a surface $\partial\mathcal{V}(t)$ enclosing a volume $\mathcal{V}(t)$, one has 
\begin{align}
    \int_{\partial \mathcal{V}(t)}\bm{\hat{n}}dA=\bm{0}.
\end{align}
Using the decomposition $\bm{\sigma}=-p\bm{I}+\bm{\tau}$, the shifted stress accounting for ambient pressure can be written as
\begin{align}
    (\bm{\sigma}+p_\infty\bm{I})\cdot\bm{\hat n}
    =
    -(p-p_\infty)\bm{\hat n}
    +
    \bm{\tau}\cdot\bm{\hat n}.
\end{align}
Thus the pressure contribution to the force on a control-volume surface appears through the gauge pressure $p-p_\infty$, while the viscous contribution is unchanged.

Substituting the above expressions, the rate of change of linear momentum becomes a balance equation
\begin{align}
    &\underbrace{\frac{d}{dt}\int_{\mathcal{V}(t)}\rho\bm{v}dV}_{\substack{\text{Rate of change of}\\\text{linear momentum}\\\text{within volume}\,\mathcal{V}(t)}}+\sum_{j\in\mathcal{J}}\underbrace{\int_{\mathcal{O}_j(t)}\bm{v}\rho(\bm{v}-\bm{u}_j)\cdot\bm{\hat{n}} dA}_{\substack{\text{Rate of momentum carried out of}\\\text{volume}\,\mathcal{V}\,\text{through surface}\,\mathcal{O}_j}}\nonumber\\
    &
    =\underbrace{\int_{\mathcal{V}(t)}\rho\bm{g}dV}_{\substack{\text{Gravitational}\\\text{force on the fluid}\\\text{within}\,\mathcal{V}(t)}}
    +
    \sum_{j\in\mathcal{J}}\underbrace{\int_{\mathcal{O}_j(t)}(\bm{\sigma}+p_\infty\bm{I})\cdot\bm{\hat{n}} dA}_{\substack{\text{Stress force}\\\text{exerted on the fluid}\\\text{through open surface}\,\mathcal{O}_j}}+\sum_{i\in\mathcal{I}}\underbrace{\int_{\mathcal{S}_i(t)}(\bm{\sigma}+p_\infty\bm{I})\cdot\bm{\hat{n}}dA}_{\substack{\text{Stress force}\\\text{exerted on the fluid}\\\text{through solid surface}\,\mathcal{S}_i}}.
\end{align}
The relative normal velocity $(\bm{v}-\bm{u})\cdot\bm{\hat n}$ vanishes on solid surfaces, so the momentum-flow term contributes only on open surfaces, whereas the stress term contributes on both open and solid surfaces.

Denote the gravitational force and the net force exerted on the fluid through the solid parts of the boundary by
\begin{align}
    \bm{F}_g
    &=
    \int_{\mathcal{V}(t)}\rho\bm{g}\,dV,
    &
    \bm{F}_\mathcal{S}
    &=
    \sum_{i\in\mathcal{I}}
    \int_{\mathcal{S}_i(t)}
    (\bm{\sigma}+p_\infty\bm{I})\cdot\bm{\hat{n}}\,dA.
\end{align}
In engineering applications, the open surfaces commonly used as inlet and outlet sections are often planar or behave as approximately planar. For such cases, define the average surface direction vector $\bar{\bm n}_j$ (not necessarily a unit vector on curved surfaces), the average surface pressure $\bar p_j$, and the mass-flow-averaged velocity $\bar{\bm v}_j$ by
\begin{align}
    A_j\bar{\bm{n}}_j&=\int_{\mathcal{O}_j(t)}\bm{\hat{n}}dA,\\
    \bar{p}_j\bar{\bm{n}}_jA_j&=\int_{\mathcal{O}_j(t)}p\bm{\hat{n}}dA,\\
    \dot{m}_j\bar{\bm{v}}_j&=\int_{\mathcal{O}_j(t)}\bm{v}\rho(\bm{v}-\bm{u}_j)\cdot\bm{\hat{n}}dA.
\end{align}
Moving the open-surface momentum-flow term to the right-hand side and using the shifted stress decomposition, define the combined open-boundary contribution
\begin{align}
    \bm{F}_\mathcal{O}=-\sum_{j\in\mathcal{J}}(\dot{m}_j\bar{\bm{v}}_j+(\bar{p}_j-p_\infty)\bar{\bm{n}}_jA_j)+\sum_{j\in\mathcal{J}}\int_{\mathcal{O}_j(t)}\bm{\tau}\cdot\bm{\hat{n}}dA.
\end{align}
The linear momentum balance can now be written in a compact form
\begin{align}\label{eq:linear-momentum-balance-control-volume}
    \frac{d\bm{P}}{dt}=\bm{F}_g+\bm{F}_\mathcal{O}+\bm{F}_\mathcal{S}.
\end{align}
The force $\bm{F}_\mathcal{S}$ is the force exerted by the solid boundary on the fluid. By Newton's third law, the force exerted by the fluid on the solid boundary is therefore $-\bm{F}_\mathcal{S}$. In applications where the solid boundary is held fixed by a support structure, this opposite force is transmitted to the supports. Similarly, $\bm{F}_\mathcal{O}$ follows the sign convention of the momentum balance for the fluid, so a thrust force exerted by the fluid on a surrounding device has the opposite sign.

\begin{example}
    Consider a jet engine operating at steady state on a test stand. The engine intakes air at the inlet open surface $\mathcal{O}_1$ of area $A_1$, with mass-flow-averaged velocity $\bar{\bm{v}}_1=\bar{v}_1\bm{\hat{x}}$ and surface-averaged static pressure $\bar{p}_1$. The fuel mass-flow rate added inside the engine is small enough to be neglected compared to the air mass flow. The combustion gases exit through the nozzle outlet open surface $\mathcal{O}_2$ of area $A_2$, with mass-flow-averaged velocity $\bar{\bm{v}}_2=\bar{v}_2\bm{\hat{x}}$ and surface-averaged static pressure $\bar{p}_2$. The ambient atmospheric pressure surrounding the engine casing is a constant~$p_\infty$.

    Define a stationary control volume $\mathcal{V}$ enclosing the fluid inside the jet engine. The boundary $\partial\mathcal{V}$ is composed of the open inlet $\mathcal{O}_1$, the open outlet $\mathcal{O}_2$, and the interior solid surfaces of the engine casing $\mathcal{S}$. Because the engine is stationary on the test stand, the boundary velocity is $\bm{u}=\bm{0}$ everywhere. The flow is steady, so $d\bm{P}/dt=\bm{0}$. Neglecting the gravitational force on the fluid, the linear momentum balance reduces to
    \begin{align*}
        \bm{0}=\bm{F}_\mathcal{O}+\bm{F}_\mathcal{S}.
    \end{align*}
    Here, $\bm{F}_\mathcal{S}$ is the force exerted by the solid engine casing on the fluid, while $-\bm{F}_\mathcal{S}$ is the force exerted by the fluid on the engine casing.

    Align the coordinate system such that the flow enters from the left and exits to the right, both in the positive $\bm{\hat{x}}$ direction. The outward-pointing unit normals of the open surfaces are then
    \begin{align*}
        \bm{\hat{n}}_1=-\bm{\hat{x}},
        \qquad
        \bm{\hat{n}}_2=\bm{\hat{x}}.
    \end{align*}
    Using conservation of mass for steady flow with no accumulation, the inlet and outlet mass-flow rates satisfy
    \begin{align*}
        \dot{m}_2=-\dot{m}_1=\dot{m},
    \end{align*}
    where $\dot{m}>0$ denotes the positive through-flow rate. Neglecting viscous stresses on the open surfaces, the open-boundary contribution in the $\bm{\hat{x}}$ direction becomes
    \begin{align*}
        F_{\mathcal{O},x}&=-\sum_{j=1}^{2}\left(\dot{m}_j\bar{v}_{j,x}+(\bar{p}_j-p_\infty)\bar{n}_{j,x}A_j\right)\\
        &=\dot{m}(\bar{v}_1-\bar{v}_2)+(\bar{p}_1-p_\infty)A_1-(\bar{p}_2-p_\infty)A_2.
    \end{align*}
    The solid-boundary force on the fluid is therefore
    \begin{align*}
        F_{\mathcal{S},x}=-F_{\mathcal{O},x}=\dot{m}(\bar{v}_2-\bar{v}_1)+(\bar{p}_2-p_\infty)A_2-(\bar{p}_1-p_\infty)A_1.
    \end{align*}
    By Newton's third law, the force exerted by the fluid on the engine casing is
    \begin{align*}
        F_{\text{fluid on engine},x}=-F_{\mathcal{S},x}=F_{\mathcal{O},x}.
    \end{align*}
    If the pressure terms are small and $\bar{v}_2>\bar{v}_1$, this force is directed in the negative $\bm{\hat{x}}$ direction. Thus, with the present coordinate choice, the propulsive thrust of the engine is opposite to the direction of the exhaust flow. The force supplied by the test stand to keep the engine stationary is opposite to the force exerted by the fluid on the engine.
\end{example}

\section{Angular momentum balance}
Along with linear momentum, the fluid contained in a control volume
\(\mathcal{V}(t)\) also carries angular momentum with respect to a chosen
origin. The corresponding balance equation is the control-volume form of
the angular-momentum principle: angular momentum can change because it is
carried through the open surfaces by mass flow and because torques act on
the fluid.

The angular momentum of the fluid inside $\mathcal{V}(t)$ is
\begin{align}
    \bm{L}(t)
    =
    \int_{\mathcal{V}(t)}\bm{x}\times\rho\bm{v}\,dV,
\end{align}
where $\bm{x}$ is the position vector measured from any chosen fixed origin. To obtain the rate of change of angular momentum, first recall the identity
\begin{align}
    \partial_t(\rho\bm{v})
    =
    \rho\frac{D\bm{v}}{Dt}
    -
    \nabla\cdot(\rho\bm{v}\otimes\bm{v}).
\end{align}
Since the tensor $\rho\bm{v}\otimes\bm{v}$ is symmetric, one has $\bm{x}\times\nabla\cdot(\rho\bm{v}\otimes\bm{v})=\nabla\cdot\left(\bm{x}\times\rho\bm{v}\otimes\bm{v}\right)$, and therefore 
\begin{align}
    \bm{x}\times\partial_t(\rho\bm{v})=\bm{x}\times\rho\frac{D\bm{v}}{Dt}-\nabla\cdot\left(\bm{x}\times\rho\bm{v}\otimes\bm{v}\right).
\end{align}
Using the Leibniz rule~\eqref{eq:leibniz-rule}, the expression above, and the divergence theorem, the rate of change of angular momentum in the control volume then becomes
\begin{align}
    \frac{d}{dt}
    \int_{\mathcal{V}(t)}
    \bm{x}\times\rho\bm{v}\,dV
    &\underset{\eqref{eq:leibniz-rule}}{=}
    \int_{\mathcal{V}(t)}
    \bm{x}\times\partial_t(\rho\bm{v})\,dV
    +
    \int_{\partial\mathcal{V}(t)}
    \bm{x}\times\rho\bm{v}
    (\bm u\cdot\bm{\hat n})\,dA
    \nonumber\\
    &=
    \int_{\mathcal{V}(t)}
    \bm{x}\times\rho\frac{D\bm{v}}{Dt}\,dV
    +
    \int_{\partial\mathcal{V}(t)}
    \bm{x}\times\rho\bm{v}
    (\bm{u}-\bm{v})\cdot\bm{\hat{n}}\,dA.
\end{align}

Next note that the symmetry of the stress tensor, $\bm{\sigma}^T=\bm{\sigma}$, enables one to write $\bm{x}\times\nabla\cdot\bm{\sigma}=\nabla\cdot(\bm{x}\times\bm{\sigma})$. Crossing the Cauchy momentum equation~\eqref{eq:momentum-collection} with $\bm{x}$, integrating over the volume $\mathcal{V}(t)$, and using the divergence theorem, one finds
\begin{align}
    \int_{\mathcal{V}(t)}
    \bm{x}\times\rho\frac{D\bm{v}}{Dt}\,dV
    =
    \int_{\mathcal{V}(t)}
    \bm{x}\times\rho\bm{g}\,dV
    +
    \int_{\partial\mathcal{V}(t)}
    \bm{x}\times(\bm{\sigma}+p_\infty\bm{I})\cdot\bm{\hat{n}}\,dA.
\end{align}
Here, $p_\infty$ is again introduced as an ambient pressure so that pressure-related torques may be expressed in terms of gauge pressure. This addition does not change the integral, since
\begin{align}
    \int_{\partial\mathcal{V}(t)}\bm{x}\times\bm{\hat{n}}\,dA=\int_{\mathcal{V}(t)}\nabla\times\bm{x}\,dV=\bm{0}.
\end{align}

Substituting the above expressions gives the control-volume angular
momentum balance
\begin{align}
    &\underbrace{
    \frac{d}{dt}
    \int_{\mathcal{V}(t)}
    \bm{x}\times\rho\bm{v}\,dV
    }_{\substack{\text{Rate of change of}\\
    \text{angular momentum}\\
    \text{within volume}\,\mathcal{V}(t)}}
    +
    \sum_{j\in\mathcal{J}}
    \underbrace{
    \int_{\mathcal{O}_j(t)}
    \bm{x}\times\bm{v}\,
    \rho(\bm{v}-\bm{u}_j)\cdot\bm{\hat{n}}\,dA
    }_{\substack{\text{Rate of angular momentum}\\
    \text{carried out of volume}\,\mathcal{V}(t)\\
    \text{through open surface}\,\mathcal{O}_j}}
    \nonumber\\
    &=
    \underbrace{
    \int_{\mathcal{V}(t)}
    \bm{x}\times\rho\bm{g}\,dV
    }_{\substack{\text{Gravitational torque}\\
    \text{on the fluid within}\,\mathcal{V}(t)}}
    +
    \sum_{j\in\mathcal{J}}
    \underbrace{
    \int_{\mathcal{O}_j(t)}
    \bm{x}\times(\bm{\sigma}+p_\infty\bm{I})\cdot\bm{\hat{n}}\,dA
    }_{\substack{\text{Stress torque exerted}\\
    \text{on the fluid through}\\
    \text{open surface}\,\mathcal{O}_j}}
    +
    \sum_{i\in\mathcal{I}}
    \underbrace{
    \int_{\mathcal{S}_i(t)}
    \bm{x}\times(\bm{\sigma}+p_\infty\bm{I})\cdot\bm{\hat{n}}\,dA
    }_{\substack{\text{Stress torque exerted}\\
    \text{on the fluid through}\\
    \text{solid surface}\,\mathcal{S}_i}}.
\end{align}
The relative normal velocity $(\bm{v}-\bm{u})\cdot\bm{\hat{n}}$ vanishes on solid surfaces, so the angular-momentum-flow term contributes only on open surfaces, whereas stress torques may act on both open and solid surfaces.

Define the angular-momentum flow rate through open surface $\mathcal{O}_j(t)$ by
\begin{align}
    \dot{\bm{L}}_j
    =
    \int_{\mathcal{O}_j(t)}
    \bm{x}\times\bm{v}\,
    \rho(\bm{v}-\bm{u}_j)\cdot\bm{\hat{n}}\,dA.
\end{align}
For engineering inlet and outlet surfaces, this integral is often approximated using representative averaged values as
\begin{align}
    \dot{\bm{L}}_j
    \approx
    \bar{\bm{x}}_j\times\dot{m}_j\bar{\bm{v}}_j.
\end{align}
Denote the gravitational torque, the open-surface stress torque, and the
solid-surface stress torque by
\begin{align}
    \bm{T}_g&=\int_{\mathcal{V}(t)}
    \bm{x}\times\rho\bm{g}\,dV,
    \\
    \bm{T}_{\mathcal{O}}&=\sum_{j\in\mathcal{J}}\int_{\mathcal{O}_j(t)}\bm{x}\times(\bm{\sigma}+p_\infty\bm{I})\cdot\bm{\hat{n}}\,dA,
    \\
    \bm{T}_{\mathcal{S}}&=\sum_{i\in\mathcal{I}}\int_{\mathcal{S}_i(t)}\bm{x}\times(\bm{\sigma}+p_\infty\bm{I})\cdot\bm{\hat{n}}\,dA.
\end{align}
The angular momentum balance can then be written in the compact form
\begin{align}\label{eq:angular-momentum-balance-control-volume}
    \frac{d\bm{L}}{dt}+\sum_{j\in\mathcal{J}}\dot{\bm{L}}_j=\bm{T}_g+\bm{T}_{\mathcal{O}}+\bm{T}_{\mathcal{S}}.
\end{align}
The torque $\bm{T}_\mathcal{S}$ is the torque exerted by the solid
boundary on the fluid. By Newton's third law, the torque exerted by the
fluid on the solid boundary is therefore $-\bm{T}_\mathcal{S}$. In
rotating machinery, this opposite torque is transmitted to or from the
shaft.

\begin{example}
    Consider a centrifugal pump impeller rotating with a constant angular velocity about the pump axis directed along $\bm{\hat{z}}$. Choose the control volume $\mathcal{V}(t)$ to enclose the fluid inside the impeller. Denote by $\mathcal{S}(t)$ the solid impeller surface, and by $\mathcal{O}_1$ and $\mathcal{O}_2$ the open inlet and outlet surfaces. The open surfaces are chosen stationary, so $\bm{u}_1=\bm{u}_2=\bm{0}$. Let their outward-pointing unit normals respectively be parallel to the cylindrical coordinate radial unit normal vector
    \begin{align*}
        \bm{\hat{n}}_1=-\bm{\hat{r}}, \qquad \bm{\hat{n}}_2=\bm{\hat{r}}.
    \end{align*}
    
    In steady flow, the angular momentum inside the control volume does not change in time, so $d\bm{L}/dt=\bm{0}$. Neglecting the torque due to gravity, $\bm{T}_g\approx\bm{0}$, the angular momentum balance reduces to
    \begin{align*}
        \sum_{j\in\mathcal{J}}\dot{\bm{L}}_j=\bm{T}_\mathcal{O}+\bm{T}_\mathcal{S},
    \end{align*}
    where $\bm{T}_\mathcal{S}$ is the torque exerted by the solid impeller surface on the fluid. Let $T_\text{shaft}=\bm{\hat{z}}\cdot\bm{T}_\mathcal{S}$ denote the scalar torque exerted by the shaft through the impeller on the fluid in the positive $\bm{\hat{z}}$-direction. The torque exerted by the fluid on the impeller has the opposite sign. Projecting the angular momentum balance onto the pump axis then gives
    \begin{align*}
        T_\text{shaft}&=-\bm{\hat{z}}\cdot\bm{T}_\mathcal{O}+\sum_{j\in\mathcal{J}}\bm{\hat{z}}\cdot\dot{\bm{L}}_j
        \\
        &=-\sum_{j\in\mathcal{J}}\int_{\mathcal{O}_j}\bm{\hat{z}}\cdot\left(\bm{x}\times(\bm{\sigma}+p_\infty\bm{I})\cdot\bm{\hat{n}}\right)dA+\sum_{j\in\mathcal{J}}\int_{\mathcal{O}_j}\bm{\hat{z}}\cdot(\bm{x}\times\bm{v})\,\rho\bm{v}\cdot\bm{\hat{n}}\,dA,
    \end{align*}
    where the last expression uses $\bm{u}_j=\bm{0}$ on the open surfaces.

    Assume that viscous stresses on the open surfaces are negligible. Using the cylindrical coordinate representation for the position $\bm{x}=r\bm{\hat{r}}+z\bm{\hat{z}}$, and recalling that the open-surface normals are parallel to $\bm{\hat{r}}$, the stress contribution to the open-surface torque is observed to vanish
    \begin{align*}
        \sum_{j\in\mathcal{J}}\int_{\mathcal{O}_j}\bm{\hat{z}}\cdot\left(\bm{x}\times(\bm{\sigma}+p_\infty\bm{I})\cdot\bm{\hat{n}}\right)dA&\approx\sum_{j\in\mathcal{J}}\int_{\mathcal{O}_j}(p_\infty-p)\underbrace{\bm{\hat{z}}\cdot(\bm{x}\times\bm{\hat{n}})}_{=0}dA=0.
    \end{align*}
    Define the mass-flow-averaged tangential velocity $\bar{v}_{\theta,j}$ on the inlet and outlet surfaces $\mathcal{O}_j$ by the relation
    \begin{align*}
        \dot{m}_j\bar{v}_{\theta,j}=\int_{\mathcal{O}_j}(\bm{\hat{\theta}}\cdot\bm{v})\rho\bm{v}\cdot\bm{\hat{n}}\,dA.
    \end{align*}
    Since $\bm{\hat{z}}\cdot(\bm{x}\times\bm{v})=r(\bm{\hat{\theta}}\cdot\bm{v})$, the angular momentum flow through the open surfaces becomes
    \begin{align*}
        \sum_{j\in\mathcal{J}}\int_{\mathcal{O}_j}\bm{\hat{z}}\cdot(\bm{x}\times\bm{v})\,\rho\bm{v}\cdot\bm{\hat{n}}\,dA&=\sum_{j\in\mathcal{J}}\int_{\mathcal{O}_j}r(\bm{\hat{\theta}}\cdot\bm{v})\rho\bm{v}\cdot\bm{\hat{n}}\,dA
        \\
        &=\sum_{j\in\mathcal{J}}r_j\dot{m}_j\bar{v}_{\theta,j}.
    \end{align*}
    Finally, use the steady-state mass balance $\dot{m}_2=-\dot{m}_1=\dot{m}$, where $\dot{m}>0$ is the positive through-flow rate of mass. The required shaft torque exerted on the fluid can then be expressed
    \begin{align*}
        T_\text{shaft}=\dot{m}\left(r_2\bar{v}_{\theta,2}-r_1\bar{v}_{\theta,1}\right),
    \end{align*}
    and one observes the shaft torque to be approximately determined by the net flux of angular momentum through the impeller.
\end{example}


\section{Energy balance}
In engineering applications, analysis of energy flows over a control volume $\mathcal{V}(t)$ often plays a key role. For carrying out such analysis, first define the total internal energy $U(t)$, total kinetic energy $E_k(t)$, and the total gravitational potential energy $E_p(t)$ the fluid has in the control volume 
\begin{align}
    U&=\int_{\mathcal{V}(t)}\rho udV, & E_k&=\int_{\mathcal{V}(t)}\rho\frac{1}{2}|\bm{v}|^2dV, & E_p&=\int_{\mathcal{V}(t)}\rho GdV.
\end{align}
Here $G(\bm{x})$ is the gravitational potential field, from which the gravitational acceleration $\bm{g}=-\nabla G$ is defined. Close to Earth's surface the potential can be approximated as $G=g(z-z_0)$, with $z_0$ some arbitrary reference height.

To calculate the time rates of change of these quantities, consider the following identities. The definition~\eqref{eq:material-derivative} of the material derivative together with the density equation~\eqref{eq:density-collection} provide
\begin{align}\label{eq:partial-internal-energy-density}
    \partial_t(\rho u)&=\rho\frac{Du}{Dt}-\nabla\cdot(\rho u\bm{v}).
\end{align}
Using the material derivative and some algebraic vector manipulations provides
\begin{align}\label{eq:partial-kinetic-energy-density}
    \partial_t\left(\rho\frac{1}{2}|\bm{v}|^2\right)&=\rho\frac{D}{Dt}\left(\frac{1}{2}|\bm{v}|^2\right)-\nabla\cdot\left(\rho\frac{1}{2}|\bm{v}|^2\bm{v}\right)\nonumber\\
    &=\bm{v}\cdot\rho\frac{D\bm{v}}{Dt}-\nabla\cdot\left(\rho\frac{1}{2}|\bm{v}|^2\bm{v}\right).
\end{align}
Using the density equation, and the fact that the gravitational potential is independent of time, provides
\begin{align}\label{eq:partial-potential-energy-density}
    \partial_t(\rho G)&=\rho\bm{v}\cdot\nabla G-\nabla\cdot(\rho G\bm{v})\nonumber\\
    &=-\bm{v}\cdot\rho\bm{g}-\nabla\cdot(\rho G\bm{v}),
\end{align}
where also the definition of the gravitational acceleration was used. Applying the Leibniz rule~\eqref{eq:leibniz-rule} to calculate the time derivative of the total internal energy, and combining the result with the identity~\eqref{eq:partial-internal-energy-density}, and divergence theorem, and the local energy equation~\eqref{eq:internal-energy-collection}, one obtains
\begin{align}
    \frac{d}{dt}\int_{\mathcal{V}(t)}\rho udV&\underset{\eqref{eq:leibniz-rule}}{=}\int_{\mathcal{V}(t)}\partial_t(\rho u)dV+\int_{\partial\mathcal{V}(t)}\rho u\bm{u}\cdot\bm{\hat{n}}dA\nonumber\\
    &\underset{\eqref{eq:partial-internal-energy-density}}{=}\int_{\mathcal{V}(t)}\rho\frac{Du}{Dt}dV-\int_{\partial\mathcal{V}(t)}\rho u(\bm{v}-\bm{u})\cdot\bm{\hat{n}}dA\nonumber\\
    &\underset{\eqref{eq:internal-energy-collection}}{=}\int_{\mathcal{V}(t)}(q_V+\bm{\sigma}:\nabla\bm{v})dV-\int_{\partial\mathcal{V}(t)}(\bm{q}+\rho u(\bm{v}-\bm{u}))\cdot\bm{\hat{n}}dA.
\end{align}
For the time derivative of kinetic energy, one similarly uses the Leibniz rule together with the divergence theorem, the identity~\eqref{eq:partial-kinetic-energy-density}, and Cauchy's momentum equation~\eqref{eq:momentum-collection} to obtain
\begin{align}
    \frac{d}{dt}\int_{\mathcal{V}(t)}\rho\frac{1}{2}|\bm{v}|^2dV&\underset{\eqref{eq:leibniz-rule}}{=}\int_{\mathcal{V}(t)}\partial_t\left(\frac{1}{2}\rho|\bm{v}|^2\right)dV+\int_{\partial\mathcal{V}(t)}\frac{1}{2}\rho|\bm{v}|^2\bm{u}\cdot\bm{\hat{n}}dA\nonumber\\
    &\underset{\eqref{eq:partial-kinetic-energy-density}}{=}\int_{\mathcal{V}(t)}\bm{v}\cdot\rho\frac{D\bm{v}}{Dt}dV-\int_{\partial\mathcal{V}(t)}\frac{1}{2}\rho|\bm{v}|^2(\bm{v}-\bm{u})\cdot\bm{\hat{n}}dA\nonumber\\
    &\underset{\eqref{eq:momentum-collection}}{=}\int_{\mathcal{V}(t)}(\bm{v}\cdot\rho\bm{g}+\bm{v}\cdot\nabla\cdot\bm{\sigma})dV-\int_{\partial\mathcal{V}(t)}\frac{1}{2}\rho|\bm{v}|^2(\bm{v}-\bm{u})\cdot\bm{\hat{n}}dA.
\end{align}
For the time derivative of gravitational potential energy one uses instead the identity~\eqref{eq:partial-potential-energy-density} to obtain
\begin{align}
    \frac{d}{dt}\int_{\mathcal{V}(t)}\rho GdV&\underset{\eqref{eq:leibniz-rule}}{=}\int_{\mathcal{V}(t)}\partial_t(\rho G)dV+\int_{\partial\mathcal{V}(t)}\rho G\bm{u}\cdot\bm{\hat{n}}dA\nonumber\\
    &\underset{\eqref{eq:partial-potential-energy-density}}{=}-\int_{\mathcal{V}(t)}\bm{v}\cdot\rho\bm{g}dV-\int_{\partial\mathcal{V}(t)}\rho G(\bm{v}-\bm{u})\cdot\bm{\hat{n}}dA.
\end{align}

In all of the time derivatives above, the relative velocity $\bm{v}-\bm{u}$ appears, multiplied with internal energy density, kinetic energy density, or gravitational energy density. Further, the gravitational acceleration appears in two of the time derivatives, with opposite sign. This suggests one should investigate the sum of the expressions, the time derivative of the total energy. Before doing so, it is useful to note that $\bm{\sigma}:\nabla\bm{v}+\bm{v}\cdot\nabla\cdot\bm{\sigma}=\nabla\cdot(\bm{v}\cdot\bm{\sigma})$. Combining the time derivatives, grouping the open and solid surfaces separately, and separating the pressure and viscous parts of stress on open surfaces lets one then write 
\begin{align}
    &\underbrace{\frac{d}{dt}\int_{\mathcal{V}(t)}\rho (u+\frac{1}{2}|\bm{v}|^2+G)dV}_{\substack{\text{Rate of change of internal, kinetic, and}\\\text{potential energy within volume}\,\mathcal{V}(t)}}+\sum_{j\in\mathcal{J}}\underbrace{\int_{\mathcal{O}_j(t)} (h+\frac{1}{2}|\bm{v}|^2+G)\rho(\bm{v}-\bm{u}_j)\cdot\bm{\hat{n}}dA}_{\substack{\text{Rate at which enthalpy, kinetic energy, and potential}\\\text{energy are carried out of volume}\,\mathcal{V}(t)\,\text{by mass}\\\text{through the boundary}\,\mathcal{O}_j(t)}}\nonumber
    \\
    &=\underbrace{\int_{\mathcal{V}(t)}q_VdV}_{\substack{\text{Net rate of heat}\\\text{absorbed}\\\text{volumetrically}}}+\sum_{i\in\mathcal{I}}\underbrace{\int_{\mathcal{S}_i(t)}(-\bm{q})\cdot\bm{\hat{n}}dA}_{\substack{\text{Net rate of heat}\\\text{received through the}\\\text{solid boundary}\,\mathcal{S}_i(t)}}+\sum_{j\in\mathcal{J}}\underbrace{\int_{\mathcal{O}_j(t)}(-\bm{q})\cdot\bm{\hat{n}}dA}_{\substack{\text{Net rate of heat}\\\text{received through the}\\\text{open boundary}\,\mathcal{O}_j(t)}}\nonumber
    \\    &+\sum_{i\in\mathcal{I}}\underbrace{\int_{\mathcal{S}_i(t)}\bm{v}\cdot\bm{\sigma}\cdot\bm{\hat{n}}\,dA}_{\substack{\text{Rate of mechanical work}\\\text{done by solid surface}\,\mathcal{S}_i(t)\\\text{on the fluid}}}
    +\sum_{j\in\mathcal{J}}\underbrace{\int_{\mathcal{O}_j(t)}\left(-\bm{u}_jp+\bm{v}\cdot\bm{\tau}\right)\cdot\bm{\hat{n}}dA}_{\substack{\text{Rate of work done by pressure due to}\\\text{boundary motion and by viscous stresses}\\\text{on the open surface}\,\mathcal{O}_j(t)}}
\end{align}
The appearance of enthalpy in the open-surface term comes from combining the advective internal-energy flow with the pressure part of the stress work, using $h=u+p/\rho$.

If one then denotes the volumetric and boundary heat rates into the fluid with
\begin{align}
    \dot{Q}_{\mathcal{V}}&=\int_{\mathcal{V}(t)}q_VdV, & \dot{Q}_{\mathcal{S}}&=\sum_{i\in\mathcal{I}}\int_{\mathcal{S}_i(t)}(-\bm{q})\cdot\bm{\hat{n}}dA, & \dot{Q}_{\mathcal{O}}&=\sum_{j\in\mathcal{J}}\int_{\mathcal{O}_j(t)}(-\bm{q})\cdot\bm{\hat{n}}dA,
\end{align}
the rates by which the fluid performs mechanical work on the solid and open boundaries with 
\begin{align}
    \dot{W}_{\mathcal{S}}&=-\sum_{i\in\mathcal{I}}\int_{\mathcal{S}_i(t)}\bm{v}\cdot\bm{\sigma}\cdot\bm{\hat{n}}dA, & \dot{W}_{\mathcal{O}}&=-\sum_{j\in\mathcal{J}}\int_{\mathcal{O}_j(t)}(-\bm{u}_jp+\bm{v}\cdot\bm{\tau})\cdot\bm{\hat{n}}dA,
\end{align}
and the mass-flow-averaged specific enthalpy $\bar{h}_j$, specific kinetic energy $\bar{e}_{k,j}$, and specific potential energy $\bar{e}_{p,j}$ on the open surface $\mathcal{O}_j(t)$ with
\begin{align}
    \dot{m}_j\bar{h}_j&=\int_{\mathcal{O}_j(t)}h\rho(\bm{v}-\bm{u}_j)\cdot\bm{\hat{n}}dA\\
    \dot{m}_j\bar{e}_{k,j}&=\int_{\mathcal{O}_j(t)}\frac{1}{2}|\bm{v}|^2\rho(\bm{v}-\bm{u}_j)\cdot\bm{\hat{n}}dA\\
    \dot{m}_j\bar{e}_{p,j}&=\int_{\mathcal{O}_j(t)}G\rho(\bm{v}-\bm{u}_j)\cdot\bm{\hat{n}}dA,
\end{align}
the energy balance can be written in a compact form
\begin{align}\label{eq:energy-balance-control-volume}
    \frac{d}{dt}(U+E_k+E_p)+\sum_{j\in\mathcal{J}}\dot{m}_j(\bar{h}_j+\bar{e}_{k,j}+\bar{e}_{p,j})=\dot{Q}_\mathcal{V}+\dot{Q}_\mathcal{S}+\dot{Q}_\mathcal{O}-\dot{W}_\mathcal{S}-\dot{W}_\mathcal{O}.
\end{align}
In many practical engineering applications involving flow through a device, the overall flow is steady over time $(d/dt)(U,E_k,E_p)\approx (0,0,0)$, the open inlet and outlet surfaces do not move ($\bm{u}_j=\bm{0}$), and the work rate of viscous stress on the open surfaces is small ($\dot{W}_\mathcal{O}\approx 0$).

\section{Entropy balance}
The final quantity to study in relation to a control volume $\mathcal{V}(t)$ is entropy. As $s(\bm{x},t)$ refers to specicif entropy per unit mass, now a field depending on spatial location and time, the total entropy contained within the control volume is defined
\begin{align}
    S(t)=\int_{\mathcal{V}(t)}\rho s dV.
\end{align}
To calculate its time rate of change, one first proceed similarly as in case of energy. The definition of material derivative together the density equation provide
\begin{align}
    \partial_t(\rho s)=\rho\frac{Ds}{Dt}-\nabla\cdot(\rho s\bm{v})
\end{align}
Then Leibniz rule~\eqref{eq:leibniz-rule} together with the divergence theorem provide
\begin{align}
    \frac{d}{dt}\int_{\mathcal{V}(t)}\rho s dV&=\int_{\mathcal{V}(t)}\partial_t(\rho s)dV+\int_{\partial\mathcal{V}(t)}\rho s\bm{u}\cdot\bm{\hat{n}}dA\nonumber\\
    &=\int_{\mathcal{V}(t)}\rho\frac{Ds}{Dt}dV-\int_{\partial\mathcal{V}(t)}\rho s(\bm{v}-\bm{u})\cdot\bm{\hat{n}}dA.
\end{align}

From here, the path differs from the analysis of energy balance. No equation of motion for the specific entropy is known yet. It is obtained, though, using the fundamental relation~\eqref{eq:fundamental-relation-u} and the existing local fluid equations. From the fundamental relation $du=Tds-pdv$ one first obtains
\begin{align}
    ds=\frac{1}{T}du+\frac{p}{T}dv=\frac{1}{T}du-\frac{p}{T\rho^2}d\rho.
\end{align}
The partial derivatives of specific entropy with respect to specific internal energy and density can then be identified as
\begin{align}
    \left(\frac{\partial s}{\partial u}\right)_\rho&=\frac{1}{T}, & \left(\frac{\partial s}{\partial \rho}\right)_u&=-\frac{p}{T\rho^2}.
\end{align}
With the help of these partial derivatives, use of the chain rule of differentiation then provides an expressions for the material derivative of entropy in tems of known expressions for internal energy and density
\begin{align}
    \rho\frac{Ds}{Dt}&=\left(\frac{\partial s}{\partial u}\right)_\rho \rho\frac{Du}{Dt}+\left(\frac{\partial s}{\partial \rho}\right)_u \rho\frac{D\rho}{Dt}=\frac{\rho}{T}\frac{Du}{Dt}-\frac{p}{T}\frac{1}{\rho}\frac{D\rho}{Dt}.
\end{align}
Substituting the density equation~\eqref{eq:density-collection} and the energy equation~\eqref{eq:internal-energy-collection}, and performing some algebraic manipulations, one then obtains the explicit expression
\begin{align}
    \rho\frac{Ds}{Dt}&=\nabla\cdot\left(\frac{-\bm{q}}{T}\right)+\frac{q_V}{T}+\frac{\lambda|\nabla T|^2}{T^2}+\frac{2\mu\bm{D}:\bm{D}+\zeta(\nabla\cdot\bm{v})^2}{T}.
\end{align}
An important observation is that the last to terms are always non-negative.

Combining the results above, treating open and solid surfaces separately, the following balance equation is then found
\begin{align}
    &\underbrace{\frac{d}{dt}\int_{\mathcal{V}(t)}\rho s dV}_{\substack{\text{Rate of entropy}\\\text{change within}\\\text{volume}\,\mathcal{V}(t)}}+\sum_{j\in\mathcal{J}}\underbrace{\int_{\mathcal{O}_j(t)}s\rho(\bm{v}-\bm{u}_j)\cdot\bm{\hat{n}}dA}_{\substack{\text{Rate at which entropy is carried}\\\text{out of volume}\,\mathcal{V}\,\text{by mass}\\\text{through open surface}\,\mathcal{O}_j}}\nonumber\\
    &=\sum_{i\in\mathcal{I}}\underbrace{\int_{\mathcal{S}_i(t)}\frac{-\bm{q}}{T}\cdot\bm{\hat{n}}dA}_{\substack{\text{Reversible entropy}\\\text{tansfer rate from}\\\text{heat flow across}\\\text{solid surface}\,\mathcal{S}_i}}+\sum_{j\in\mathcal{J}}\underbrace{\int_{\mathcal{O}_j(t)}\frac{-\bm{q}}{T}\cdot\bm{\hat{n}}dA}_{\substack{\text{Reversible entropy}\\\text{transfer rate from}\\\text{heat flow across}\\\text{open surface}\,\mathcal{O}_j}}+\underbrace{\int_{\mathcal{V}(t)}\frac{q_V}{T}dV}_{\substack{\text{Reversible entropy}\\\text{transfer rate from}\\\text{volumetric heating}}}\nonumber\\
    &\quad+\underbrace{\int_{\mathcal{V}(t)}\left(\frac{\lambda|\nabla T|^2}{T^2}+\frac{2\mu\bm{D}:\bm{D}+\zeta(\nabla\cdot\bm{v})^2}{T}\right)dV}_{\substack{\text{Irreversible entropy generation rate from}\\\text{non-homogeneous temperature and velocity,}\\\text{always non-negative}\,\geq 0}}
\end{align}
The expression contains several terms which deserve individual notation. Denote the heat-related reversible entropy transfer rates and the irreversible entropy generation rate with
\begin{align}
    \dot{S}_{\text{heat},\mathcal{S}}&=\sum_{i\in\mathcal{I}}\int_{\mathcal{S}_i(t)}\frac{-\bm{q}}{T}\cdot\bm{\hat{n}}dA\\
    \dot{S}_{\text{heat},\mathcal{O}}&=\sum_{j\in\mathcal{J}}\int_{\mathcal{O}_j(t)}\frac{-\bm{q}}{T}\cdot\bm{\hat{n}}dA\\
    \dot{S}_{\text{heat},\mathcal{V}}&=\int_{\mathcal{V}(t)}\frac{q_V}{T}dV\\
    \dot{S}_{\text{irr},\mathcal{V}}&=\int_{\mathcal{V}(t)}\left(\frac{\lambda|\nabla T|^2}{T^2}+\frac{2\mu\bm{D}:\bm{D}+\zeta(\nabla\cdot\bm{v})^2}{T}\right)dV\geq 0
\end{align}
and define the mass-flow-averaged specific entropy $\bar{s}_j$ on the open surface $\mathcal{O}_j(t)$ with 
\begin{align}
    \dot{m}_j\bar{s}_j=\int_{\mathcal{O}_j(t)}s\rho(\bm{v}-\bm{u}_j)\cdot\bm{\hat{n}}dA.
\end{align}
The entropy balance over the volume $\mathcal{V}(t)$ can then be written in a compact form
\begin{align}
    \frac{dS}{dt}+\sum_{j\in\mathcal{J}}\dot{m}_j\bar{s}_j=\dot{S}_{\text{heat},\mathcal{S}}+\dot{S}_{\text{heat},\mathcal{O}}+\dot{S}_{\text{heat},\mathcal{V}}+\dot{S}_{\text{irr},\mathcal{V}}
\end{align}


\section{Theoretical maximum efficiency for closed flows}
Consider a flow process in a closed volume so that the solid boundary fully encloses the volume $\mathcal{S}=\partial\mathcal{V}$. Since the volume is closed, the are no open boundaries through which mass flow could happen. The energy and entropy balances then reduce to 
\begin{align}
    \frac{d}{dt}(U+E_k+E_p)&=\dot{Q}_{\partial \mathcal{V}}+\dot{Q}_{\mathcal{V}}-\dot{W}_{\partial \mathcal{V}}\\
    \frac{dS}{dt}&=\dot{S}_{\text{heat},\partial \mathcal{V}}+\dot{S}_{\text{heat},\mathcal{V}}+\dot{S}_{\text{irr},\mathcal{V}}
\end{align}
To investigate theoretical efficiency of reciprocating engines, heat pumps, and power plants, one considers cyclic systems.

Assume the flow is cyclic, with one cycle happening during $t\in[0,\Delta t]$. Assume the total internal, kinetic, and potential energies as well as the total entropy to be cyclic, meaning $(U,E_k,E_p,S)(t+\Delta t)=(U,E_k,E_p,S)(t)$ for any time $t$. Introduce the positive and negative parts of the local surface heat flux and volumetric heat source through the functions
\begin{align}
    [-\bm{q}\cdot\bm{\hat{n}}]^{+}&=\max(0,-\bm{q}\cdot\bm{\hat{n}}), & 
    [-\bm{q}\cdot\bm{\hat{n}}]^{-}&=\min(0,-\bm{q}\cdot\bm{\hat{n}}),\\
    [q_V]^{+}&=\max(0,q_V), &
    [q_V]^{-}&=\min(0,q_V),
\end{align}
These distinguish local heat addition and heat rejection contributions via the splits
\begin{align}
    -\bm{q}\cdot\bm{\hat{n}}&=[-\bm{q}\cdot\bm{\hat{n}}]^{+}+[-\bm{q}\cdot\bm{\hat{n}}]^{-}, & &\bm{x}\in\partial\mathcal{V},\quad t\in[0,\Delta t]\\
    q_V&=[q_V]^{+}+[q_V]^{-}, & &\bm{x}\in\mathcal{V},\quad t\in[0,\Delta t].
\end{align}
The total heat injected and rejected over the cycle can now be expressed as the space-time integrals
\begin{align}
    Q_\text{in}&=\int_0^{\Delta t}\left(\int_{\partial\mathcal{V}}[-\bm{q}\cdot\bm{\hat{n}}]^{+}dA+\int_{\mathcal{V}}[q_V]^{+}dV\right)dt,\\
    Q_\text{out}&=\int_0^{\Delta t}\left(\int_{\partial\mathcal{V}}[-\bm{q}\cdot\bm{\hat{n}}]^{-}dA+\int_{\mathcal{V}}[q_V]^{-}dV\right)dt,
\end{align}
and the total heat-related entropy transferred into and out of the volume as the integrals
\begin{align}
    S_\text{heat,in}&=\int_0^{\Delta t}\left[\int_{\partial\mathcal{V}}\frac{[-\bm{q}\cdot\bm{\hat{n}}]^{+}}{T}dA+\int_{\mathcal{V}}\frac{[q_V]^{+}}{T}dV\right]dt\\
    S_\text{heat,out}&=\int_0^{\Delta t}\left[\int_{\partial\mathcal{V}}\frac{[-\bm{q}\cdot\bm{\hat{n}}]^{-}}{T}dA+\int_{\mathcal{V}}\frac{[q_V]^{-}}{T}dV\right]dt
\end{align}
Finally, define the average temperatures at which heat injection and rejection happen as ratios of heat-related entrop transfer to heat transfer
\begin{align}
    \frac{1}{\bar{T}_\text{in}}&=\frac{S_\text{heat,in}}{Q_\text{in}}, & 
    \frac{1}{\bar{T}_\text{out}}&=\frac{S_\text{heat,out}}{Q_\text{out}}
\end{align}
Both averaged temperatures are positive since the corresponding heat and entropy transfers have identical signs.

Integrate then the energy balance over the cycle to get
\begin{align}
    \int_0^{\Delta t}\frac{d}{dt}(U+E_k+E_p)dt&=\int_0^{\Delta t}\left(\dot{Q}_{\partial \mathcal{V}}+\dot{Q}_{\mathcal{V}}-\dot{W}_{\partial \mathcal{V}}\right)dt\nonumber\\
    \implies 0&=Q_\text{in}+Q_\text{out}-\int_0^{\Delta t}\dot{W}_{\partial\mathcal{V}}dt
\end{align}
and perform similar integration of the entropy balance to get
\begin{align}
    \int_{0}^{\Delta t}\frac{dS}{dt}dt&=\int_0^{\Delta t}\left(\dot{S}_{\text{heat},\partial \mathcal{V}}+\dot{S}_{\text{heat},\mathcal{V}}+\dot{S}_{\text{irr},\mathcal{V}}\right)dt\nonumber\\    
    \implies 0&=S_\text{heat,in}+S_\text{heat,out}+\int_0^{\Delta t}\dot{S}_{\text{irr},\mathcal{V}}dt
\end{align}
The efficiency $\eta$ at which the cycle converts heat to mechanical work is defined as the ratio of net mechanical work produced to the heat injected. With the integrated energy and entropy balances and the average temperatures, this ratio can now be expressed as
\begin{align}
    \eta=\frac{\int_0^{\Delta t}\dot{W}_\mathcal{V}dt}{Q_\text{in}}=1-\frac{\bar{T}_\text{out}}{\bar{T}_\text{in}}\left(1+\frac{\int_0^{\Delta t}\dot{S}_{\text{irr},\mathcal{V}}dt}{S_\text{heat,in}}\right).
\end{align}

One may observe that the efficiency is maximized by minimizing both the irreversible entropy generation and the ratio $\bar{T}_\text{out}/\bar{T}_\text{in}$ of average heat rejection and injection temperatures. For maximum efficiency, heat should therefore be injected at as high average temperature as possible and rejected at as low average temperature as possible. The theoretical maximum is achieved in the idealized quasistatic limit, with irreversible entropy generation vanishing ($\int_0^{\Delta t}\dot{S}_{\text{irr},\mathcal{V}}dt\ll S_\text{heat,in}$), and all heat injection and rejection occurring at fixed temperatures that correspond to the maximum and minimum temperatures of the flow
\begin{align}
    T_\text{min}=\min(\{T(\bm{x},t)|\bm{x}\in\partial\mathcal{V}\cup\mathcal{V},t\in[0,\Delta t]\}),\\
    T_\text{max}=\max(\{T(\bm{x},t)|\bm{x}\in\partial\mathcal{V} \cup\mathcal{V},t\in[0,\Delta t]\}).
\end{align}
This limit leads to expression
\begin{align}
    \eta_\text{Carnot}=1-\frac{T_\text{min}}{T_\text{max}},
\end{align}
known as the Carnot efficiency. Since the heat addition and rejection occur at different fixed temperatures, the remaining parts of the cycle must change the fluid temperature without heat transfer. With zero entropy generation, these parts are isentropic, meaning entropy remains constant. Such idealized flow process, containing two isothermal parts and two isentropic parts, is referred to as Carnot process. 

\begin{example}
    For water, one possible depiction of a Carnot process would utilize phase change to maintain temperature constant during heat addition and removal. In a $T$-$s$ diagram, a theoretical Carnot process could therefore look like the one below drawn with the dashed orange curve:    
    \begin{center}
    \begin{tikzpicture}
    \node[above right,inner sep = 0] (image) at (0,0) {
        \includegraphics[width=0.85\linewidth]{\rootfolder/figs/fig-Ts-p-x-H2O.pdf}
    };
    \begin{scope}[relative to node=image]
        % \foreach \ratio in {-1,-0.9,...,1}{
        %    \draw[help lines,purple] (-1,\ratio) -- (1,\ratio);
        %    \draw[help lines,purple] (\ratio,-1) -- (\ratio,1);
        % }
        % \draw (-1,0) -- (1,0);
        % \draw (0,-1) -- (0,1);
        
        \coordinate (A) at (0.055,-0.23);
        \coordinate (B) at (0.055,-0.7);
        \coordinate (C) at (-0.36,-0.7);
        \coordinate (D) at (-0.36,-0.23);

        % Process point annotations
        \fill (A) circle (2pt) node[above right=0.15cm and 0.15cm, circle,fill=white,draw=black,inner sep=1pt] {1};
        \fill (B) circle (2pt) node[below right=0.15cm and 0.15cm, circle,fill=white,draw=black,inner sep=1pt] {2};
        \fill (C) circle (2pt) node[below left=0.15cm and 0.15cm, circle,fill=white,draw=black,inner sep=1pt] {3};
        \fill (D) circle (2pt) node[above left=0.15cm and 0.15cm, circle,fill=white,draw=black,inner sep=1pt] {4};

        % Carnot rectanle
        \draw[orange, ultra thick,dashed] (B)--(C)|-(A)-| cycle;

    \end{scope}
    \end{tikzpicture}
    \end{center}
\end{example}

\section{Power plant and heat pump components}
A useful application for the balance equations of energy and entropy is found in system-level process analysis of power plants and heat pumps. In such analyses, various engineering components are encountered that can be treated as steady-flow open systems. At the system level, the flows through the components are often approximated as black boxes, assuming single values for the thermodynamic state variables and property functions in the inlet and outlet surfaces. The details of the exact dynamics inside the boxes is then ignored. Sometimes one-dimensional description for the flow through components is used, though, and the average themodynamic state along the flow through a component is parametrized with respect to, e.g., pressure or the heat rate exchanged. Effectively, the one-dimensional approximation lets one describe the components as paths in thermodynamic state space.

\paragraph{Single-phase heat exchangers:}
Two streams of fluid at different temperatures, separated by a solid material surface, exchange heat. This principle is utilized to construct heat exchangers needed in various energy engineering applications. In single-phase heat exchangers, both streams retain their phase throughout the heat exchange process, while their temperatures change from their corresponding inlet values to the outlet values. In such processes, both streams are treated as individual steady-flow open systems that interact via the solid heat exchange structure. The stream gaining heat receives it from the solid structure while the stream losing heat yields it to the solid structure. Theoretically, the solid stucture separating the streams is chosen as thin as possible with as large a surface area as possible to enable efficient heat transfer. In practice, the possibly large pressure difference between the streams limits structural and material choices. For streams with small pressure differences, thin flat metal plates can be used as the separating structure but for streams with signicant pressure difference the most common option is an array of circular pipes. Assuming the different streams to be labeled with $A$ and $B$, and the inlet and outlet denoted by indexes 1 and 2, Fig.~\ref{fig:heat-exchanger} illustrates a potential counter-flow single-phase heat exchanger.
\begin{figure}[!htb]
    \centering
    \input{\rootfolder/code/fig-counter-flow-heat-exchanger.tex}
    \caption{A counter-flow single-phase heat exchanger with colors indicating stream temperatures.}
    \label{fig:heat-exchanger}
\end{figure}

As heat exchangers effectively contain no moving parts, there is no mechanical work involved on any of the solid surfaces, meaning $\dot{W}_\mathcal{S}=0$. Typically, also the mechanical work rate associated with the open inlet and outlet surfaces is neglected, meaning $\dot{W}_\mathcal{O}\approx 0$. The change in the enthalpy between the stream inlet and outlet is much larger than the change in the stream kinetic or potential energy, meaning $\bar{h}_2-\bar{h}_1\gg\bar{e}_{p,2}-\bar{e}_{p,1}$ and $\bar{h}_2-\bar{h}_1\gg\bar{e}_{k,2}-\bar{e}_{k,1}$. The energy balance then reduces to the conditions
\begin{align}
    \dot{m}_A(h_{A,2}-h_{A,1})=\dot{Q}_{A}\\
    \dot{m}_B(h_{B,2}-h_{B,1})=\dot{Q}_{B}.    
\end{align}
While the solid heat exchanger structure is exposed to the surrounding environment, the structure is optimized so that the heat loss rate to the environment can be neglected compared to the heat exchange rate between the streams. Therefore one may assume 
\begin{align}
    \dot{Q}_A+\dot{Q}_B&\approx0.
\end{align}
Typically, also the potential pressure drop due to viscous effects on the solid walls is neglegted, making the process at system level isobaric for both streams. This is justified if the heat exchanger is used in a larger system containing also turbines, pumps, or compressors that typically results in much larger pressure changes in the total system.

Heat exchangers are often illustrated drawing $T$-$\dot{Q}$ diagrams. In such diagrams, one draws the temperatures of both streams as if they were perfectly one-dimensional and parametrized with respect to the accumulated heat transfer rate along the streams that is taken positive. Consequently, the change in temperature of the stream gaining heat is aligned with the direction of the $\dot{Q}$-axis while the stream losing heat is read from the opposite direction. Such diagrams are often supplemented with additional axes for stream enthalpy changes along the $\dot{Q}$-axis. As heat only conducts from higher temperature to lower temperature through the solid heat exchanger structure, for the system-level $T$-$\dot{Q}$ description of the heat exchanger to be physically possible, the temperature curves should then never cross. 

\begin{example}\label{example:single-phase-heat-exchanger}
    In a CO$_2$ heat pump, an internal heat exchanger operating in a single phase can be is used to improve effiency between different phases of the cycle. The stream $A$ of CO$_2$ that is cooled enters the heat exchanger at $T_{A,1}=40^\circ\text{C}$ at a pressure of $p_{A,1}=120\,\text{bar}$. At the outlet, the stream is at temperature of $T_{A,2}=20^\circ\text{C}$ with no meaningful pressure change. Simultaneously the stream $B$ of CO$_2$ that is to be heated enters the heat exchanger as saturated vapor at a temperature of $T_{B,1}=15^\circ\text{C}$. If the mass flow rates of the streams are $\dot{m}_A=0.2\,\text{kg}/\text{s}$ and $\dot{m}_B=0.3\,\text{kg}/\text{s}$, at what temperature does the stream $B$ exit the heat exchanger?

    The problem can be solved using the energy balance equations for the two streams. This is conveniently done by writing a short script. First load the data and collect known information    \lstinputlisting[language=Python,firstline=1,lastline=25,firstnumber=1]{\rootfolder/code/example-single-phase-heat-exchanger.py}
    Then use $\dot{Q}_A+\dot{Q}_B=0$ and the energy balance equations for both streams to get
    \begin{align*}
        h_{B,2}=h_{B,1}+\frac{\dot{m}_A}{\dot{m}_B}(h_{A,1}-h_{A,2})
    \end{align*}
    and use the exit temperature and enthalpy to calculate temperature
    \lstinputlisting[language=Python,firstline=27,lastline=35,firstnumber=27]{\rootfolder/code/example-single-phase-heat-exchanger.py}
    Running the script then provide $T_{B,2}=32.50^\circ\text{C}$.

    The associated $T$-$\dot{Q}$ diagram for the heat exchanger is illustrated below, with the stream B gaining heat. Note that the temperature curves are not perfectly straight, demonstrating the isobaric temperature dependency of enthalpy is not necessarily exactly linear and the isobaric heat capacity a constant.
    \begin{center}
    \begin{tikzpicture}
    \node[above right,inner sep = 0] (image) at (0,0) {
        \includegraphics[width=0.85\linewidth]{\rootfolder/figs/fig-single-phase-heat-exchanger-example.pdf}
    };
    \begin{scope}[relative to node=image]
        % \foreach \ratio in {-1,-0.9,...,1}{
        %    \draw[help lines,purple] (-1,\ratio) -- (1,\ratio);
        %    \draw[help lines,purple] (\ratio,-1) -- (\ratio,1);
        % }
        % \draw (-1,0) -- (1,0);
        % \draw (0,-1) -- (0,1);
        
        \coordinate (BB) at (0.945,0.55);
        \coordinate (BA) at (-0.8,-0.245);
        \coordinate (AA) at (0.945,0.9);
        \coordinate (AB) at (-0.8,-0.015);
        
        \fill (BB) circle (2pt) node[right=0.15cm, circle,fill=white,draw=black,inner sep=1pt] {B,2};
        \fill (BA) circle (2pt) node[above left=0.18cm and 0.25cm, circle,fill=white,draw=black,inner sep=1pt] {B,1};
        \fill (AA) circle (2pt) node[right=0.15cm, circle,fill=white,draw=black,inner sep=1pt] {A,1};
        \fill (AB) circle (2pt) node[above left=0.18cm and 0.25cm, circle,fill=white,draw=black,inner sep=1pt] {A,2};
       
    \end{scope}
    \end{tikzpicture}
    \end{center}
\end{example}

\paragraph{Condensing and evaporating heat exchangers:} Single-phase heat exchangers are not the only option for transferring heat between two fluid streams. Under appropriate conditions, one or both of the streams can undergo also a phase transition during heat transfer. In such cases one specifically refers to condensing or evaporating heat exchangers that are fundamental in steam power plants and heat pumps.

In phase-change heat exchangers, the stream that receives heat absorbs it at a nearly constant pressure, just as in the single-phase heat exchangers, but undergoes potentially up to three distinct thermal regimes. First, the stream can be heated sensibly as a subcooled liquid until it reaches its saturation temperature. Second, the stream undergoes a phase change at constant temperature, absorbing latent heat until the saturated liquid is fully converted into saturated vapor. Finally, the stream may be further heated into the superheated vapor region. Simultaneously, the stream that yields the heat may either cool down in a single phase, or similarly undergo potentially up to three thremal regimes from superheated vapor to subcooled liquid but in reversed order. 

Geometrically, the evaporation or condensation process is gravity driven. In condensers used in steam power plants, cold water is circulated in the coolant piping to enable heat removal from the steam, causing droplets and ultimately a thin film to form on the cooling pipe surfaces. Once the liquid film grows large enough, the excess liquid that surface tension can no longer hold together drops down to the bottom of the condenser, forming a pool of saturated liquid. Conversely, in pool-type evaporating steam generators the heat providing piping, with a hot heat yielding stream passing through it, is immersed in liquid water. As the water boils, the saturated vapor bubbles rise up up to the steam dome while the denser liquid stays at the bottom of the pool. Both condensing and evaporating heat exchangers are illustrated in Fig.~\ref{fig:condenser-evaporator} with the stream A undergoing the phase transition and stream B undergoing temperature change in single phase.
\begin{figure}[!htb]
    \centering
    \input{\rootfolder/code/fig-condenser-evaporator.tex}
    \caption{Condensing (left) and evaporating (right) heat exchangers, with colors indicating stream temperatures.}
    \label{fig:condenser-evaporator}
\end{figure}

\begin{example}
    In a steam power plant condenser, the liquid coolant water is designed to enter the heat exchanger at $T_{B,1}=20^\circ\text{C}$ and $p_{B}=1.1\,\text{bar}$ and exit at $T_{B,2}=35^\circ\text{C}$ while the low-pressure steam is expected to enter at $T_{A,1}=40^\circ\text{C}$ and $p_A=0.05\,\text{bar}$ and condense into saturated liquid at the outlet $x_{A,2}=0$. The mass flow rate of the steam is expected to be $\dot{m}_A=300\,\text{kg}/\text{s}$. 
    
    To estimate the necessary coolant mass flow rate, one calculates the inlet and outlet enthalpy values and uses the energy balance equations. This can be implemented conveniently in a short script
    \lstinputlisting[language=Python,firstline=1,lastline=28,firstnumber=1]{\rootfolder/code/example-condenser.py}
    Running the script provides $\dot{m}_{B}=11\,656\,\text{kg}/\text{s}$. At first glance, everything looks fine. The incoming steam is at a high temperature than the designed coolant exit temperature, and the saturation temperature of $32.9^\circ\text{C}$ at $0.05\,\text{bar}$ is higher than the designed coolant inlet temperature of $20^\circ\text{C}$. 
    However, if one extends the script by calculating the enthalpy values with respect to cumulative heat transfer rate and the respective stream temperatures 
    \lstinputlisting[language=Python,firstline=30,lastline=35,firstnumber=30]{\rootfolder/code/example-condenser.py}
    the resulting $T$-$\dot{Q}$ diagram shows the stream temperature curves to cross. The intended design of the condenser will not work. Instead, the coolant outlet temperature should be designed low enough for the curves not to cross, requiring a larger coolant mass flow rate.
    \begin{center}
    \begin{tikzpicture}
    \node[above right,inner sep = 0] (image) at (0,0) {
        \includegraphics[width=0.85\linewidth]{\rootfolder/figs/example-condenser.pdf}
    };
    \begin{scope}[relative to node=image]
        % \foreach \ratio in {-1,-0.9,...,1}{
        %    \draw[help lines,purple] (-1,\ratio) -- (1,\ratio);
        %    \draw[help lines,purple] (\ratio,-1) -- (\ratio,1);
        % }
        % \draw (-1,0) -- (1,0);
        % \draw (0,-1) -- (0,1);
        
        \coordinate (BB) at (0.95,0.615);
        \coordinate (BA) at (-0.765,-0.245);
        \coordinate (AA) at (0.95,0.9);
        \coordinate (AB) at (-0.765,0.49);
        
        \fill (BB) circle (2pt) node[right=0.15cm, circle,fill=white,draw=black,inner sep=1pt] {B,2};
        \fill (BA) circle (2pt) node[above left=0.18cm and 0.25cm, circle,fill=white,draw=black,inner sep=1pt] {B,1};
        \fill (AA) circle (2pt) node[right=0.15cm, circle,fill=white,draw=black,inner sep=1pt] {A,1};
        \fill (AB) circle (2pt) node[above left=0.18cm and 0.25cm, circle,fill=white,draw=black,inner sep=1pt] {A,2};
       
    \end{scope}
    \end{tikzpicture}
    \end{center}
\end{example}

\paragraph{Steam and gas turbines:} As black boxes, steam and gas turbines convert the enthalpy of the entering high pressure and high temperature working fluid into mechanical work causing the working fluid to exit at lower pressure and temperature due to reduction in the fluid enthalpy. A schematic illustration of a turbine is provided in Fig.~\ref{fig:turbine}.
\begin{figure}[!htb]
    \centering
    \input{\rootfolder/code/fig-turbine.tex}
    \caption{A turbine with colors indicating working fluid temperature (and pressure).}
    \label{fig:turbine}
\end{figure}

Within the turbine, the working fluid passes through subsequent turbine stages consisting of alternating stationary stator blades and rotating rotor blades. Passing through the stator, the varying area of the flow channel is designed to convert part of the working fluid enthalpy into additional kinetic energy, a process during which both the fluid pressure and temperature reduce. Passing through the rotor blades, the fluid direction changes, causing a pressure difference on the rotor blade surfaces, forces on the blades, a change of the flow angular momentum, and torque on the turbine shaft. Simultaneously, the absolute kinetic energy of the flow is reduced and transferred to the turbine shaft as rotational energy. Depending on the stage type (impulse or reaction), the rotor blade channel can also be shaped to result in a partial conversion of the fluid enthalpy to kinetic energy (reaction type) and cause further reduction in temperature and pressure. Passing through multiple stages, the fluid exits the turbine at low pressure and reduced temperature. In steam turbines, the exiting steam is typically close to the saturated vapor line, with a vapor quality often slightly below unity. In gas turbines the exit temperature is often quite high. 


As the kinetic energy is increased in the stator phases and subsequently reduced in the rotor phases, the total change in the working fluid kinetic energy between the turbine inlet and outlet is often negligible compared to the total enthalpy change $\bar{e}_{k,2}-\bar{e}_{k,1}\ll \bar{h}_2-\bar{h}_1$. In steam and gas turbines, also the change in the potential energy is negligible compared to the enthalpy change $\bar{e}_{p,2}-\bar{e}_{p,1}\ll \bar{h}_2-\bar{h}_1$. While there always exists a nonzero heat transfer rate to the turbine surroundings due to the surroundings being a lower temparature than the steam or gas passing through the turbine, this rate is limited by the heat conductivity of the solid turbine casing and typically much smaller than the mechanical work rate $\dot{Q}_{\mathcal{S}}\ll\dot{W}_\mathcal{S}$. The same applies applies to the heat transfer rate accross the open inlet and outlet surfaces $\dot{Q}_\mathcal{O}\ll\dot{W}_\mathcal{S}$. As the inlet and outlet surfaces are static, there is no pressure work rate on them associated with the surface motion, and the viscous work rate on the open surfaces from the fluid motion can be assumed small compared to the mechanical work rate on the blades $\dot{W}_\mathcal{O}\ll\dot{W}_\mathcal{S}$. As a result, the steady-flow energy balance across the turbine reduces to 
\begin{align}
    \dot{m}(h_2-h_1)=-\dot{W}
\end{align}
In heat exchangers, the heat related entropy transfer rate dominates the irreversible entropy generation rate. In turbines, however, no significant heat transfer rate exists to the environment and the entropy change between the inlet and outlet primarily results from the non-homogeneous local working fluid temperature and velocity profiles. 
\begin{align}
    \dot{m}(s_2-s_1)=\dot{S}_\text{irr}
\end{align}
At system-level, an ideal turbine would therefore be isentropic. Turbine performance with respect to an ideal turbine is measured with a so-called isentropic turbine efficiency that is defined by comparing the work rate of a real turbine ($\dot{W}$) to that of a hypothetical isentropic turbine ($\dot{W}_s$) with the formula reducing to the process inlet and outlet enthalpy values
\begin{align}
    \eta_\text{s,turbine}=\frac{\dot{W}}{\dot{W}_\text{s}}=\frac{h_1-h_2}{h_1-h_{2s}}.
\end{align}
The comparison is performed assuming the isentropic turbine inlet conditions to equal the inlet conditions of the real turbine, while at the exit only the pressure is assumed equal. 
\begin{example}
    Steam enters a back-pressure turbine at $T_1=550^\circ\text{C}$ and $p_1=200\,\text{bar}$ and exits at $p_2=1\,\text{bar}$ and $x_2=0.9$. The mass-flow rate of steam through the turbine is $\dot{m}=500\,\text{kg}/\text{s}$.

    Per the provided values for temperature, pressure, and vapor quality, the inlet and outlet specific enthalpies can be calculated with a short script:
    \lstinputlisting[language=Python,firstline=1,lastline=20,firstnumber=1]{\rootfolder/code/example-turbine.py}
    The resulting work rate is $\dot{W}=102.11\,\text{MW}$. The isentropic efficiency can be obtained by first calculating the entropy at the turbine inlet and using it together with the outlet pressure
    \lstinputlisting[language=Python,firstline=22,lastline=26,firstnumber=22]{\rootfolder/code/example-turbine.py}
    The resulting isentropic efficiency of the turbine is $\eta_\text{s,turb}=0.86$. The process curve corresponding to a hypothetical isentropic turbine is illustrated below together with a dashed straight line connecting true turbine inlet and outlet states. Note that the dashed line does not represent a true process state path.
    \begin{center}
    \begin{tikzpicture}
    \node[above right,inner sep = 0] (image) at (0,0) {
        \includegraphics[width=0.85\linewidth]{\rootfolder/figs/example-turbine.pdf}
    };
    \begin{scope}[relative to node=image]
        % \foreach \ratio in {-1,-0.9,...,1}{
        %    \draw[help lines,purple] (-1,\ratio) -- (1,\ratio);
        %    \draw[help lines,purple] (\ratio,-1) -- (\ratio,1);
        % }
        % \draw (-1,0) -- (1,0);
        % \draw (0,-1) -- (0,1);
        
        \coordinate (A) at (0.66,0.623);
        \coordinate (B) at (0.267,-0.122);
        \coordinate (Bs) at (0.203,-0.122);
        
        \fill (A) circle (2pt) node[right=0.18cm, circle,fill=white,draw=black,inner sep=1pt] {1};
        \fill (B) circle (2pt) node[below right=0.18cm and 0.18cm, circle,fill=white,draw=black,inner sep=1pt] {2};
        \fill (Bs) circle (2pt) node[below left=0.18cm and 0.18cm, circle,fill=white,draw=black,inner sep=1pt] {2s};
       
    \end{scope}
    \end{tikzpicture}
    \end{center}
\end{example}

\paragraph{Compressors:} In system-level analysis, compressors perform the inverse function of turbines. They consume mechanical work to increase the enthalpy of the working fluid, resulting in higher pressure and temperature at the exit. A schematic illustration of an axial compressor is provided in Fig.~\ref{fig:compressor}.
\begin{figure}[!htb]
    \centering
    \input{\rootfolder/code/fig-compressor.tex}
    \caption{An axial compressor with colors indicating working fluid temperature and pressure.}
    \label{fig:compressor}
\end{figure}

Within an axial compressor, the fluid passes through alternating stages of rotating rotor blades and stationary stator blades, just as in turbines. The rotor blades impart kinetic energy to the fluid by increasing its angular momentum, effectively increasing the relative velocity of the flow. Subsequently, the stator blades act as diffusers, decelerating the fluid to convert the kinetic energy gained in the rotor into enthalpy. This process of enthalpy addition repeats through multiple rotor--stator stages, progressively increasing the fluid enthalpy which results in temperature and pressure rise. 

Similar to turbines, the net changes in kinetic and potential energy between the inlet and outlet are generally negligible compared to the net enthalpy change, and the heat transfer rate to the surroundings is minimized by the casing. The steady-flow energy balance across the compressor simplifies to the same expression as for a turbine
\begin{align}
    \dot{m}(h_2-h_1) = -\dot{W}.
\end{align}
Irreversible entropy generation happens as in turbines, driven by viscous effects and flow non-homogeneities. As a theoretical limiting case, also an ideal compressor is isentropic. Compressor performance can therefore be quantified similarly as in turbines by comparing the ideal isentropic work rate to the real compressor work rate to compress the working fluid to the same pressure at the exit
\begin{align}
    \eta_\text{s,compressor} = \frac{\dot{W}_\text{s}}{\dot{W}} = \frac{h_{2s}-h_1}{h_2-h_1}.
\end{align}
Contrary to the turbine, the numerator and denominator have switched locations. An isentropic compressor requires less power input than a real compressor, while an isentropic turbine producess more power than a real turbine.

\begin{example}
    In a CO$_2$ heat pump, the working fluid arrives to the compressor as saturated vapor at a temperature of $T_1=-5^\circ\text{C}$. At the compressor exit the working fluid is at a pressure of $p_2=100\,\text{bar}$ and at a temperature of $T_2=95^\circ\text{C}$. The mass flow rate of CO$_2$ in the heat pump is $\dot{m}=0.5\,\text{kg}/\text{s}$ .

    Per the provided values for temperature, pressure, and vapor quality, the inlet and outlet specific enthalpies can be calculated with a short script:
    \lstinputlisting[language=Python,firstline=1,lastline=20,firstnumber=1]{\rootfolder/code/example-compressor.py}
    The resulting work rate is $\dot{W}=-31.5\,\text{kW}$, with the negative value illustrating that the compressor requires work input. The isentropic efficiency can be obtained by first calculating the entropy at the compressor inlet and using it together with the outlet pressure
    \lstinputlisting[language=Python,firstline=22,lastline=26,firstnumber=22]{\rootfolder/code/example-compressor.py}
    The resulting isentropic efficiency of the compressor is $\eta_\text{s,comp}=0.79$. The process curve corresponding to a hypothetical isentropic compressor is illustrated below together with a dashed straight line connecting the true compressor inlet and outlet states. Note that the dashed line does not represent a true process state path.
    \begin{center}
    \begin{tikzpicture}
    \node[above right,inner sep = 0] (image) at (0,0) {
        \includegraphics[width=0.85\linewidth]{\rootfolder/figs/example-compressor.pdf}
    };
    \begin{scope}[relative to node=image]
        % \foreach \ratio in {-1,-0.9,...,1}{
        %    \draw[help lines,purple] (-1,\ratio) -- (1,\ratio);
        %    \draw[help lines,purple] (\ratio,-1) -- (\ratio,1);
        % }
        % \draw (-1,0) -- (1,0);
        % \draw (0,-1) -- (0,1);
        
        \coordinate (A) at (0.38,-0.61);
        \coordinate (B) at (0.657,0.53);
        \coordinate (Bs) at (0.6,0.53);
        
        \fill (A) circle (2pt) node[left=0.18cm, circle,fill=white,draw=black,inner sep=1pt] {1};
        \fill (B) circle (2pt) node[above right=0.18cm and 0.18cm, circle,fill=white,draw=black,inner sep=1pt] {2};
        \fill (Bs) circle (2pt) node[above left=0.18cm and 0.18cm, circle,fill=white,draw=black,inner sep=1pt] {2s};
       
    \end{scope}
    \end{tikzpicture}
    \end{center}    
\end{example}


\paragraph{Pumps:} Many engineering applications from water treatment plants to steam power plants have pumps. Where compressors are used to pressurize gaseous fluids, the role of pumps is to move and pressurize liquid fluids. Like compressors, they require work input to accomplish the pressurization. An illustration of centrifugal pump is provided in Fig.~\ref{fig:pump}.
\begin{figure}[!htb]
    \centering
    \input{\rootfolder/code/fig-pump.tex}
    \caption{A centrifugal pump with inlet at the impeller eye.}
    \label{fig:pump}
\end{figure}

Inside a centrifugal pump, liquid water enters impeller blade structure through the impeller eye. The rotation of the blades causes the flow to generate angular momentum and additional kinetic energy before passing through the impeller outlet and into the shaped inner side of the pump casing functioning as a diffuser. In the diffuser, the additional kinetic energy of the liquid is efficiently converted into a pressure increase due to the near-incompressible nature of the liquid phase. The temperature of the liquid barely increases. By impeller blade design, the ratio of added work between kinetic energy and static pressure related energy may be altered for different applications.

At the power plant system-engineering level, modeling of pumping systems ignores the changes in the kinetic and potential energies between the inlet and the outlet, and assumes the heat rate through the pump casing to be ignorable in comparison to the mechanical work rate needed for the pump-system operation. The energy balance therefore again reduces to 
\begin{align}
    \dot{m}(h_2-h_1)=-\dot{W}
\end{align}
In other applications, such as in waste water treatment plants, pump systems are often primarily used to move liquids and not necessarily to provide significant static pressure differences between potentially long piping network inlets and outlets. In such cases, hydrostatic pressure differences originating from potential energy changes and potential volumetric flow rate changes along the piping network must be accounted for in the energy balance. 

As the flow inside the pump and the associated piping network contains flow gradients due to viscous effects, some irreversible entropy generation always takes place and typically dominates over heat-associated entropy transfer rate. Similarly as in a compressor, an ideal pump system would be isentropic, and the efficiency of the pump at system level modeling can be quantified with an isentropic efficiency by comparing the isentropic work rate to the real needed work rate
\begin{align}
    \eta_\text{s,pump}=\frac{\dot{W}_{s}}{\dot{W}}=\frac{h_{2s}-h_1}{h_2-h_1}
\end{align}
As the liquid inside a pump system is often nearly incompressible, meaning $\rho\approx\text{constant}$, a practical estimate for the isentropic enthalpy change in a pump system can be obtained with the expected pressure change alone
\begin{align}
    h_{2s}-h_1\approx \rho^{-1}(p_2-p_1).
\end{align}
An estimate of isentropic pump work rate, in systesm dominated by static pressure rise, is then obtained by multiplying the volumetric flow rate $\dot{V}=\rho^{-1}\dot{m}$ with the pressure change
\begin{align}
    W_\text{s}\approx \dot{m}\rho^{-1}(p_1-p_2)=\dot{V}(p_1-p_2)
\end{align}
\begin{example}
    A centrifugal pump system used in a steam power plant pressurizes saturated water from $p_1=0.1\,\text{bar}$ to $p_2=200\,\text{bar}$ with the water temperature rising $T_2-T_1=2^\circ\text{C}$ in the process.

    \lstinputlisting[language=Python,firstline=1,lastline=31,firstnumber=1]{\rootfolder/code/example-pump.py}
    The resulting isentropic efficiencies are $\eta_\text{s,pump}=0.7848$ and $\eta_\text{s,pump,approx}=0.7876$, meaning the approximation is very good. The process curve corresponding to a hypothetical isentropic pump is nearly vertical in the associated $p$-$h$ diagram, credited to the nearly constant temperature and density. The straight dashed line connecting the true pump inlet and outlet state points, not describing the true process state path, is difficult to separate from the isentropic pump curve.
    \begin{center}
    \begin{tikzpicture}
    \node[above right,inner sep = 0] (image) at (0,0) {
        \includegraphics[width=0.85\linewidth]{\rootfolder/figs/example-pump.pdf}
    };
    \begin{scope}[relative to node=image]
        % \foreach \ratio in {-1,-0.9,...,1}{
        %    \draw[help lines,purple] (-1,\ratio) -- (1,\ratio);
        %    \draw[help lines,purple] (\ratio,-1) -- (\ratio,1);
        % }
        % \draw (-1,0) -- (1,0);
        % \draw (0,-1) -- (0,1);
        
        \coordinate (A) at (-0.677,-0.453);
        \coordinate (B) at (-0.667,0.627);
        \coordinate (Bs) at (-0.672,0.627);
        
        \fill (A) circle (2pt) node[below=0.18cm, circle,fill=white,draw=black,inner sep=1pt] {1};
        \fill (B) circle (2pt) node[above right=0.18cm and 0.18cm, circle,fill=white,draw=black,inner sep=1pt] {2};
        \fill (Bs) circle (2pt) node[above left=0.18cm and 0.18cm, circle,fill=white,draw=black,inner sep=1pt] {2s};
       
    \end{scope}
    \end{tikzpicture}
    \end{center}
\end{example}

\paragraph{Expansion valves:} Inside refrigeration systems and heat pumps, expansion valves perform the functional inverse of a compressor or pump by reducing the pressure of the working fluid. At the system level, they are treated as steady-flow open systems that receive a high-pressure fluid and output it at a lower pressure. The geometric mechanism relies on forcing the fluid through a severe restriction, such as a narrow orifice, a capillary tube, or a needle valve regulated by a spring mechanism, as illustrated in Fig.~\ref{fig:expansion-valve}. As the fluid passes through this narrow opening, viscous effects result in pressure drop without requiring moving mechanical parts that extract work from the system. Consequently, a portion of high-pressure subcooled or saturated liquid passing through the valve flashes immediately into a low-temperature, low--pressure two--phase liquid--vapor mixture.
\begin{figure}[!htb]
    \centering
    \input{\rootfolder/code/fig-expansion-valve.tex}
    \caption{An expansion valve, with colors indicating the temperature and pressure of the flow.}
    \label{fig:expansion-valve}
\end{figure}

Because the valve contains no rotating shafts or moving boundaries that cross the control surface, no mechanical work is performed on the solid boundaries ($\dot{W}_\mathcal{S} = 0$). The mechanical work rate across the open inlet and outlet boundaries is also negligible ($\dot{W}_\mathcal{O} \approx 0$). Furthermore, due to the highly compact size of the device and the rapid transit of the fluid through the restriction, the residence time of the fluid inside the devince is short, allowing the heat transfer rate to the surrounding environment to be neglected ($\dot{Q} \approx 0$). Assuming changes in kinetic and potential energy between the inlet and outlet surfaces are small, the mass and steady-flow energy balances reduce to
\begin{align}
    \left.\begin{aligned}
        &\dot{m}_1=\dot{m}_{2,\text{liq}}+\dot{m}_{2,\text{vap}}\\
    &\dot{m}_1h_1\approx \dot{m}_{2,\text{liq}}h_{2,\text{liq}}+\dot{m}_{2,\text{vap}}h_{2,\text{vap}}
    \end{aligned}\right\} \implies h_1=(1-x)h_{2,\text{liq}}+x\,h_{2,\text{vap}}=h_2.
\end{align}
At the system level before separating the outlet streams $2_\text{liq}$ and $2_\text{vap}$, an expansion valve can be modeled as isenthalpic, i.e., as a constant--enthalpy process.  While enthalpy is approximately conserved across the device, viscous effects results in irreversible entropy generation ($\dot{S}_\text{irr} > 0$). 

\begin{example}
    In a heat pump, the mix of saturated liquid and vapor passing through the orifice of the expansion valve is separated into vapor and liquid streams, and the vapor stream directed to the compressor while the liquid stream first passes through the evaporator. Assume the CO$_2$ refrigerant enters the expansion valve at $T_1=20^\circ\text{C}$ and $p_1=90\,\text{bar}$ with a mass-flow rate of $\dot{m}=1\,\text{kg}/\text{s}$. At the expansion valve outlet, pressure has dropped to $p_2=35\,\text{bar}$. 
    
    In such case, the vapor quality can be calculated at the valve outlet and used to calculate the liquid and vapor stream mass-flow rates. This is conveniently done in a short script:    \lstinputlisting[language=Python,firstline=1,lastline=26,firstnumber=1]{\rootfolder/code/example-expansion-valve.py}
    Running the script reveals the outlet vapor quality to be $x_2=0.192$. Of the input stream of $\dot{m}=1\,\text{kg}/\text{s}$, evaporator therefore receives $0.808\,\text{kg}/\text{s}$ of saturated liquid while the vapor stream directed to the compressor equals $0.192\,\text{kg}/\text{s}$. As the process at the system level is approximated as isenthalpic, its corresponding state path can be drawn in the CO$_2$ $p$-$h$ diagram.
    \begin{center}
    \begin{tikzpicture}
    \node[above right,inner sep = 0] (image) at (0,0) {
        \includegraphics[width=0.85\linewidth]{\rootfolder/figs/example-expansion-valve.pdf}
    };
    \begin{scope}[relative to node=image]
        % \foreach \ratio in {-1,-0.9,...,1}{
        %    \draw[help lines,purple] (-1,\ratio) -- (1,\ratio);
        %    \draw[help lines,purple] (\ratio,-1) -- (\ratio,1);
        % }
        % \draw (-1,0) -- (1,0);
        % \draw (0,-1) -- (0,1);
        
        \coordinate (A) at (-0.41,0.37);
        \coordinate (B) at (-0.41,-0.54);
        \coordinate (Bliq) at (-0.6,-0.54);
        \coordinate (Bvap) at (0.38,-0.54);       

        \draw[dashed,blue,ultra thick] (Bliq)--(Bvap);
        
        \fill (A) circle (2pt) node[above=0.18cm, circle,fill=white,draw=black,inner sep=1pt] {1};
        \fill (B) circle (2pt) node[below = 0.18cm, circle,fill=white,draw=black,inner sep=1pt] {2};
        \fill (Bliq) circle (2pt) node[below = 0.18cm, circle,fill=white,draw=black,inner sep=1pt] {$2_{\text{liq}}$};
        \fill (Bvap) circle (2pt) node[below = 0.18cm, circle,fill=white,draw=black,inner sep=1pt] {$2_{\text{vap}}$};
        
    \end{scope}
    \end{tikzpicture}
    \end{center}
\end{example}


\section{Modeling of Power Plants}
Useful application for the system-level analysis of different components is found in power plants, such as steam power plants using the closed steady-flow Rankine cycle and gas power plants using the open steady-flow Brayton process. 

The simplified Rankine cycle is illustrated in Fig.~\ref{fig:rankine-cycle}. It contains four components and consists of the following main phases: 
\begin{figure}[!htb]
    \centering
    \input{\rootfolder/code/fig-rankine.tex}
    \caption{Rankine cycle.}
    \label{fig:rankine-cycle}
\end{figure}
\begin{itemize}
    \item[1--2] \textbf{Turbine:} Hot steam at high pressure and temperature (1) enters the turbine. Within the turbine, the steam enthalpy is converted to kinetic enerfy stage by stage and to rotational energy of the turbine shaft. The steam exits (2) the turbine at low pressure and temperature, close to the saturated vapor line, with a vapor quality often slightly below unity. 

    \item[2--3] \textbf{Condenser:} Arriving at the condenser (2) at low temperature and pressure, a small mass fraction of the steam might already be saturated liquid. Within the condenser, heat is removed from the flow to condense the rest of the vapor into saturated liquid. At the exit (3), water is saturated liquid. During this process, the coolant mass flow rate and the allowed temperature change have to be matched with the desired heat transfer rate enabling the condensing of the working fluid. 

    \item[3--4] \textbf{Pump:} Liquid water arrives (3) at a pumping system. At the pump exit (4), the liquid is near the same temperature as in the inlet and at the same pressure as at the turbine inlet (1) and about to enter the steam generator or boiler, the choice depending on the heat source.

    \item[4--1] \textbf{Steam generator / Boiler:} High-pressure liquid water enters (4) the steam generator, where it absorbs heat at a nearly constant pressure from an external high-temperature source (such as a combustion chamber or a primary nuclear coolant loop). At the outlet (1), the high-pressure, high-temperature steam is ready to enter the turbine and complete the cycle.
\end{itemize}

To estimate the efficiency of such process, the net work rate and the heat rate into the process are needed. These can be calculated by first analyzing each of the Rankine cycle components individually using the energy balance~\eqref{eq:energy-balance-control-volume}. Given the process simplifications, the energy balances over the components of a Rankine-cycle reduce to
\begin{align}
    \dot{m}(h_2-h_1)&=-\dot{W}_{1,2}, & &\text{turbine}\\
    \dot{m}(h_3-h_2)&=\dot{Q}_{2,3}, & &\text{condenser}\\
    \dot{m}(h_4-h_3)&=-\dot{W}_{3,4}, & &\text{pump}\\
    \dot{m}(h_1-h_4)&=\dot{Q}_{4,1}, & &\text{heat source}.
\end{align}
The cycle efficiency can therefore be calculated as
\begin{align}
    \eta_{\text{rankine}}&=\frac{\dot{W}_{1,2}+\dot{W}_{3,4}}{\dot{Q}_{4,1}}=\frac{h_1-h_2-(h_4-h_3)}{h_1-h_4}
\end{align}
\begin{example}
    Steam enters a turbine at a temperature of $T_1=550^\circ\text{C}$ and a pressure of $p_1=200\,\text{bar}$. At the turbine exit the steam pressure has dropped to $p_2=0.1\,\text{bar}$ and the vapor quality is $x_2=0.9$. After the condenser, the saturated liquid water enters the pump. At the pump exit, the temperature has increased by $T_4-T_3=5^\circ\text{C}$. The mass-flow rate of steam in the cycle is $\dot{m}=500\,\text{kg}/\text{s}$.

    To estimate the cycle efficiency, one first finds out the specific enthalpy values at the process points. This is most conveniently done using either \verb|pyromat| or \verb|coolprop|:
    \lstinputlisting[language=Python,firstline=1,lastline=36,firstnumber=1]{\rootfolder/code/example-rankine.py} 
    Once the enthalpy values are obtained, the process can be illustrated in the $p$-$h$ diagram as below. The solid line from state 4 to state 1 refers to the passing of water through the steam generator or boiler where the water pressure remains constant. The solid line from state 2 to state 3 refers to the water passing through the condenser where the pressure remains constant as well. The dashed lines connecting the states 1 and 2 and also the states 3 and 4 refer to water passing through the turbine and the pump. These lines are drawn dashed as they do not present the accurate average state of water passing the turbine. 
    \begin{center}
    \begin{tikzpicture}
    \node[above right,inner sep = 0] (image) at (0,0) {
        \includegraphics[width=0.85\linewidth]{\rootfolder/figs/example-rankine-ph.pdf}
    };
    \begin{scope}[relative to node=image]
        % \foreach \ratio in {-1,-0.9,...,1}{
        %    \draw[help lines,purple] (-1,\ratio) -- (1,\ratio);
        %    \draw[help lines,purple] (\ratio,-1) -- (\ratio,1);
        % }
        % \draw (-1,0) -- (1,0);
        % \draw (0,-1) -- (0,1);
        
        \coordinate (A) at (0.66,0.625);
        \coordinate (B) at (0.223,-0.453);
        \coordinate (C) at (-0.677,-0.453);
        \coordinate (D) at (-0.66,0.625);
        
        \fill (A) circle (2pt) node[above right=0.15cm and 0.15cm, circle,fill=white,draw=black,inner sep=1pt] {1};
        \fill (B) circle (2pt) node[below right=0.15cm and 0.15cm, circle,fill=white,draw=black,inner sep=1pt] {2};
        \fill (C) circle (2pt) node[below left=0.15cm and 0.15cm, circle,fill=white,draw=black,inner sep=1pt] {3};
        \fill (D) circle (2pt) node[above left=0.15cm and 0.15cm, circle,fill=white,draw=black,inner sep=1pt] {4};
       
    \end{scope}
    \end{tikzpicture}
    \end{center}
    The cycle efficiency and the heat and work rates of different components can then be calculated from the enthalpies and the mass-flow rate
    \lstinputlisting[language=Python,firstline=38,lastline=51,firstnumber=38]{\rootfolder/code/example-rankine.py}         
    For this specific cycle the thermal effiency is $\eta_\text{Rankine}=0.32$. The work rate in the turbine is $\dot{W}_{12}=525.8\,\text{MW}$, the positive value indicating the turbine produces work. The work rate in the pump $\dot{W}_{34}=-19.0\,\text{MW}$, the negative value indicating the pump requires work input. The heat rates for the condenser and steam generator are $\dot{Q}_{23}=-1076.4\,\text{MW}$ and $\dot{Q}_{41}=1583.2\,\text{MW}$, respectively, with the negative value for the condenser indicating heat is removed from the flow and the positive value in the steam generator indicating that heat is added to the flow. 
\end{example}

The simplified Brayton process is illustrated in Fig.~\ref{fig:brayton-process}. It contains three components and consists of the following main phases: 
\begin{figure}[!htb]
    \centering
    \input{\rootfolder/code/fig-brayton.tex}
    \caption{Brayton process.}
    \label{fig:brayton-process}
\end{figure}
\begin{itemize}
    \item[1--2] \textbf{Compressor:} Low-pressure ambient air at state (1) is drawn into the compressor, which typically consists of multiple axial or centrifugal stages (here depicted as axial). Passing through the stages, work input is needed to increase the air enthalpy. The compressed air exits the unit at high pressure at state (2) and at moderate temperature. 

    \item[2--3] \textbf{Combustion chamber / Combustor:} High-pressure air at state (2) enters the combustor, where fuel is continuously injected and ignited. This sets off an exothermic chemical reaction that delivers a high heat rate to the continuous flow. In ideal system-level modeling, this heat addition is assumed to occur at a constant pressure, neglecting minor aerodynamic pressure drops caused by fluid viscosity. Unlike the phase-change behavior observed in a steam generator, the working fluid remains entirely in the gaseous phase, but its temperature, specific volume, and volumetric flow rate increase drastically. This massive thermal expansion accelerates the gas toward the exit, yielding a high-velocity, high-temperature flue gas mixture at state (3) while maintaining the high operating pressure. The system-level energy balance treats this process as performing zero mechanical work, and the changes in macro-scale kinetic and potential energy between the inlet and outlet are neglected compared to the massive thermal energy input that drives the flow enthalpy change.

    \item[3--4] \textbf{Turbine:} The high-temperature, high-pressure gas mixture at state (3) enters the gas turbine to undergo expansion. The fluid mechanical process is similar to that inside a steam turbine, and for system-level modeling, the thermodynamic properties of the flue gas mixture are often assumed the same as of air. In contrast to steam power plants, where a minor fraction of the turbine mechanical power is used for efficient pressurization of the liquid water, a large fraction of the mechanical power generated by the gas turbine is transmitted internally via a common shaft to drive the compressor---a characteristic known as a high back-work ratio. The other significant difference to steam power plants is that the exhaust gas exits the turbine at state (4) at low pressure but still retains a high temperature. 
\end{itemize}

Similarly as in the analysis of a Rankine-cycle, the components of a Brayton process are treated individually as open steady-flow systems. Given the process simplifications, the energy balances over the components of a Brayton process reduce to
\begin{align}
    \dot{m}(h_2-h_1)&=-\dot{W}_{1,2}, & &\text{compressor}\\
    \dot{m}(h_3-h_2)&=\dot{Q}_{2,3}, & &\text{combustion chamber}\\
    \dot{m}(h_4-h_3)&=-\dot{W}_{3,4}, & &\text{turbine}
\end{align}
The process efficiency can therefore be calculated from the specific enthalpies of the different states similarly as for the Rankine-cycle
\begin{align}
    \eta_{\text{Brayton}}&=\frac{\dot{W}_{3,4}+\dot{W}_{1,2}}{\dot{Q}_{2,3}}=\frac{h_3-h_4-(h_2-h_1)}{h_3-h_2}
\end{align}
\begin{example}
    Ambient air enters the compressor of a Brayton process at a pressure of $p_1=1.0\,\text{bar}$ and at a temperature of $T_1=20^\circ\text{C}$ with a mass-flow rate of $\dot{m}=150\,\text{kg}/\text{s}$. After the compressor, the air pressure is $p_2=10\,\text{bar}$ and temperature $T_2=330^\circ\text{C}$. At the outlet from the combustion chamber, before the turbine inlet, the flue gas temperature is $T_3=1200^\circ\text{C}$. In the turbine, the pressure drops to the ambient level of $p_4=1\,\text{bar}$ while the flue gas temperature remains at $T_4=625^\circ\text{C}$. 

    The process efficiency can be estimated from the specific enthalpy values of gaseous air during different states: 
    \lstinputlisting[language=Python,firstline=1,lastline=43,firstnumber=1]{\rootfolder/code/example-brayton.py} 
    For this specific process the thermal efficiency is $\eta_\text{Brayton}=0.357$. The work rate in the compressor is $\dot{W}_{12}=-158.6\,\text{MW}$, the negative value indicating the compressor to require mechanical work. The work rate of the turbine is $\dot{W}_{34}=335.7\,\text{MW}$, the positive value indicating the turbine produces work. Note that roughly half of the turbine work rate is consumed directly by the compressor. 
\end{example}

The example Rankine cycle achieved a thermal efficiency of $\eta_\text{Rankine}=0.32$ in a process where the minimum and maximum temperatures of the working fluid where roughly $T_\text{min}=46^\circ\text{C}$ and $T_\text{max}=550^\circ\text{C}$. If all of the heat addition would happen at the maximum temperature, all of the heat removal at the minimum temperature, and the turbine and pump operated isentropically with no irreversible entropy generation, the theoretical maximum achievable efficiency would be the Carnot efficiency with a value of $\eta_\text{Carnot}=0.612$. The achieved value is roughly only half of the theoretical maximum. Similarly, if the minimum and maximum working fluid temperatures in the example Brayton process would be taken as the ambient air inlet temperature $T_\text{min}=20^\circ\text{C}$ and the flue gas combustion chamber outlet temperature $T_\text{max}=1200^\circ\text{C}$, the theoretical Carnot efficiency for such process would be as high as $\eta_\text{Carnot}=0.80$. Clearly neither of the example processes is maximizing their potential yet. The difference to hypothetical Carnot processes can be visualized illustrating the processes in respective $T$-$s$ diagrams. 
\begin{example}
    The $T$-$s$ diagram for the example Rankine process is illustrated below, with the process curve drawn in blue and the state points merked separately. Accompanying the process curve is the corresponding Carnot rectangle associated with the process minimum and maximum temperature, drawn with dashed orange curve. 
    
    \begin{center}
    \begin{tikzpicture}
    \node[above right,inner sep = 0] (image) at (0,0) {
        \includegraphics[width=0.85\linewidth]{\rootfolder/figs/example-rankine-Ts.pdf}
    };
    \begin{scope}[relative to node=image]
        % \foreach \ratio in {-1,-0.9,...,1}{
        %    \draw[help lines,purple] (-1,\ratio) -- (1,\ratio);
        %    \draw[help lines,purple] (\ratio,-1) -- (\ratio,1);
        % }
        % \draw (-1,0) -- (1,0);
        % \draw (0,-1) -- (0,1);
        
        \coordinate (A) at (0.112,0.343);
        \coordinate (B) at (0.262,-0.687);
        \coordinate (C) at (-0.69,-0.687);
        \coordinate (D) at (-0.685,-0.679);

        % Process point annotations
        \fill (A) circle (2pt) node[above right=0.15cm and 0.15cm, circle,fill=white,draw=black,inner sep=1pt] {1};
        \fill (B) circle (2pt) node[below right=0.15cm and 0.15cm, circle,fill=white,draw=black,inner sep=1pt] {2};
        \fill (C) circle (2pt) node[below left=0.15cm and 0.15cm, circle,fill=white,draw=black,inner sep=1pt] {3};
        \fill (D) circle (2pt) node[above left=0.15cm and 0.15cm, circle,fill=white,draw=black,inner sep=1pt] {4};

        % Carnot rectanle
        \draw[orange, ultra thick,dashed] (B)--(C)|-(A)-| cycle;

    \end{scope}
    \end{tikzpicture}
    \end{center}

    One observes the area the blue Rankine-process curve spans not to fill the entire Carnot rectangle. In fact, the area corresponds roughtly to half of the area of the dashed-orange rectangle, i.e., to the ratio of the efficiencies $\eta_\text{Rankine}$ and $\eta_\text{Carnot}$ for the process. It is impossible for the Rankine cycle to fill the Carnot rectangle since water above the pressure of $5000\,\text{bar}$ and temperature $T=0^\circ\text{C}$ effectively becomes solid and no longer can flow through even a theoretical heat exchanger. 

    The $T$-$s$ diagram of the example Brayton process is illustrated below, with the process curve is drawn in blue. If the process were closed by cooling the turbine exhaust gas from state point 4 isobarically in a heat exchanger back to the temperature of state point point 1, the blue curve would be extend along the $1\,\text{bar}$ isobar to form a loop and enclose an area. The enclosed area would significantly smaller that the area spanned by the orange dashed rectangle corresponding to the respective Carnot process. From this illustration, one understands that filling the rectangle area is difficult since heat addition and removel happen during isobaric processes and in gaseous phase those processes do not align with isothermal processes. Even 
    \begin{center}
    \begin{tikzpicture}
    \node[above right,inner sep = 0] (image) at (0,0) {
        \includegraphics[width=0.85\linewidth]{\rootfolder/figs/example-brayton-Ts.pdf}
    };
    \begin{scope}[relative to node=image]
        % \foreach \ratio in {-1,-0.9,...,1}{
        %    \draw[help lines,purple] (-1,\ratio) -- (1,\ratio);
        %    \draw[help lines,purple] (\ratio,-1) -- (\ratio,1);
        % }
        % \draw (-1,0) -- (1,0);
        % \draw (0,-1) -- (0,1);
        
        \coordinate (A) at (-0.465,-0.7);
        \coordinate (B) at (-0.405,-0.32);
        \coordinate (C) at (0.44,0.73);
        \coordinate (D) at (0.513,0.033);

        % Process point annotations
        \fill (A) circle (2pt) node[left=0.15cm, circle,fill=white,draw=black,inner sep=1pt] {1};
        \fill (B) circle (2pt) node[above=0.15cm, circle,fill=white,draw=black,inner sep=1pt] {2};
        \fill (C) circle (2pt) node[above=0.15cm, circle,fill=white,draw=black,inner sep=1pt] {3};
        \fill (D) circle (2pt) node[right=0.15cm, circle,fill=white,draw=black,inner sep=1pt] {4};

        % Carnot rectanle
        \draw[orange, ultra thick,dashed] (A)|-(C)-|(D)|- cycle;

    \end{scope}
    \end{tikzpicture}
    \end{center}
\end{example}


\section{Modeling of Refrigerators and Heat Pumps}
In addition to power plants, the steady-state balance equations have their use in the system-level analysis of Refrigerators and heat pumps that utilize external work to transfer heat from lower temperature to higher temperature. Both refrigerators and heat pumps use exactly the same working principle and are effectively based on inversion of the Rankine process. 

A simplified sub-critical refrigeration or heat pump cycle is illustrated in Fig.~\ref{fig:heat-pump}. It requires four basic components and consists of the following main phases:
\begin{figure}[!htb]
    \centering
    \input{\rootfolder/code/fig-heat-pump.tex}
    \caption{A sub-critical refrigerator or heat pump process with colors indicative of temperature.}
    \label{fig:heat-pump}
\end{figure}
\begin{itemize}
    \item[1--2] \textbf{Compressor:} Low-pressure, low-temperature refrigerant stream (5) from the evaporator in a state of saturated or slightly superheated vapor is joined to the stream ($4_\text{vap}$) of equal pressure saturated vapor from the expansion valve before entering the compressor at state (1). As in the Brayton cycle, mechanical work is consumed to compress the refrigerant, resulting in a high-pressure, high-temperature superheated vapor at the compressor exit (2). In sub-critical processes, the exit pressure is below the refrigerant critical point pressure.

    \item[2--3] \textbf{Condenser:} The hot, high-pressure vapor enters the condenser (2), where it releases heat to the an external high temperature stream. During this process, the refrigerant undergoes a phase change from superheated vapor to saturated or subcooled liquid at constant pressure, while the other stream gains in temperature. At the exit (3), the refrigerant is a high-pressure liquid.

    \item[3--4] \textbf{Expansion valve:} The high-pressure liquid at (3) passes through an expansion valve. This is an isenthalpic process where the pressure of the refrigerant is rapidly reduced. As a result, the refrigerant enters a low-pressure, low-temperature liquid-vapor mixture state (4) before being split into saturated liquid ($4_\text{liq}$) and vapor ($4_\text{vap}$) streams. The vapor quality of the mixture varies based on the cycle details.

    \item[4--1] \textbf{Evaporator:} The low-temperature, low-pressure saturated liquid stream ($4_\text{liq}$) is directed to the evaporator where it absorbs heat from the a cold external stream at constant pressure. The saturated liquid stream evaporates into saturated or slightly superheated vapor state (5). The saturated vapor stream ($4_\text{vap}$) from the expansion valve bypasses the evaporator and joins the evaporator stream (5) to form the state (1) ready to enter the compressor to complete the cycle.
\end{itemize}

The components of the cycle are treated as open steady-flow systems. Given the process simplifications, the balance equations over the compressor, condenser, and the expansion valve provide
\begin{align}
    \dot{m}(h_2-h_1)&=-\dot{W}_{1,2}, & &\text{compressor}\\
    \dot{m}(h_3-h_2)&=\dot{Q}_{2,3}, & &\text{condenser}\\
    h_4-h_3&=0, & &\text{expansion valve}.
\end{align}
Between the states (4) and (1), one considers the splitting and rejoining of the stream. Mass balance provides $\dot{m}=\dot{m}_\text{liq}+\dot{m}_\text{vap}$, with $\dot{m}_\text{liq}$ and $\dot{m}_\text{vap}$ describing the (positively valued) mass-flow rates for the streams ($4_\text{liq}$) and ($4_\text{vap}$), respectively. Energy balance provides three equations: the split into two streams $\dot{m}h_4=\dot{m}_\text{liq}h_{4_\text{liq}}+\dot{m}_\text{vap}h_{4_\text{vap}}$; the evaporator $\dot{m}_\text{liq}(h_5-h_{4_\text{liq}})=\dot{Q}_{4_\text{liq},5}$; and the joining of the streams back together $\dot{m}h_1=\dot{m}_{4_\text{vap}}h_{4_\text{vap}}+\dot{m}_{\text{liq}}h_5$. Subtracting the first from the last equation, and using the middle one then provides
\begin{align}
    \dot{m}(h_1-h_4)=\dot{Q}_{4_\text{liq},5}, & &\text{evaporator}.
\end{align}
At system level analysis, it is therefore not always necessary to known the details of the vapor and liquid streams.

The measurement of efficiency in refrigerators and heat pumps differs from that of power plants. In power plants, one pays for heat and gets work. In refrigerators and heat pumps, one pays for work and gets either cooling or heating. Instead of thermal efficiency, one is interested in the coefficient of performance, the COP-value. In terms of the process described above, they can be defined
\begin{align}
    \text{COP}_\text{heating}&=\frac{-\dot{Q}_{2,3}}{-\dot{W}_{1,2}}=\frac{h_2-h_3}{h_2-h_1},\\
    \text{COP}_\text{cooling}&=\frac{\dot{Q}_{4_\text{liq},5}}{-\dot{W}_{1,2}}=\frac{h_1-h_4}{h_2-h_1}.
\end{align}
The minus-signs have been incorporated to account for the used sign convention, with the compressor work rate and condenser heat rate being negative.

\begin{example}\label{example:sub-critical-heat-pump}
    Heat pump uses R-1234ze as refrigerant. The evaporator pressure is adjusted to $2\,\text{bar}$ so that the temperature at compressor inlet is $1^\circ\text{C}$, slightly above the refrigerant saturation temperature of $-2.1^\circ\text{C}$ at the same pressure. With the expansion valve, the pressure for the condenser is adjusted to $10\,\text{bar}$, and the refrigerant enters the expansion valve at a temperature of $47^\circ\text{C}$, slightly below the saturation temperature of $50.1^\circ\text{C}$ at the condenser pressure. The compressor isentropic efficiency has been estimated to be $0.75$.

    To estimate the COP-value for heating, one first finds out the enthalpy values at the process state points. This is most conveniently done with a short script.
    \lstinputlisting[language=Python,firstline=1,lastline=31,firstnumber=1]{\rootfolder/code/example-heat-pump.py} 
    From the enthalpy values, one can then calculate the COP-value for heating and print out also the vapor quality at the expansion valve flash tank:
    \lstinputlisting[language=Python,firstline=33,lastline=38,firstnumber=33]{\rootfolder/code/example-heat-pump.py} 
    Running the script provides $\text{COP}_\text{heating}=3.95$ and a vapor quality of 0.37 after the expansion valve. For every kW of compressor power, one therefore receives 4\,kW of heating rate. Further, roughly one-third of the mass flow rate in the process bypasses the evaporator. The process path is illustrated in the R-1234ze $p$-$h$ diagram below in blue color.
    \begin{center}
    \begin{tikzpicture}
    \node[above right,inner sep = 0] (image) at (0,0) {
        \includegraphics[width=0.85\linewidth]{\rootfolder/figs/example-heat-pump-ph.pdf}
    };
    \begin{scope}[relative to node=image]
        % \foreach \ratio in {-1,-0.9,...,1}{
        %    \draw[help lines,purple] (-1,\ratio) -- (1,\ratio);
        %    \draw[help lines,purple] (\ratio,-1) -- (\ratio,1);
        % }
        % \draw (-1,0) -- (1,0);
        % \draw (0,-1) -- (0,1);
        
        \coordinate (A) at (0.41,-0.533);
        \coordinate (B) at (0.63,0.037);
        \coordinate (C) at (-0.227,0.037);
        \coordinate (D) at (-0.227,-0.533);

        % Process point annotations
        \fill (A) circle (2pt) node[below right=0.15cm, circle,fill=white,draw=black,inner sep=1pt] {1};
        \fill (B) circle (2pt) node[above right=0.15cm, circle,fill=white,draw=black,inner sep=1pt] {2};
        \fill (C) circle (2pt) node[above left=0.15cm, circle,fill=white,draw=black,inner sep=1pt] {3};
        \fill (D) circle (2pt) node[below left=0.15cm, circle,fill=white,draw=black,inner sep=1pt] {4};

    \end{scope}
    \end{tikzpicture}
    \end{center}
    
\end{example}

In the heat pump described in example~\ref{example:sub-critical-heat-pump}, the majority of heat transfer from the refrigerant to the external stream in the condenser is yielded via a phase transition, at constant temperature. The stream receiving the heat typically remains in single phase, and the end-use of it may vary from one application to another. Generally, heat pumps that yield heat primarly via phase transition are best suited for processes where the receiving stream does not need to undergo a significant temperature change. Such processes involve for example the heating of the air circulating inside apartments. In applications where the receiving stream is to undergo significant temperature change, also the refrigerant can be allowed remain in single phase during heat yielding. This enables use of transcritical process cycles, where the refrigerant is compressed beyond the critical point pressure. Such processes enable high temperatures to be reached even with refrigerants with moderately low critical point temperature.

A transcritical heat pump with an additional internal heat exchanger is illustrated in Fig.~\ref{fig:heat-pump-transcritical}. Inplace of a condenser, a regular gas cooling heat exchanger is used as the primary heat exchanger instead. The process consists of the following phases:
\begin{figure}[!htb]
    \centering
    \input{\rootfolder/code/fig-heat-pump-transcritical.tex}
    \caption{A transcritical heat pump with an internal heat exchanger to superheat the stream from evaporator.}
    \label{fig:heat-pump-transcritical}
\end{figure}
\begin{itemize}
    \item[1--2] \textbf{Internal heat exchanger (cold side):} The low-pressure vapor stream at state (1), formed by rejoining the evaporator stream (7) and the expansion valve vapor stream ($6_\text{vap}$), enters the cold side of the internal heat exchanger (IHX). The stream is superheated at constant pressure to state (2) by absorbing heat from the high-pressure side. This superheating shifts the compressor inlet state further into the vapor region to prevent liquid slugging and improves the process COP-value.

    \item[2--3] \textbf{Compressor:} The superheated vapor enters the compressor at state (2), where mechanical work is consumed to compress the refrigerant to a supercritical pressure well above its critical point. The refrigerant exits at state (3) as a highly superheated, high-pressure supercritical fluid.

    \item[3--4] \textbf{Gas cooler:} The supercritical fluid enters the primary heat exchanger at state (3). Because the process operates above the critical pressure, the refrigerant does not condense at a constant temperature plateau. Instead, it remains a single-phase gas and experiences a substantial, continuous temperature glide at constant pressure while rejecting heat to the external high-temperature stream, exiting at state (4).

    \item[4--5] \textbf{Internal heat exchanger (hot side):} The high-pressure fluid at state (4) passes through the hot side of the IHX. It transfers heat to the low-pressure side, enabling significant vapor side superheating (1--2), and cools down further to state (5).

    \item[5--6] \textbf{Expansion valve and flash tank:} The high-pressure fluid at state (5) undergoes an isenthalpic expansion through the expansion valve into the flash tank at state (6). The sudden pressure drop converts the fluid into a low-pressure, low-temperature liquid-vapor mixture, which naturally separates into a saturated liquid stream ($6_\text{liq}$) and a saturated vapor stream ($6_\text{vap}$) due to gravity.

    \item[6--1] \textbf{Evaporator and mixing:} The saturated liquid stream ($6_\text{liq}$) is routed to the evaporator, where it absorbs heat from a low-temperature external source at constant pressure, evaporating completely into a vapor state (7). This stream is then mixed with the bypassed flash vapor ($6_\text{vap}$) to return the refrigerant to state (1), closing the transcritical loop.
\end{itemize}

In the transcritical process, the balance equations function exactly in the same manner as in the subcrictical process, and the COP-values are determined by the enthalpy differences over the compressor for the required work rate and over the primary gas cooler for heating processes and the evaporator for cooling processes. 
\begin{example}
    In a transcritical CO$_2$ heat pump containing internal heat exchanger, the evaporator pressure is set to $30\,\text{bar}$ and the compressor outlet pressure to $90\,\text{bar}$. In the gas cooler, the CO$_2$ exit temperature is $50^\circ\text{C}$ before the hot side inlet of the internal heat exchanger. The temperature of the joined evaporator and expansion valve vapor streams is $-4^\circ\text{C}$ before the cold side inlet of the internal heat exchanger, slightly above the saturation temperature of $-5.6^\circ\text{C}$ for the respective evaporator pressure. For the internal heat exchanger to function properly, the temperature difference between the streams at inlet and outlet is designed to be at least $4^\circ\text{C}$. The compressor isentropic efficiency is estimated at $0.75$.

    To estimate the process heating COP-value, one starts by determining the known states at the gas cooler exit and the internal heat exchanger cold side inlet:
    \lstinputlisting[language=Python,firstline=1,lastline=17,firstnumber=1]{\rootfolder/code/example-heat-pump-transcritical.py} 
    Given the minimum stream temperature difference for the internal heat exchanger inlet and outlet, one determines the potentially maximum possible cold side outlet temperature, the potentially minimum possible hot side outlet temperature, and their correponding state enthalpies. Since the mass flow rates to the internal heat exchanger are equal for the hot and cold streams, the absolute values of the enthalpy changes must be equal too. This sets the mutual enthalpy exchange between the streams to be the smaller of the two potential maximal enthalpy changes:
    \lstinputlisting[language=Python,firstline=19,lastline=33,firstnumber=19]{\rootfolder/code/example-heat-pump-transcritical.py} 
    Once the true enthalpy exchange in the internal heat exchanger is known, one may proceed to determine the stream states at the hot and cold side outlets, calculate the compressor outlet state, the expansion valve outlet state, and the resulting COP-value
    \lstinputlisting[language=Python,firstline=35,lastline=58,firstnumber=35]{\rootfolder/code/example-heat-pump-transcritical.py} 
    Running the script prints $\text{COP}_\text{heating}=1.97$. The process state path is illustrated below in a CO$_2$ $p$-$h$ diagram, the internal heat exchanger in orange color and the rest of the path in blue. The figure reveals that without the internal heat exchanger, the expansion valve would produce a mixture with barely any liquid in it and make it difficult injecting heat into the process in the first place. Further, due to the superheating of the vapor, the process path is stretched to the right, significantly increasing the achievable highest temperature.
    \begin{center}
    \begin{tikzpicture}
    \node[above right,inner sep = 0] (image) at (0,0) {
        \includegraphics[width=0.85\linewidth]{\rootfolder/figs/example-heat-pump-transcritical-ph.pdf}
    };
    \begin{scope}[relative to node=image]
        % \foreach \ratio in {-1,-0.9,...,1}{
        %    \draw[help lines,purple] (-1,\ratio) -- (1,\ratio);
        %    \draw[help lines,purple] (\ratio,-1) -- (\ratio,1);
        % }
        % \draw (-1,0) -- (1,0);
        % \draw (0,-1) -- (0,1);
        
        \coordinate (A) at (0.285,-0.617);
        \coordinate (B) at (0.51,-0.617);
        \coordinate (C) at (0.83,0.365);
        \coordinate (D) at (0.20,0.365);
        \coordinate (E) at (-0.03,0.365);
        \coordinate (F) at (-0.03,-0.617);

        % Process point annotations
        \fill (A) circle (2pt) node[below=0.15cm, circle,fill=white,draw=black,inner sep=1pt] {1};
        \fill (B) circle (2pt) node[below right=0.15cm, circle,fill=white,draw=black,inner sep=1pt] {2};
        \fill (C) circle (2pt) node[above right=0.15cm, circle,fill=white,draw=black,inner sep=1pt] {3};
        \fill (D) circle (2pt) node[above=0.15cm, circle,fill=white,draw=black,inner sep=1pt] {4};
        \fill (E) circle (2pt) node[above left=0.15cm, circle,fill=white,draw=black,inner sep=1pt] {5};
        \fill (F) circle (2pt) node[below left=0.15cm, circle,fill=white,draw=black,inner sep=1pt] {6};

    \end{scope}
    \end{tikzpicture}
    \end{center}    
\end{example}

\section{Modeling of reciprocating engines}
A convenient and useful application for the mass and energy balance equations is found in the analysis of reciprocating engines. The four-stroke internal combustion engine process is illustrated in Fig.~\ref{fig:internal-combustion-engine} and consists of the following phases and events:
\begin{itemize}
    \item[1--3] \textbf{Intake Stroke:} With the piston at or near the Top Dead Center, the intake valve opens (1). The piston moves downward, resulting in a small pressure difference between the engine chamber and the intake manifold, drawing fresh air or an air--fuel mixture into the engine cylinder (1--3), after which the intake valve closes (3) at or near the Bottom Dead Center.
    \item[3--6] \textbf{Compression Stroke:} Once the valves are closed (3), the piston moves upward to compress the trapped air or the air--fuel mixture (3--6). In direct-injection engines, fuel is injected during this phase. Near the end of the compression stroke (6), just before the Top Dead Center, the fuel is ignited either by a spark plug or by the high temperature rise from the rapid compression causing the fuel to self-ignite.
    \item[6--7/8] \textbf{Combustion and Heat Addition:} Combustion then occurs as the piston passes through the Top Dead Center (6--7/8). If the combustion happens fast (6--7), pressure, temperature, and internal energy of the gas mixture all rise at a nearly constant volume. If the combustion is more controlled taking longer (6--8), the phase carries partially over the following expansion stroke. With the cylinder volume expanding as the gas receives energy, pressure and temperature rise are less sharp.
    \item[7--9] \textbf{Expansion Stroke:} The high-pressure combustion gases expand, pushing the piston downward and performing work on the crankshaft (7--9). Near the end of the expansion phase, slightly before the Bottom Dead Center, the exhaust valve opens~(9).
    \item[9--10] \textbf{Exhaust Blowdown:} The pressure difference between the engine chamber and the exhaust manifold causes a large portion of the exhaust gas to rush out spontaneously (9--10), resulting in a rapid drop in the remaining gas pressure and temperature.
    \item[10--12] \textbf{Exhaust Stroke:} The piston moves upward again, mechanically pushing most of the remaining burned gases out of the cylinder (10--12) until the exhaust valve closes near the Top Dead Center, allowing the cycle to repeat.
\end{itemize}
Accurate modeling of internal combustion engines involves detailed spatial and temperoral computational analysis of the fluid flow inside the engine cylinder and the chemical composition changes involved in the combustion process. A reasonable understanding of the working principle of such engines, including an estimate for the engine thermal efficiency and the timing of heat release phase, can nevertheless be achieved also with simplified models. Such simplified modeling involves experimental thermodynamic data from the engine process and certain assumptions. 
\begin{figure}[!htb]
    \centering
    \input{\rootfolder/code/fig-engine.tex}
    \caption{A four-stroke internal combustion engine process.}
    \label{fig:internal-combustion-engine}
\end{figure}

The approximations one introduces involve the neglect of chemical composition changes. The energy released in combustion is modeled as volumetric addition of heat. This enables the use of single-substance fluid equations and thermodynamic relations. Typically, the working fluid is approximated as an ideal gas with properties mimicking those of air. In the energy-balance equation~\eqref{eq:energy-balance-control-volume}, one neglects the kinetic and gravitational potential energy terms ($dE_k/dt\ll dU/dt$, $dE_p/dt\ll dU/dt$, $\bar{e}_{k,j}\ll\bar{h}_{j}$, $\bar{e}_{p,j}\ll\bar{h}_{j}$), assuming them to be negligible in comparison to internal energy and enthalpy. The work rates on the open intake and exhaust surfaces are assumed small in comparison to the work rate on the solid engine surface ($\dot{W}_{\mathcal{O}}\ll\dot{W}_{\mathcal{S}}$). The viscous effects are neglected and the work is assumed to result from pressure alone. Then, pressure, internal energy, and the inflow and outflow enthalpies are assumed to be uniform in space within the engine volume and on the intake and exhaust surfaces. As a results, the work rate reduces to $\dot{W}_{\mathcal{S}}\approx pdV/dt$ and the entire energy balance equation to
\begin{align}
    \frac{dU}{dt}+\dot{m}_\text{out}(t)h_\text{out}(t)-\dot{m}_\text{in}(t)h_\text{in}(t)=\dot{Q}(t)-p(t)\frac{dV}{dt}.
\end{align}
The energy equation is then accompanied by the mass balance equation
\begin{align}
    \frac{dm}{dt}+\dot{m}_\text{out}(t)-\dot{m}_\text{in}(t)=0,
\end{align}
with $m$ describing the air mass inside the engine chamber. The sign convention for the inflow and outflow mass-flow rates $\dot{m}_\text{in}$ and $\dot{m}_\text{out}$ is here chosen so that they both remain non-negative.

With the working fluid assumed ideal gas, the specific enthalpy for the inflow and outflow become $h_\text{in}=c_p(T_\text{in})T_\text{in}$ and $h_\text{out}=c_p(T)T$ where $T_\text{in}$ is the temperature assumed for the intake air. The air temperature $T$ in the cylinder relates to the pressure, volume, and mass via the ideal gas equation of state $T=(pV)/(mR_s)$, and the internal energy for the volume can be written $U=mc_vT$. In experiments, the data is typically recorded with respect to the crank angle $\theta$ that rotates with the angular velocity
\begin{align}
    \omega=\frac{d\theta}{dt}=\frac{2\pi\,N_\text{RPM}}{60}\,(\text{rad}/\text{s})=\frac{360\,N_\text{RPM}}{60}\,(^\circ/\text{s}),
\end{align}
with $N_\text{RPM}$ denoting the number of revolutions per minute. Transforming to crank angle, the energy and mass balances are written
\begin{align}
    \omega\frac{dU}{d\theta}+\dot{m}_\text{out}(\theta)h_\text{out}(\theta)-\dot{m}_\text{in}h_\text{in}&=\dot{Q}(\theta)-\omega p(\theta)\frac{dV}{d\theta}\\
    \omega\frac{dm}{d\theta}+\dot{m}_\text{out}(\theta)-\dot{m}_\text{in}(\theta)&=0
\end{align}
Provided experimentally measured data $p(\theta)$, $V(\theta)$, $\dot{m}_\text{in}(\theta)$, and $\dot{m}_\text{out}(\theta)$, one may calculate the net work and estimate the injected heat
\begin{align}
    W&=\int_{\theta_0}^{\theta_0+\Delta\theta}p(\theta)\frac{dV}{d\theta}d\theta, &Q_\text{in}&=\int_{\theta_0}^{\theta_0+\Delta\theta}\max(\dot{Q}(\theta),0)\frac{d\theta}{\omega},
\end{align}
where the cycle period is $\Delta\theta=4\pi=720^\circ$. From the net work and injected heat, the engine efficiency can then be estimated as
\begin{align}
    \eta=\frac{W}{Q_\text{in}}.
\end{align}
Characteristic pressure, volume, and mass-flow-rate data for a four-stroke engine is illustrated in Fig.~\ref{fig:four-stroke-pVmdot-data} as a function of the crankshaft angle. The corresponding $p$-$V$ diagram of the cycle is shown in Fig.~\ref{fig:four-stroke-pV-diagram}.
\begin{figure}[!htb]
    \centering
    \includegraphics[width=0.85\linewidth]{\rootfolder/figs/fig-engine-pVmdot-theta.pdf}
    \caption{Four-stroke engine data as a function of crankshaft angle over the cycle period $\theta\in[-180^\circ,540^\circ]$.}
    \label{fig:four-stroke-pVmdot-data}
\end{figure}
\begin{figure}[!htb]
    \centering
    \includegraphics[width=0.6\linewidth]{\rootfolder/figs/fig-engine-pV.pdf}
    \caption{A $p$-$V$ diagram of the data of Fig.~\ref{fig:four-stroke-pVmdot-data} showing both the compression--expansion and the exhaust--intake loops.}
    \label{fig:four-stroke-pV-diagram}
\end{figure}

\begin{example}
    Given a point-wise data set $\{\theta_i,p_i,V_i,\dot{m}_{\text{in},i},\dot{m}_{\text{out},i}\}_{i=0}^N$ corresponding to the four-stroke cycle shown in Fig.~\ref{fig:four-stroke-pVmdot-data}, and the engine $N_\text{RPM}$ value, estimation of the engine efficiency can be performed in the following manner. 
    
    Load libraries to handle numerical data, to perform numerical integration, and to evaluate ideal gas data:
    \lstinputlisting[language=Python,firstline=36,lastline=46,firstnumber=1]{\rootfolder/code/fig-engine.py}
    Read the data and estimate the mass of air within the engine. For this, assume air to be at temperature $T_\text{ivc}$ right after the closing of the intake valve ($\theta_\text{ivc}$), and calculate the corresponding instantaneous mass $m_\text{ivc}$ from the equation of state $m_\text{ivc}=p_\text{ivc}V_\text{ivc}/(T_\text{ivc}R_s)$. Then calculate the change in the mass from the beginning of the data set
    \begin{align*}
        \Delta m_{[\theta_0,\theta]}=\int_{\theta_0}^{\theta}(\dot{m}_\text{in}-\dot{m}_\text{out})\frac{d\theta}{\omega}
    \end{align*}
    and evaluate the mass at any angle $\theta$ from
    \begin{align*}
        m(\theta)=m_\text{ivc}-\Delta m_{[\theta_0,\theta_\text{ivc}]}+\Delta m_{[\theta_0,\theta]}
    \end{align*}
    \lstinputlisting[language=Python,firstline=48,lastline=77,firstnumber=13]{\rootfolder/code/fig-engine.py}
    Reconstruct the gas temperature using pressure, volume and mass data, calculate internal energy, enthalpy, work and heat rates, and integrate the net work and heat injected. Finally, estimate the engine efficiency
    \lstinputlisting[language=Python,firstline=79,lastline=95,firstnumber=44]{\rootfolder/code/fig-engine.py}
    The resulting effiency is $\eta\approx 0.46$, and the calculated data is visualized below. From the curves for $\dot{Q}$ and $\dot{W}$, one clearly recognizes the heat injection phase and how the compression phase results in negative work rate. The intake and exhaust phases also require and produce some work, but the net work requirement from exhaust--intake loop is relatively moderate. This is reflected in the $p$-$V$ diagram of Fig.~\ref{fig:four-stroke-pV-diagram} where the area of the compression--expansion loop is observed to be significantly larger than the area of the exhaust--intake loop.
    \begin{center}
        \includegraphics[width=0.85\linewidth]{\rootfolder/figs/example-engine-mTQdotWdot-theta.pdf}
    \end{center}
    
\end{example}

\section{Summary}

The conservation laws and thermodynamic relations developed in the
previous chapters can be combined to analyze engineering systems
containing flowing fluids. A volume with solid and open surfaces provides
a general control volume for such analyses. The mass, momentum, angular
momentum, energy, and entropy balances describe the accumulation,
transport, transfer, and generation of the corresponding physical
quantities within the volume.

For many engineering devices, the general balance equations can be
simplified by introducing suitable modeling assumptions. Steady flow
removes temporal accumulation, while negligible changes in kinetic and
potential energy, negligible heat losses, or negligible mechanical work
at open surfaces can reduce the balances to simple relations between the
inlet and outlet states. For a steady-flow component with negligible
changes in kinetic and potential energy and negligible work associated
with the open surfaces, the energy balance reduces to
\begin{align}
    \dot{m}(h_2-h_1) = \dot{Q}-\dot{W}.
\end{align}
This relation connects the enthalpy change of the flowing fluid to the
heat and mechanical power rates of the component and provides the basis
for modeling turbines, compressors, pumps, valves, and heat exchangers.

Combining component models makes it possible to analyze complete
thermodynamic cycles. Power plants convert heat into mechanical work,
while refrigerators and heat pumps use mechanical work to transfer heat
between thermal reservoirs. Their performance can be characterized by
thermal efficiency and coefficient of performance, with the thermal
efficiency of a cycle defined by
\begin{align}
    \eta = \frac{\dot{W}_{\mathrm{net}}}{\dot{Q}_{\mathrm{in}}}.
\end{align}
Entropy balance provides an additional measure of the irreversible
behaviour of the processes and establishes an upper limit for achievable
performance. In the ideal quasistatic limit, heat addition and rejection
at the highest and lowest available temperatures lead to the Carnot
efficiency
\begin{align}
    \eta_{\mathrm{Carnot}}
    = 1-\frac{T_{\min}}{T_{\max}}.
\end{align}

The same framework can also be applied to systems for which detailed
cycle state points are not known a priori. Experimentally measured
pressure, volume, and mass-flow-rate data can be combined with the
balance equations and thermodynamic property relations to estimate the
net work, heat input, and efficiency of reciprocating engines. In this
way, the conservation laws, thermodynamic closure, and experimentally
determined or library-based property data form a unified framework for
analyzing and evaluating the performance of practical energy-conversion
devices.